\documentclass[a4paper,12pt]{article}

%---------------------------------
% Кодировка и язык
\usepackage[T2A]{fontenc}
\usepackage[utf8]{inputenc}
\usepackage[russian]{babel}
\usepackage[unicode, colorlinks=true, linkcolor=blue, urlcolor=blue]{hyperref}
\usepackage{cmap} % для поиска кириллицы в PDF

%---------------------------------
% Математика и графика
\usepackage{amsmath, amssymb, amsthm}
\usepackage{graphicx}
\usepackage[most]{tcolorbox}
\usepackage[dvipsnames]{xcolor}
\usepackage{tikz}
\usetikzlibrary{arrows, decorations.pathreplacing, positioning, shapes.geometric}
\usepackage{pgfplots}
\pgfplotsset{compat=1.18}
\usepackage{caption}
\captionsetup{hypcap=false}

%---------------------------------
% Геометрия страницы
\usepackage{geometry}
\geometry{left=2.5cm, right=2.5cm, top=2.5cm, bottom=2.5cm}

%---------------------------------
% Списки и заголовки
\usepackage{enumitem}
\usepackage{titlesec}
\usepackage{tocloft}

%---------------------------------
% Колонтитулы
\usepackage{fancyhdr}
\usepackage{lastpage}

%---------------------------------
% Команда для защиты математики в оглавлении
\newcommand{\tsubsection}[2]{\subsection[#1]{\texorpdfstring{#2}{#1}}}

%---------------------------------
% Настройка содержания
\renewcommand{\contentsname}{Содержание}
\renewcommand{\cfttoctitlefont}{\hfill\Large\bfseries}
\renewcommand{\cftaftertoctitle}{\hfill}
\setlength{\cftbeforetoctitleskip}{0pt}
\setlength{\cftaftertoctitleskip}{1em}
\setlength{\cftsecindent}{0em}
\setlength{\cftsubsecindent}{1.5em}
\setlength{\cftbeforesecskip}{0.5em}
\setlength{\cftbeforesubsecskip}{0.3em}

%---------------------------------
% Заголовки разделов
\titleformat{\section}
{\normalfont\Large\bfseries}{\thesection.}{1em}{}

\titleformat{\subsection}
{\normalfont\large\bfseries}{\thesubsection.}{1em}{}

%---------------------------------
% Стиль для доказательств, замечаний и следствий
\newcommand{\myproof}{\textbf{Доказательство.}\ }
\newcommand{\myremark}{\textbf{Замечание.}\ }
\newcommand{\mycorollary}{\textbf{Следствие.}\ }

%---------------------------------
% Цветные боксы для разных типов информации
\newtcolorbox{definitionbox}[2][]{colback=blue!5, colframe=blue!60!black, fonttitle=\bfseries, title=Определение #2, #1}
\newtcolorbox{theorembox}[2][]{colback=red!5, colframe=red!60!black, fonttitle=\bfseries, title=Теорема #2, #1}
\newtcolorbox{lemmabox}[2][]{colback=orange!5, colframe=orange!60!black, fonttitle=\bfseries, title=Лемма #2, #1}
\newtcolorbox{corollarybox}[2][]{colback=green!5, colframe=green!60!black, fonttitle=\bfseries, title=Следствие #2, #1}
\newtcolorbox{examplebox}[2][]{colback=gray!5, colframe=gray!50!black, fonttitle=\bfseries, title=Пример #2, #1}
\newtcolorbox{remarkbox}[2][]{colback=yellow!5, colframe=yellow!60!black, fonttitle=\bfseries, title=Замечание #2, #1}
\newtcolorbox{notebook}[2][]{colback=purple!5, colframe=purple!60!black, fonttitle=\bfseries, title=Заметочка #2, #1}
\newtcolorbox{propertybox}[2][]{colback=cyan!5, colframe=cyan!60!black, fonttitle=\bfseries, title=Свойство #2, #1}
\newtcolorbox{codebox}[2][]{colback=gray!10, colframe=gray!80!black, fonttitle=\bfseries, title=Листинг #2, #1}
\newtcolorbox{algorithmbox}[2][]{colback=green!5, colframe=green!60!black, fonttitle=\bfseries, title=Алгоритм #2, #1}


\newtcolorbox{importantbox}[2][]{colback=red!10, colframe=red!80!black, fonttitle=\bfseries, title=Важно! #2, #1}

%---------------------------------
% Колонтитулы
\pagestyle{fancy}
\fancyhf{}
\renewcommand{\headrulewidth}{0pt}
\fancyfoot[L]{\footnotesize Конспект по дискретной математике}
\fancyfoot[R]{\footnotesize \thepage\ / \pageref{LastPage}}

%---------------------------------
% Типографика
\sloppy

%-------------------------------------------------
\begin{document}
	
	% Титульная страница
	\begin{titlepage}
		\centering
		{\LARGE\bfseries Подготовка к экзамену по дискретной математике\par}
		\vspace{0.5cm}
		{\Large ПИ СПБГУ\par}
		\vspace{0.3cm}
		{\Large Семестр 1\par}
		
		\vspace{2cm}
		
		{\large Январь 2026\par}
		
		\vfill
		
		{\footnotesize составлен по учебнику Шмидта и др.\par}
	\end{titlepage}
	
	\newpage
	\tableofcontents
	
	
	
	
	
% ==================== РАЗДЕЛ 1: ТЕОРИЯ МНОЖЕСТВ ====================
\newpage
\section{Множества и операции над ними}

\subsection{Множества, способы задания и операции над ними. Разбиение множеств. Произведение разбиений, теорема о его существовании. Принцип математической индукции. Мощность множества всех подмножеств. Прямое произведение множеств (1 вопрос)}

\begin{definitionbox}{1.1: Множество}
	\textbf{Множеством} называется неопределяемое понятие, коллекция объектов произвольной природы, рассматриваемая как единое целое. Для обозначения используются прописные латинские или греческие буквы. Объекты, составляющие множество, называются его \textbf{элементами}.
\end{definitionbox}

\noindent \textbf{Обозначения и примеры:}
\begin{itemize}
	\item $a \in A$ — элемент $a$ \textit{содержится} в множестве $A$ (является его элементом); $a \notin A$ — элемент $a$ \textit{не содержится} в $A$.
	\item \textbf{Пример:} $\Phi = \{1, \lambda, \text{red}, \{IV, \text{white}\}\}$. Здесь $\{IV, \text{white}\} \in \Phi$, но $IV \notin \Phi$ (так как $IV$ не является элементом $\Phi$, хотя и содержится внутри элемента $\{IV, \text{white}\}$).
	\item $|A|$ — \textbf{мощность} (кардинальное число) множества $A$. Для конечного множества это число его элементов.
	\item $\varnothing$ — \textbf{пустое множество} ($|\varnothing| = 0$), не содержащее ни одного элемента.
\end{itemize}

\begin{definitionbox}{1.2: Подмножество}
	Множество $B$ называется \textbf{подмножеством} множества $A$ ($B \subseteq A$), если каждый элемент множества $B$ является также элементом множества $A$:
	\[
	B \subseteq A \ \Longleftrightarrow\ \forall x \, (x \in B \Rightarrow x \in A).
	\]
	Если при этом $B \neq A$ (т.е. $\exists x \in A : x \notin B$), то $B$ называется \textbf{собственным подмножеством} ($B \subset A$ или $B \subsetneqq A$).
	
	\textbf{Замечание:} Пустое множество $\varnothing$ является подмножеством любого множества. Любое множество является подмножеством самого себя.
\end{definitionbox}

\begin{center}
	\begin{tikzpicture}
		\draw[thick, blue] (0,0) circle (2cm) node[above] {$A$};
		\draw[thick, red] (-0.5,0) circle (1cm) node[above] {$B$};
		\node at (0,-2.5) {$B \subseteq A$};
	\end{tikzpicture}
	\captionof{figure}{Подмножество $B \subseteq A$}
\end{center}

\subsubsection{Способы задания множеств}
Существует несколько основных способов описать множество:
\begin{enumerate}
	\item \textbf{Перечислением (явным списком):} Все элементы множества записываются в фигурных скобках через запятую. Например, $A = \{1, 2, 3, 4, 5\}$, $B = \{\text{красный}, \text{зелёный}, \text{синий}\}$.
	\item \textbf{Интуитивное (с многоточием):} Используется, когда множество содержит много элементов, но закономерность их образования очевидна. Например, $\{1, 2, \dots, 10\}$, $\{2, 4, 6, \dots\}$ — множество четных положительных чисел.
	\item \textbf{Описанием характеристического свойства (условием выбора):} Задается как $\{x \in U : P(x)\}$, где $U$ — универсальное множество (или множество, из которого производится выбор), а $P(x)$ — некоторое условие (предикат), которому должны удовлетворять элементы. Например, $\{x \in \mathbb{N} : x \text{ делится на } 2\}$ — множество четных натуральных чисел.
	\item \textbf{Порождающей процедурой:} Множество задается правилом, порождающим его элементы. Например, множество всех строк, которые можно получить из заданного алфавита.
	\item \textbf{Результатом операций:} Новое множество задается как результат применения операций (объединение, пересечение и т.д.) к уже известным множествам.
\end{enumerate}

\subsubsection{Операции над множествами}
Пусть даны множества $A$ и $B$.
\begin{itemize}
	\item \textbf{Объединение:} $A \cup B = \{x : x \in A \lor x \in B\}$ (элементы, принадлежащие хотя бы одному из множеств).
	\item \textbf{Пересечение:} $A \cap B = \{x : x \in A \land x \in B\}$ (элементы, принадлежащие обоим множествам одновременно).
	\item \textbf{Разность:} $A \setminus B = \{x : x \in A \land x \notin B\}$ (элементы $A$, не принадлежащие $B$).
	\item \textbf{Симметрическая разность:} $A \triangle B = (A \setminus B) \cup (B \setminus A)$ (элементы, принадлежащие ровно одному из множеств).
\end{itemize}

\begin{center}
	\begin{tikzpicture}
		% Объединение
		\begin{scope}[shift={(0,2.5)}]
			\draw[thick, blue] (0,0) circle (1.5cm);
			\draw[thick, blue] (2,0) circle (1.5cm);
			\node at (1,0) {$A \cup B$};
			\node at (-0.8,0.8) {$A$};
			\node at (2.8,0.8) {$B$};
		\end{scope}
		
		% Пересечение
		\begin{scope}[shift={(5,2.5)}]
			\draw[thick, blue] (0,0) circle (1.5cm);
			\draw[thick, blue] (2,0) circle (1.5cm);
			\begin{scope}
				\clip (0,0) circle (1.5cm);
				\fill[red, opacity=0.5] (2,0) circle (1.5cm);
			\end{scope}
			\node at (1,0) {$A \cap B$};
			\node at (-0.8,0.8) {$A$};
			\node at (2.8,0.8) {$B$};
		\end{scope}
		
		% Разность A\B
		\begin{scope}[shift={(0,-2)}]
			\draw[thick, blue] (0,0) circle (1.5cm);
			\draw[thick, blue] (2,0) circle (1.5cm);
			\begin{scope}
				\clip (0,0) circle (1.5cm);
				\fill[red, opacity=0.5] (0,0) circle (1.5cm);
				\clip (2,0) circle (1.5cm);
				\fill[white] (2,0) circle (1.5cm);
			\end{scope}
			\node at (0.5,0) {$A \setminus B$};
			\node at (-0.8,0.8) {$A$};
			\node at (2.8,0.8) {$B$};
		\end{scope}
		
		% Симметрическая разность
		\begin{scope}[shift={(5,-2)}]
			\draw[thick, blue] (0,0) circle (1.5cm);
			\draw[thick, blue] (2,0) circle (1.5cm);
			\begin{scope}
				\fill[red, opacity=0.5] (0,0) circle (1.5cm);
				\fill[red, opacity=0.5] (2,0) circle (1.5cm);
			\end{scope}
			\begin{scope}
				\clip (0,0) circle (1.5cm);
				\fill[white] (2,0) circle (1.5cm);
			\end{scope}
			\node at (1,-0.5) {$A \triangle B$};
			\node at (-0.8,0.8) {$A$};
			\node at (2.8,0.8) {$B$};
		\end{scope}
	\end{tikzpicture}
	\captionof{figure}{Операции над множествами (диаграммы Эйлера-Венна)}
\end{center}

\begin{propertybox}{1.3: Свойства объединения и пересечения}
	Операции \textbf{объединения} ($\cup$) и \textbf{пересечения} ($\cap$) множеств обладают свойствами:
	\begin{itemize}
		\item \textbf{Коммутативность:} $A \cup B = B \cup A$, $A \cap B = B \cap A$.
		\item \textbf{Ассоциативность:} $(A \cup B) \cup C = A \cup (B \cup C)$, $(A \cap B) \cap C = A \cap (B \cap C)$.
		\item \textbf{Дистрибутивность:} $A \cup (B \cap C) = (A \cup B) \cap (A \cup C)$, $A \cap (B \cup C) = (A \cap B) \cup (A \cap C)$.
	\end{itemize}
\end{propertybox}

\noindent Для удобства записи используются обозначения:
\[
A_1 \cup A_2 \cup \dots \cup A_n = \bigcup_{i=1}^{n} A_i, \quad A_1 \cap A_2 \cap \dots \cap A_n = \bigcap_{i=1}^{n} A_i
\]

\begin{definitionbox}{1.4: Множество всех подмножеств (Булеан)}
	Пусть $A$ — произвольное множество. \textbf{Множеством всех подмножеств} (или \textbf{булеаном}) множества $A$ называется множество, элементами которого являются все возможные подмножества множества $A$:
	\[ \mathcal{P}(A) := \{ X : X \subseteq A \} \]
	Также используется обозначение $2^A$.
\end{definitionbox}

\begin{examplebox}{1.5: Примеры с булеаном}
	Рассмотрим множество $\Phi = \{1, \lambda, \text{red}, \{IV, \text{white}\}\}$.
	\begin{itemize}
		\item Утверждение ``$\{red\} \subseteq 2^\Phi$'' — \textbf{ложно}, так как $\{red\}$ — это множество, содержащее элемент $red$. Чтобы оно было подмножеством $2^\Phi$, каждый его элемент ($red$) должен быть элементом $2^\Phi$, т.е. $red$ должен быть подмножеством $\Phi$. Но $red$ не является множеством (в данном контексте) или не является подмножеством $\Phi$.
		\item Утверждение ``$\{red\} \in 2^\Phi$'' — \textbf{истинно}, так как $\{red\}$ — это множество, содержащее элемент $red$, и $\{red\} \subseteq \Phi$, значит, оно является элементом булеана.
		\item $\varnothing \in 2^\Phi$, так как $\varnothing$ — подмножество любого множества.
		\item $\{1, \{IV, \text{white}\}\} \in 2^\Phi$, так как это подмножество $\Phi$.
	\end{itemize}
	\textbf{Еще пример:} Для множества $A = \{a, b\}$ булеан $2^A = \{\varnothing, \{a\}, \{b\}, \{a, b\}\}$, его мощность $|2^A| = 4$.
\end{examplebox}

\begin{theorembox}{1.6: Мощность булеана}
	Для любого конечного множества $A$ мощность его булеана равна $2^{|A|}$: $|2^A| = 2^{|A|}$.
\end{theorembox}

\subsection{Разбиение множеств}

\begin{definitionbox}{1.7: Разбиение множества}
	Пусть $A$ — непустое множество. Набор (семейство) $\Lambda = \{\Lambda_1, \Lambda_2, \dots, \Lambda_n\}$ называется \textbf{разбиением} множества $A$, если выполняются следующие условия:
	\begin{enumerate}
		\item $\forall i \ \Lambda_i \subseteq A$ (каждый элемент разбиения — подмножество $A$);
		\item $\forall i\ \Lambda_i \neq \varnothing$ (все блоки разбиения непусты);
		\item $\forall i \neq j\ \Lambda_i \cap \Lambda_j = \varnothing$ (блоки попарно не пересекаются, т.е. дизъюнктны);
		\item $\bigcup_{i=1}^n \Lambda_i = A$ (объединение всех блоков дает исходное множество $A$).
	\end{enumerate}
\end{definitionbox}

\noindent \textbf{Пример:} Для $A = \{1,2,3,4,5\}$ разбиением будет $\Lambda = \{\{1,2\}, \{3,5\}, \{4\}\}$.

\begin{definitionbox}{1.8: Измельчение разбиения}
	Пусть $\Lambda = \{\Lambda_1, \dots, \Lambda_m\}$ и $K = \{K_1, \dots, K_n\}$ — два разбиения одного и того же множества $A$.
	Говорят, что разбиение $\Lambda$ является \textbf{измельчением} разбиения $K$ (или $\Lambda$ мельче $K$), если каждый блок разбиения $\Lambda$ целиком содержится в некотором блоке разбиения $K$:
	\[
	\forall i \in \{1,\dots,m\}\ \exists j \in \{1,\dots,n\} : \Lambda_i \subseteq K_j.
	\]
	\textbf{Замечание:} Любое разбиение является измельчением самого себя. Интуитивно, измельчение получается путем дальнейшего дробления блоков исходного разбиения на более мелкие части.
\end{definitionbox}

\noindent \textbf{Пример:} Для множества $A = \{1,2,3,4\}$ разбиение $\Lambda = \{\{1\}, \{2\}, \{3,4\}\}$ является измельчением разбиения $K = \{\{1,2\}, \{3,4\}\}$, так как $\{1\} \subseteq \{1,2\}$, $\{2\} \subseteq \{1,2\}$, $\{3,4\} \subseteq \{3,4\}$.

\subsubsection*{Разбиение в прообразах функций}
Важный способ получения разбиения связан с функциями. Пусть $f: M \to R$ — функция из множества $M$ в множество $R$. Рассмотрим для каждого значения $r \in R$ его \textbf{полный прообраз}:
\[
f^{-1}(r) = \{ x \in M : f(x) = r \}.
\]
Тогда набор непустых полных прообразов $\mathcal{A}_f = \{ f^{-1}(r) : r \in R, f^{-1}(r) \neq \varnothing \}$ образует \textbf{разбиение} множества $M$. Элементы $M$ попадают в один блок тогда и только тогда, когда их значения под действием $f$ совпадают.


\begin{definitionbox}{1.9: Произведение разбиений}
	Пусть $\Lambda$ и $K$ — два разбиения множества $A$. Их \textbf{произведением} (или пересечением) называется такое разбиение $\Pi$ множества $A$, которое является \textbf{наиболее крупным} (т.е. содержащим самые большие возможные блоки) среди всех разбиений, которые являются измельчениями одновременно и $\Lambda$, и $K$. Другими словами, $\Pi$ — это "наибольшее общее измельчение" разбиений $\Lambda$ и $K$.
\end{definitionbox}

\begin{theorembox}{1.10: О существовании произведения разбиений}
	Для любых двух разбиений $\Lambda$ и $K$ множества $A$ их произведение $\Pi$ существует и состоит из всех возможных непустых пересечений блоков этих разбиений:
	\[ \Pi = \{ \Lambda_i \cap K_j \mid \Lambda_i \in \Lambda, K_j \in K,\ \Lambda_i \cap K_j \neq \varnothing \} \]
\end{theorembox}

\myproof
Пусть даны два разбиения $\Lambda = \{\Lambda_1, \dots, \Lambda_m\}$ и $K = \{K_1, \dots, K_n\}$ множества $A$. Рассмотрим семейство множеств $\Pi = \{\Lambda_i \cap K_j : i = 1..m, j = 1..n,\ \Lambda_i \cap K_j \neq \varnothing\}$.

Докажем, что $\Pi$ — разбиение (т.е. удовлетворяет пунктам из определения 1.7):
\begin{enumerate}
	\item \textbf{Непустота блоков:} По построению мы включаем в $\Pi$ только непустые пересечения.
	\item \textbf{Попарная дизъюнктность:} Предположим, что два различных элемента $\Pi$, например $X = \Lambda_i \cap K_j$ и $Y = \Lambda_p \cap K_q$, пересекаются. Тогда существует элемент $x \in X \cap Y$.
	Так как $x \in \Lambda_i \cap K_j$, то $x \in \Lambda_i$ и $x \in K_j$.
	Так как $x \in \Lambda_p \cap K_q$, то $x \in \Lambda_p$ и $x \in K_q$.
	Поскольку $\Lambda$ — разбиение, его блоки не пересекаются, поэтому из $x \in \Lambda_i \cap \Lambda_p$ следует, что $i = p$.
	Аналогично, поскольку $K$ — разбиение, из $x \in K_j \cap K_q$ следует, что $j = q$.
	Значит, $X$ и $Y$ — это одно и то же пересечение. Противоречие с предположением, что они различны. Следовательно, любые два различных элемента $\Pi$ не пересекаются.
	\item \textbf{Полнота покрытия:} Возьмем произвольный элемент $x \in A$. Так как $\Lambda$ — разбиение, $x$ принадлежит ровно одному его блоку, скажем $\Lambda_i$. Так как $K$ — разбиение, $x$ принадлежит ровно одному его блоку, скажем $K_j$. Тогда $x$ принадлежит их пересечению $\Lambda_i \cap K_j$, которое непусто (так как содержит $x$) и, следовательно, является элементом $\Pi$. Таким образом, каждый элемент $A$ покрывается одним из блоков $\Pi$.
\end{enumerate}
Мы доказали, что $\Pi$ — разбиение. Очевидно, что $\Pi$ является измельчением $\Lambda$ (так как каждый блок $\Pi$ есть подмножество некоторого $\Lambda_i$) и измельчением $K$. Пусть $\Pi'$ — любое другое общее измельчение $\Lambda$ и $K$. Тогда любой блок $\Pi'$ обязан целиком лежать в пересечении какого-то блока из $\Lambda$ и какого-то блока из $K$, то есть быть подмножеством некоторого $\Lambda_i \cap K_j$. Следовательно, $\Pi'$ является измельчением $\Pi$, что делает $\Pi$ самым крупным общим измельчением. \qed

\begin{examplebox}{1.11: Пример произведения разбиений}
	Пусть $A = \{0,1,\dots,11\}$. Определим два разбиения $A$:
	\begin{itemize}
		\item $\Lambda$ разбивает $A$ по остатку от деления на 3 (условно):
		$\Lambda_1 = \{0,1,2,3\}$,
		$\Lambda_2 = \{4,5,6,7,8,9\}$,
		$\Lambda_3 = \{10,11\}$.
		\item $K$ разбивает $A$ на две половины:
		$K_1 = \{0,1,2,3,4,5\}$,
		$K_2 = \{6,7,8,9,10,11\}$.
	\end{itemize}
	Найдем произведение $\Pi$, взяв все возможные непустые пересечения:
	\begin{itemize}
		\item $\Lambda_1 \cap K_1 = \{0,1,2,3\}$ (непусто)
		\item $\Lambda_1 \cap K_2 = \{4,5\}?$ Нет, $\Lambda_1$ не пересекается с $K_2$, т.к. $\Lambda_1 = \{0,1,2,3\}$, а $K_2 = \{6,\dots,11\}$.
		\item $\Lambda_2 \cap K_1 = \{4,5\}$ (непусто)
		\item $\Lambda_2 \cap K_2 = \{6,7,8,9\}$ (непусто)
		\item $\Lambda_3 \cap K_1 = \varnothing$
		\item $\Lambda_3 \cap K_2 = \{10,11\}$ (непусто)
	\end{itemize}
	Таким образом, произведение:
	\[ \Pi = \{ \{0,1,2,3\}, \{4,5\}, \{6,7,8,9\}, \{10,11\} \} \]
	Видно, что $\Pi$ действительно мельче и $\Lambda$, и $K$.
\end{examplebox}

\subsection{Принцип математической индукции}

Математическая индукция — метод математического доказательства, который используется, чтобы доказать истинность некоторого утверждения для всех натуральных чисел (т.е. имеющих вид «Для любого натурального $n$ справедливо свойство $P(n)$»).

\begin{theorembox}{1.12: Принцип математической индукции}
	Пусть $P(n)$ — некоторое утверждение, зависящее от натурального числа $n$. Если выполняются два условия:
	\begin{enumerate}
		\item \textbf{База индукции:} Утверждение $P(1)$ истинно.
		\item \textbf{Шаг индукции:} Для любого натурального $k$ из предположения, что утверждения $P(1), P(2), \dots, P(k)$ истинны (или, в простейшей форме, из истинности $P(k)$), следует истинность $P(k+1)$.
	\end{enumerate}
	то утверждение $P(n)$ истинно для всех натуральных чисел $n$ (т.е. множество $\Sigma = \{ n \in \mathbb{N} \mid P(n) \text{ истинно} \}$ совпадает с $\mathbb{N}$).
\end{theorembox}

\noindent \textbf{Пояснение на примере:} Докажем, что $1 + 2 + \dots + n = \frac{n(n+1)}{2}$.
\begin{itemize}
	\item \textbf{База ($n=1$):} $1 = \frac{1\cdot2}{2} = 1$. Верно.
	\item \textbf{Шаг:} Предположим, что формула верна для $n=k$, т.е. $1+2+\dots+k = \frac{k(k+1)}{2}$.
	Докажем для $n=k+1$:
	$1+2+\dots+k+(k+1) = (1+2+\dots+k) + (k+1) = \frac{k(k+1)}{2} + (k+1) = \frac{k(k+1) + 2(k+1)}{2} = \frac{(k+1)(k+2)}{2}$.
	Что и требовалось доказать.
	\item \textbf{Вывод:} Формула верна для всех $n \in \mathbb{N}$.
\end{itemize}


\begin{propertybox}{Замечание (о базе индукции)}
	В некоторых случаях удобнее переопределять принцип математической индукции так, чтобы базой было не число $1$, а число $0$. В таком случае вместо множества натуральных чисел $\mathbb{N}$ используется множество $\mathbb{N}_0 = \{0,1,2,\dots\}$ (натуральные числа с нулём)\\
		
	\textbf{Пример:} При доказательстве свойств, связанных с мощностью конечных множеств (как в теореме 1.13), база $n=0$ (пустое множество) часто оказывается естественнее, чем $n=1$.
\end{propertybox}


\subsection{Мощность множества всех подмножеств (доказательство)}

\begin{theorembox}{1.13: Мощность булеана (доказательство)}
	Для любого конечного множества $A$ верно $|2^A| = 2^{|A|}$.
\end{theorembox}

\myproof (Методом математической индукции по $n = |A|$):
\begin{itemize}
	\item \textbf{База:} $n=0$. Тогда $A = \varnothing$. Булеан пустого множества содержит только само пустое множество: $2^\varnothing = \{\varnothing\}$. Следовательно, $|2^\varnothing| = 1 = 2^0$. База доказана.
	\item \textbf{Шаг индукции:} Предположим, что утверждение верно для любого множества мощности $n$, т.е. $|2^X| = 2^n$, если $|X| = n$ (гипотеза индукции).
	Рассмотрим произвольное множество $A$ такое, что $|A| = n+1$. Выберем и зафиксируем в нем некоторый элемент $a \in A$.
	Разобьем все подмножества $A$ на два непересекающихся класса:
	\begin{enumerate}
		\item Подмножества, \textbf{не содержащие} элемент $a$. Это в точности все подмножества множества $A \setminus \{a\}$, мощность которого равна $n$. По гипотезе индукции, количество таких подмножеств равно $2^n$.
		\item Подмножества, \textbf{содержащие} элемент $a$. Каждое такое подмножество можно единственным образом представить как $Y \cup \{a\}$, где $Y$ — некоторое подмножество множества $A \setminus \{a\}$ (т.е. $Y \subseteq A \setminus \{a\}$). Количество различных $Y$ равно $|2^{A \setminus \{a\}}| = 2^n$. Следовательно, количество подмножеств, содержащих $a$, также равно $2^n$.
	\end{enumerate}
	Эти два класса не пересекаются и в сумме дают все подмножества $A$. Поэтому общее число подмножеств:
	\[ |2^A| = 2^n + 2^n = 2 \cdot 2^n = 2^{n+1}. \]
	Шаг индукции доказан.
\end{itemize}
Таким образом, согласно принципу математической индукции, формула $|2^A| = 2^{|A|}$ верна для всех конечных множеств $A$. \qed

\subsection{Прямое (декартово) произведение}

Прямое, или декартово произведение двух множеств — множество, элементами которого являются все возможные упорядоченные пары элементов исходных множеств.


\begin{propertybox}{Обобщение прямого произведения}
	Понятие прямого произведения естественно обобщается на произведение множеств, наделённых дополнительной структурой (алгебраической, топологической, порядковой и так далее). Важно, что прямое произведение часто \textbf{наследует} структуры исходных множеств, позволяя конструировать новые объекты с уже знакомыми свойствами.\\
	
	\textbf{Примеры:}
	\begin{itemize}
		\item Если $(G,\cdot)$ и $(H,\star)$ — группы, то на $G \times H$ можно ввести групповую операцию покомпонентно:
		\[
		(g_1,h_1) \ast (g_2,h_2) = (g_1 \cdot g_2,\; h_1 \star h_2).
		\]
		Получается новая группа — \textit{прямое произведение групп}.
		
		\item Если $(X,\tau_X)$ и $(Y,\tau_Y)$ — топологические пространства, то на $X \times Y$ можно ввести \textit{топологию произведения}. Её базой являются произведения открытых множеств $U \times V$, где $U \in \tau_X$, $V \in \tau_Y$.
		
		\item Для упорядоченных множеств на произведении вводят \textit{порядок произведения}: покомпонентный или лексикографический.
	\end{itemize}
\end{propertybox}


\begin{definitionbox}{1.14: Упорядоченная пара}
	Пусть $a \in A$, $b \in B$. \textbf{Упорядоченной парой} называется объект $(a, b)$, существенным свойством которого является то, что порядок следования элементов важен.
	\[ (a_1, b_1) = (a_2, b_2) \ \Longleftrightarrow\ a_1 = a_2 \ \text{и} \ b_1 = b_2 \]
\end{definitionbox}

\begin{definitionbox}{1.15: Кортеж (вектор)}
	Обобщением понятия пары является \textbf{упорядоченный кортеж} (или вектор) длины $n$: $(a_1, a_2, \dots, a_n)$, где $a_i \in A_i$.
	Два кортежа равны тогда и только тогда, когда они имеют одинаковую длину и соответствующие компоненты равны:
	\[ (a_1, \dots, a_n) = (b_1, \dots, b_m) \ \Longleftrightarrow\ n = m \ \text{и} \ \forall i: a_i = b_i \]
\end{definitionbox}

\begin{definitionbox}{1.16: Прямое произведение}
	\textbf{Прямым (декартовым) произведением} множеств $A$ и $B$ называется множество всех упорядоченных пар, в которых первый элемент принадлежит $A$, а второй — $B$:
	\[ A \times B := \{ (a, b) \mid a \in A,\ b \in B \} \]
	В общем случае для $n$ множеств:
	\[ A_1 \times A_2 \times \dots \times A_n := \{ (a_1, a_2, \dots, a_n) \mid a_i \in A_i \} \]
	Если все множества одинаковы ($A_1 = A_2 = \dots = A_n = A$), то используется обозначение степени:
	\[ A^n = A \times A \times \dots \times A \ (\text{$n$ раз}) \]
\end{definitionbox}

\begin{examplebox}{1.17: Примеры}
	\begin{itemize}
		\item Пусть $A = \{1,2\}$, $B = \{x,y\}$. Тогда
		$A \times B = \{(1,x), (1,y), (2,x), (2,y)\}$.
		$B \times A = \{(x,1), (x,2), (y,1), (y,2)\}$.
		Видно, что в общем случае $A \times B \neq B \times A$.
		\item Пусть $B = \{0, 1\}$ (двоичное множество). Тогда $B^m$ состоит из всех возможных двоичных кортежей (строк) длины $m$. Так как на каждой из $m$ позиций может стоять 0 или 1, то $|B^m| = 2^m$.
		\item $\mathbb{R}^2 = \mathbb{R} \times \mathbb{R}$ — это всем известная декартова плоскость, где каждая точка задается парой координат $(x, y)$.
	\end{itemize}
\end{examplebox}

\begin{center}
	\begin{tikzpicture}[scale=0.8]
		\draw[->] (-1,0) -- (4,0) node[right] {$A$};
		\draw[->] (0,-1) -- (0,4) node[above] {$B$};
		
		\foreach \x in {1,2,3} {
			\draw[gray, dashed] (\x,0) -- (\x,3);
		}
		\foreach \y in {1,2,3} {
			\draw[gray, dashed] (0,\y) -- (3,\y);
		}
		
		\foreach \x in {1,2,3} {
			\foreach \y in {1,2,3} {
				\fill[blue] (\x,\y) circle (2pt);
			}
		}
		
		\node at (1,-0.3) {$a_1$};
		\node at (2,-0.3) {$a_2$};
		\node at (3,-0.3) {$a_3$};
		\node at (-0.3,1) {$b_1$};
		\node at (-0.3,2) {$b_2$};
		\node at (-0.3,3) {$b_3$};
		
		\node at (1.5,1.5) {$(a_i, b_j)$};
	\end{tikzpicture}
	\captionof{figure}{Иллюстрация декартова произведения конечных множеств $A \times B$. Каждая точка с целочисленными координатами соответствует упорядоченной паре.}
\end{center}
	
	
	
	
	
	
	
	
	
	
	
	
	% ==================== РАЗДЕЛ 2: БИНАРНЫЕ ОТНОШЕНИЯ ====================
	\newpage
	\section{Бинарные отношения и их свойства}
	
	\subsection{Бинарные отношения, свойства бинарных отношений. Отношение порядка. Частичный и линейный порядки, отношение строгого порядка (2 вопрос)}
	
	\begin{definitionbox}{2.1: Бинарное отношение}
		\textbf{Бинарным (двуместным) отношением} $R$ на непустых и конечных множествах $A$ и $B$ называется любое подмножество их прямого произведения: $R \subseteq A \times B$.
		
		Если $A = B$, говорят о \textbf{бинарном отношении на множестве $A$}: $R \subseteq A \times A$.
	\end{definitionbox}
	
	\noindent Для записи используют два формата:
	\begin{itemize}
		\item $(x, y) \in R$ — стандартная запись;
		\item $xRy$ — инфиксная запись.
	\end{itemize}
	
	\begin{examplebox}{2.2: Примеры бинарных отношений}
		\begin{itemize}
			\item $A = \{1,2,3,4\}$, $P = \{(1,2), (1,3), (2,4), (4,4)\}$
			\item $A = \{\text{яблоко, груша, персик, банан, киви}\}$, \\
			$P = \{(\text{яблоко, персик}), (\text{банан, киви}), (\text{груша, киви})\}$
			\item $R = \{(a,b) \in \mathbb{N} \times \mathbb{N} \mid a \le b\}$ (отношение «меньше или равно»)
			\item $R = \{(a,b) \in \mathbb{N} \times \mathbb{N} \mid a \vdots b\}$ (отношение «делится на»)
		\end{itemize}
	\end{examplebox}
	
	\myremark{Важно:} Чаще отношения задаются правилом: $R = \{(x,y) \mid x + y > 2\}$.
	
	\subsection{Свойства бинарных отношений}
	
	Пусть $R$ — бинарное отношение на множестве $A$.
	
	\begin{definitionbox}{2.3: Рефлексивность и антирефлексивность}
		\begin{itemize}
			\item $R$ \textbf{рефлексивно}, если $\forall a \in A: (a,a) \in R$ (пример: $\ge$)
			\item $R$ \textbf{антирефлексивно}, если $\forall a \in A: (a,a) \notin R$ (пример: $>$)
		\end{itemize}
	\end{definitionbox}
	
	\begin{definitionbox}{2.4: Симметричность...}
		\begin{itemize}
			\item $R$ \textbf{симметрично}, если $(a,b) \in R \Rightarrow (b,a) \in R$ (пример: $=$)
			\item $R$ \textbf{антисимметрично}, если $(a,b) \in R \land (b,a) \in R \Rightarrow a = b$ (пример: $\le$)
			\item $R$ \textbf{асимметрично}, если $(a,b) \in R \Rightarrow (b,a) \notin R$ (пример: $>$)
		\end{itemize}
	\end{definitionbox}
	
	\begin{definitionbox}{2.5: Транзитивность}
		$R$ \textbf{транзитивно}, если $(a,b) \in R \land (b,c) \in R \Rightarrow (a,c) \in R$ (примеры: $>$, $<$, $\le$)
	\end{definitionbox}
	
	\begin{center}
		\begin{tabular}{|l|l|}
			\hline
			\textbf{Свойство} & \textbf{Определение} \\ \hline
			Рефлексивность & $\forall a \in A: (a,a) \in R$ \\ \hline
			Антирефлексивность & $\forall a \in A: (a,a) \notin R$ \\ \hline
			Симметричность & $(a,b) \in R \Rightarrow (b,a) \in R$ \\ \hline
			Антисимметричность & $(a,b) \in R \land (b,a) \in R \Rightarrow a = b$ \\ \hline
			Асимметричность & $(a,b) \in R \Rightarrow (b,a) \notin R$ \\ \hline
			Транзитивность & $(a,b) \in R \land (b,c) \in R \Rightarrow (a,c) \in R$ \\ \hline
		\end{tabular}
	\end{center}
	
	
	
\begin{definitionbox}{2.6: Свойства бинарных отношений (на пальцах)}
	Представь, что у нас есть множество $A$ (например, группа людей или набор чисел), а отношение $R$ — это некое правило, по которому мы их связываем.
	
	\textbf{1. Рефлексивность и антирефлексивность}
	\begin{itemize}
		\item \textbf{Рефлексивность:} Каждый элемент связан сам с собой.
		\begin{itemize}
			\item \textit{Суть:} Если ты посмотришь в зеркало, ты увидишь себя.
			\item \textit{Пример:} Отношение $\le$ на числах. $5 \le 5$ — правда.
			\item \textit{На графе:} У каждой вершины есть «петля».
		\end{itemize}
		
		\item \textbf{Антирефлексивность:} Ни один элемент не связан сам с собой.
		\begin{itemize}
			\item \textit{Суть:} Запрет на связь с самим собой.
			\item \textit{Пример:} Отношение «быть строго выше ростом». Ты не можешь быть выше себя.
			\item \textit{На графе:} Ни у одной вершины нет петли.
		\end{itemize}
	\end{itemize}
	
	\textbf{2. Симметричность, антисимметричность и асимметричность}
	\begin{itemize}
		\item \textbf{Симметричность:} Если $a$ связано с $b$, то и $b$ связано с $a$.
		\begin{itemize}
			\item \textit{Суть:} Полная взаимность.
			\item \textit{Пример:} Отношение «родственники». Если Иван — родственник Петра, то и Петр — родственник Ивана.
		\end{itemize}
		
		\item \textbf{Антисимметричность:} Если $a$ связано с $b$ \textbf{и} $b$ связано с $a$, то $a = b$.
		\begin{itemize}
			\item \textit{Суть:} Взаимность только при совпадении.
			\item \textit{Пример:} Отношение $\le$. Если $x \le y$ и $y \le x$, то $x = y$.
		\end{itemize}
		
		\item \textbf{Асимметричность:} Если $a$ связано с $b$, то $b$ не связано с $a$.
		\begin{itemize}
			\item \textit{Суть:} Игра в одни ворота.
			\item \textit{Пример:} Отношение «отец». Если А — отец Б, то Б не может быть отцом А.
		\end{itemize}
	\end{itemize}
	
	\textbf{3. Транзитивность}
	\begin{itemize}
		\item \textbf{Определение:} Если $a$ связано с $b$ и $b$ связано с $c$, то $a$ связано с $c$.
		\item \textit{Суть:} Правило домино.
		\item \textit{Примеры:}
		\begin{itemize}
			\item \textbf{Числа:} Если $a < b$ и $b < c$, то $a < c$.
			\item \textbf{Родословная:} Если А — предок Б, а Б — предок В, то А — предок В.
		\end{itemize}
		\item \textit{Контрпример:} Отношение «дружит с». Дружба не транзитивна.
	\end{itemize}
	
\end{definitionbox}
	
	\subsection{Отношения порядка}
	
	\begin{definitionbox}{2.7: Частичный порядок}
		Бинарное отношение $R$ над множеством А — \textbf{частичный порядок}, если оно:
		\begin{enumerate}
			\item рефлексивно;
			\item транзитивно;
			\item антисимметрично.
		\end{enumerate}
	\end{definitionbox}
	
	\begin{definitionbox}{2.8: Строгий частичный порядок}
		$R$ — \textbf{строгий частичный порядок}, если оно:
		\begin{enumerate}
			\item транзитивно;
			\item асимметрично.
		\end{enumerate}
	\end{definitionbox}
	
	\begin{definitionbox}{2.9:(Строгий) Линейный порядок}
		Частичный порядок $R$ называется \textbf{линейным}, если 
		$\forall a,b \in A, a \neq b: (a,b) \in R \lor (b,a) \in R$.
		
		\textbf{Строгий линейный порядок} — строгий частичный порядок с тем же свойством.
	\end{definitionbox}

	
	\begin{examplebox}{2.10: Примеры порядков}
		\begin{itemize}
			\item \textbf{Частичный, но не линейный:} Отношение ``делит'' на множестве $\{1,2,3,4,5,6\}$.
			\begin{itemize}
				\item Пример: $(2,6) \in R$ (2 делит 6), $(3,6) \in R$.
				\item Почему не линейный: 4 и 6 несравнимы (4 не делит 6, и 6 не делит 4).
			\end{itemize}
			\item \textbf{Линейный:} Отношение ``не больше'' ($\le$) на $\mathbb{N}$.
			\begin{itemize}
				\item Пример: $(3,5) \in R$, $(5,5) \in R$.
			\end{itemize}
			\item \textbf{Строгий линейный:} Отношение ``меньше'' ($<$) на $\mathbb{N}$.
			\begin{itemize}
				\item Пример: $(3,5) \in R$, но $(5,5) \notin R$.
			\end{itemize}
		\end{itemize}
	\end{examplebox}
	
	\begin{notebook}{2.11: Связь между порядками}
		\begin{itemize}
			\item Если $\preceq$ — частичный порядок (например, $\le$), то $a \prec b \iff (a \preceq b \land a \neq b)$ — строгий частичный порядок ($<$).
			\item Если $\prec$ — строгий частичный порядок (например, $<$), то $a \preceq b \iff (a \prec b \lor a = b)$ — частичный порядок ($\le$).
		\end{itemize}
	\end{notebook}




	
	
	
	
	
	
	
	
	
	
	
	% ==================== РАЗДЕЛ 3: ТОПОЛОГИЧЕСКАЯ СОРТИРОВКА ====================
	\newpage
	\section{Топологическая сортировка и экстремальные элементы}
	
	\subsection{Топологическая сортировка. Минимальные и максимальные элементы. Теорема о существовании топологической сортировки (3 вопрос)}
		
	\subsubsection{Топологическая сортировка}
	
	\begin{definitionbox}{3.9: Топологическая сортировка}
		Пусть $A$ — множество, частично упорядоченное отношением $R$ (строгим или нестрогим). \textbf{Топологической сортировкой} множества $A$ называется такой (строгий) линейный порядок $Q$ на $A$, что:
		\[ \forall a,b \in A: (a,b) \in R \Rightarrow (a,b) \in Q \]
		
		Иными словами, топологическая сортировка — это линейное упорядочение элементов множества, которое не нарушает исходный частичный порядок: если элемент $a$ предшествует $b$ в исходном порядке, то в результирующем линейном порядке $a$ должен стоять раньше $b$.
	\end{definitionbox}
	
	\myremark{Простыми словами:} Топологическая сортировка выстраивает элементы в очередь так, чтобы все зависимости (отношения порядка) были удовлетворены — предки всегда стоят раньше потомков.
	
	
	\subsubsection{Экстремальные элементы частично упорядоченных множеств}
		
	Пусть $(A, R)$ — частично упорядоченное множество (ЧУМ), где $R$ — отношение (строгого или нестрогого) частичного порядка.
	
	\begin{definitionbox}{3.1: Минимальный элемент}
		Элемент $m \in A$ называется \textbf{минимальным}, если в $A$ не существует элемента, меньшего чем $m$:
		\[ \forall x \in A: x \neq m \Rightarrow (x, m) \notin R \]
		(Для нестрогого порядка: $\forall x \in A: (x, m) \in R \Rightarrow x = m$).
		
		Иными словами, \textit{минимальный} — тот, меньше которого нет.
	\end{definitionbox}
	
	\begin{definitionbox}{3.2: Наименьший элемент}
		Элемент $m_0 \in A$ называется \textbf{наименьшим}, если он меньше любого другого элемента множества:
		\[ \forall x \in A: x \neq m_0 \Rightarrow (m_0, x) \in R \]
		
		Иными словами, \textit{наименьший} — тот, что меньше любого другого.
	\end{definitionbox}
	
	\begin{definitionbox}{3.3: Максимальный элемент}
		Элемент $M \in A$ называется \textbf{максимальным}, если в $A$ не существует элемента, большего чем $M$:
		\[ \forall x \in A: x \neq M \Rightarrow (M, x) \notin R \]
		
		Иными словами, \textit{максимальный} — тот, больше которого нет.
	\end{definitionbox}
	
	\begin{definitionbox}{3.4: Наибольший элемент}
		Элемент $M_0 \in A$ называется \textbf{наибольшим}, если он больше любого другого элемента множества:
		\[ \forall x \in A: x \neq M_0 \Rightarrow (x, M_0) \in R \]
		
		Иными словами, \textit{наибольший} — тот, что больше любого другого.
	\end{definitionbox}
	
	\begin{importantbox}{3.5: Важное различие}
		\begin{itemize}
			\item В частично упорядоченном множестве (ЧУМ) \textbf{минимальных} элементов может быть несколько, но \textbf{наименьших} — не более одного (аналогично с \textbf{максимальными} и \textbf{наибольшими}).
			\item В линейно упорядоченном множестве (ЛУМ) \textbf{минимальный} элемент будет лишь один, и он же будет \textbf{наименьшим} (аналогично с \textbf{максимальным} и \textbf{наибольшим}).
		\end{itemize}
	\end{importantbox}
	
	\begin{examplebox}{3.6: Пример экстремальных элементов}
		
		Рассмотрим частичный порядок на множестве $A = \{1,2,3,4\}$.
		
		\medskip
		\textbf{1. Частичный порядок.}
		
		\begin{center}
			\begin{tikzpicture}[>=stealth, scale=1]
				
				\node[draw, circle, fill=red!30] (1) at (0,2) {1};
				\node[draw, circle, fill=red!30] (2) at (0,0) {2};
				\node[draw, circle, fill=red!30] (3) at (2,1) {3};
				\node[draw, circle, fill=red!30] (4) at (4,1) {4};
				
				\draw[->] (1) -- (3);
				\draw[->] (1) -- (4);
				\draw[->] (2) -- (3);
				\draw[->] (2) -- (4);
				\draw[->] (3) -- (4);
				
			\end{tikzpicture}
		\end{center}
		
		Здесь:
		\begin{itemize}
			\item $1$ и $2$ — \textbf{минимальные}, но не наименьшие;
			\item $4$ — \textbf{максимальный} и \textbf{наибольший}.
		\end{itemize}
		
		\medskip
		\textbf{2. Топологическая сортировка.}
		
		Построим топологическую сортировку на $A$:
		
		\begin{center}
			\begin{tikzpicture}[>=stealth, scale=1]
				
				\node[draw, circle, fill=red!30] (1) at (0,0) {1};
				\node[draw, circle, fill=red!30] (2) at (2,1) {2};
				\node[draw, circle, fill=red!30] (3) at (4,1) {3};
				\node[draw, circle, fill=red!30] (4) at (6,0) {4};
				
				\draw[->] (1) -- (2);
				\draw[->] (1) -- (3);
				\draw[->] (1) -- (4);
				\draw[->] (2) -- (3);
				\draw[->] (2) -- (4);
				\draw[->] (3) -- (4);
				
			\end{tikzpicture}
		\end{center}
		
		В полученном линейном порядке:
		\begin{itemize}
			\item $1$ — минимальный и наименьший;
			\item $4$ — максимальный и наибольший.
		\end{itemize}
		
	\end{examplebox}
	
	\subsubsection{Лемма о существовании минимального элемента}
	
	\begin{lemmabox}{3.7: О существовании минимального/максимального элемента}
		Во всяком \textbf{конечном} непустом частично упорядоченном множестве $X$ существует минимальный (и, аналогично, максимальный) элемент.
	\end{lemmabox}
	
	\myproof (для минимального элемента, для максимального элемента доказательство аналогично) методом математической индукции по $|X|$.
	
	\textbf{База индукции:} $|X| = 1$, т.е. $X = \{x\}$. Очевидно, $x$ является минимальным элементом.
	
	\textbf{Шаг индукции:} Пусть утверждение верно для всех множеств мощности $n$. Рассмотрим множество $X$ с $|X| = n+1$, $n \geq 1$. Выберем произвольный элемент $x \in X$ и рассмотрим $X' = X \setminus \{x\}$. По предположению индукции в $X'$ существует минимальный элемент, обозначим его $x_0$.
	
	Рассмотрим два случая:
	\begin{itemize}
		\item Если $(x, x_0) \in R$ (в нестрогом порядке это означает $x \le x_0$), то $x$ является минимальным элементом $X$. Действительно, предположим, что существует $y \in X$, $y \neq x$, такой что $(y, x) \in R$. Тогда $y \in X'$ и $(y, x_0) \in R$ по транзитивности (так как $(y, x) \in R$ и $(x, x_0) \in R$), что противоречит минимальности $x_0$ в $X'$.
		\item Иначе $(x, x_0) \notin R$. Тогда $x_0$ остается минимальным элементом $X$, так как для любого $y \in X'$ $(y, x_0) \notin R$ по построению, а $(x, x_0) \notin R$ по условию.
	\end{itemize}
	
	Таким образом, в любом случае минимальный элемент существует. \qed
	
	\begin{corollarybox}{3.8}
		В конечном непустом ЧУМ всегда существует также и максимальный элемент (доказательство аналогично, с заменой порядка на обратный).
	\end{corollarybox}

	
	\subsubsection{Теорема о существовании топологической сортировки}
	
	\begin{theorembox}{3.10: О существовании топологической сортировки}
		Для любого \textbf{конечного} частично упорядоченного относительно R множества A $(A, R)$ существует топологическая сортировка.
	\end{theorembox}
	
	\myproof \textbf{(Конструктивное доказательство)}
	
	Обозначим $A_0 := A$. Построим последовательность множеств $A_i$ и элементов $m_i$ следующим итерационным процессом:
	
	\begin{enumerate}
		\item Если $A_i = \emptyset$, процесс завершен: топологическая сортировка $T_i = \emptyset$.
		\item Если $A_i \neq \emptyset$, то по лемме 3.7 в $A_i$ существует минимальный элемент (относительно исходного порядка $R$). Выберем один из них и обозначим его $m_i$.
		\item Определим $A_{i+1} := A_i \setminus \{m_i\}$.
		\item Построим отношение $T_i$ рекуррентно:
		\[ T_i := \{(m_i, a) \mid a \in A_{i+1}\} \cup T_{i+1} \]
		(т.е. $m_i$ ставится в отношение со всеми элементами, которые будут выбраны позже).
	\end{enumerate}
	
	Продолжаем процесс, пока $A_i \neq \emptyset$. В итоге получим $T := T_0$ — искомую топологическую сортировку. Докажем, что $T$ действительно является линейным порядком, согласованным с $R$.\\
	
	
	\textbf{Лемма:} Если $(a, b) \in R$, то существуют такие индексы $i, j \in \{0, 1, \dots, |A|-1\}$, что $a = m_i$, $b = m_j$, причем $i \le j$.
	
	\textit{Доказательство леммы:} Существование индексов следует из построения (каждый элемент когда-то будет выбран как минимальный). Предположим, что $j < i$. Тогда $a \in A_j$ (так как $a$ еще не удален), но $b$ уже удален на шаге $j$. Однако $b$ был минимальным в $A_j$, а $(a, b) \in R$ означает, что $a$ меньше $b$ (или предшествует ему), что противоречит минимальности $b$. Следовательно, $i \le j$. $\square$
	
	Теперь проверим свойства $T$:
	
	\begin{itemize}
		\item \textbf{Рефлексивность:} По построению $T$ не содержит пар вида $(x, x)$, но при построении линейного порядка мы можем добавить их искусственно или считать, что рефлексивность выполняется тривиально, так как $x$ предшествует сам себе в любом порядке.
		
		\item \textbf{Транзитивность:} Пусть $(a, b) \in T$ и $(b, c) \in T$. Тогда по построению существуют индексы $i, j, k$ такие, что $a = m_i$, $b = m_j$, $c = m_k$, причем $i < j$ (из $(a, b) \in T$) и $j < k$ (из $(b, c) \in T$). Следовательно, $i < k$, а значит $c \in A_i$ и по построению $(a, c) \in T$.
		
		\item \textbf{Антисимметричность:} Пусть $(a, b) \in T$ и $(b, a) \in T$. Тогда существуют индексы $i, j$ такие, что $a = m_i$, $b = m_j$. Из $(a, b) \in T$ следует $i \le j$, из $(b, a) \in T$ следует $j \le i$. Значит $i = j$, откуда $a = b$.
		
		\item \textbf{Линейность (полнота):} Для любых $a, b \in A$ существуют индексы $i, j$ такие, что $a = m_i$, $b = m_j$. Если $i < j$, то по построению $(a, b) \in T$; если $j < i$, то $(b, a) \in T$.
		
		\item \textbf{Согласованность с $R$:} Если $(a, b) \in R$, то по лемме $a = m_i$, $b = m_j$ с $i \le j$, следовательно $(a, b) \in T$.
	\end{itemize}
	
	Таким образом, $T$ является искомой топологической сортировкой. \qed
	
	\begin{remarkbox}{3.11: Замечание о построении}
		Топологическая сортировка, вообще говоря, \textbf{не единственна}. Если на каком-то шаге есть несколько минимальных элементов, выбор любого из них даст валидную сортировку. Разные последовательности выбора приводят к разным топологическим сортировкам.
	\end{remarkbox}
	
	\begin{examplebox}{3.12: Пример построения топологической сортировки}
		\begin{center}
			\begin{tikzpicture}[>=stealth, scale=1]
				
				\node[draw, circle, fill=red!30] (1) at (0,2) {1};
				\node[draw, circle, fill=red!30] (2) at (0,0) {2};
				\node[draw, circle, fill=red!30] (3) at (2,1) {3};
				\node[draw, circle, fill=red!30] (4) at (4,1) {4};
				
				\draw[->] (1) -- (3);
				\draw[->] (1) -- (4);
				\draw[->] (2) -- (3);
				\draw[->] (2) -- (4);
				\draw[->] (3) -- (4);
				
			\end{tikzpicture}
		\end{center}
		
		Рассмотрим множество $A = \{1, 2, 3, 4\}$ с частичным порядком $R$, заданным следующим образом: $1R2$, $1R3$, $1R4$, $2R4$, $3R4$ (элемент 1 предшествует всем, 2 и 3 предшествуют 4, при этом 2 и 3 несравнимы).\\
		
		\textbf{Первый вариант сортировки} (выбираем минимальные в порядке 1, 2, 3, 4):
		\begin{align*}
			A_0 &= \{1,2,3,4\}, \; m_0 = 1 \\
			A_1 &= \{2,3,4\}, \; T_0 = \{(1,2),(1,3),(1,4)\} \cup T_1 \\
			A_1 &: \; m_1 = 2 \\
			A_2 &= \{3,4\}, \; T_1 = \{(2,3),(2,4)\} \cup T_2 \\
			A_2 &: \; m_2 = 3 \\
			A_3 &= \{4\}, \; T_2 = \{(3,4)\} \cup T_3 \\
			A_3 &: \; m_3 = 4 \\
			A_4 &= \emptyset, \; T_3 = \emptyset
		\end{align*}
		Итог: $T = \{(1,2),(1,3),(1,4),(2,3),(2,4),(3,4)\}$, что соответствует порядку $1 \prec 2 \prec 3 \prec 4$.\\
		
		\textbf{Второй вариант сортировки} (выбираем минимальные в порядке 1, 3, 2, 4):
		\begin{align*}
			A_0 &= \{1,2,3,4\}, \; m'_0 = 1 \\
			A'_1 &= \{2,3,4\}, \; T'_0 = \{(1,2),(1,3),(1,4)\} \cup T'_1 \\
			A'_1 &: \; m'_1 = 3 \\
			A'_2 &= \{2,4\}, \; T'_1 = \{(3,2),(3,4)\} \cup T'_2 \\
			A'_2 &: \; m'_2 = 2 \\
			A'_3 &= \{4\}, \; T'_2 = \{(2,4)\} \cup T'_3 \\
			A'_3 &: \; m'_3 = 4 \\
			A'_4 &= \emptyset, \; T'_3 = \emptyset
		\end{align*}
		Итог: $T' = \{(1,2),(1,3),(1,4),(3,2),(3,4),(2,4)\}$, что соответствует порядку $1 \prec 3 \prec 2 \prec 4$.
	\end{examplebox}
	
	\begin{notebook}{3.13: Замечание о рефлексивности исходного порядка}
		При построении топологической сортировки важно учитывать, является ли исходный порядок $R$ рефлексивным или нет. Если $R$ — нестрогий порядок (рефлексивный), то $(x,x) \in R$, что тривиально выполняется. Если $R$ — строгий порядок (антирефлексивный), то $(x,x) \notin R$, что также не создает проблем, так как в результирующем линейном порядке мы можем явно не включать рефлексивные пары.
		
		Общая формула для построения на $i$-м шаге:
		\[ m_i = \min(A_i, R) \]
		\[ A_{i+1} = \begin{cases} 
			A_i \setminus \{m_i\}, & i < n \\ 
			\emptyset, & i \ge n 
		\end{cases} \]
		\[ T_i = \{(m_i, a) \mid a \in A_{i+1}\} \cup \{(m_i, m_j) \mid (m_i, m_j) \in R\} \cup T_{i+1} \]
	\end{notebook}
	
	
	
	
	
	
	
	
	
	
	% ==================== РАЗДЕЛ 4: РАСПИСАНИЯ И КРИТИЧЕСКИЕ ПУТИ ====================
	\newpage
	\section{Теория расписаний и критические пути}
	
	\subsection{Расписание, длина и высота расписания. Цепи и критические пути. Теорема о существовании кратчайшего расписания (4 вопрос)}
	
	Рассматривается задача выполнения набора работ $A$, между которыми задано отношение строгого частичного порядка $R$ (отношение зависимости: $(a,b) \in R$ означает, что работа $a$ должна быть выполнена до работы $b$). Каждая работа выполняется ровно одну единицу времени.
	
	\subsubsection{Основные определения}
	
	\begin{definitionbox}{4.1: Расписание}
		Пусть $A$ — непустое множество, строго упорядоченно относительно $R$. 
		\textbf{Расписанием} называется разбиение множества $A$ на непустые подмножества 
		$\{A_1, A_2, \dots, A_n\}$ такие, что:
		\[ \forall a,b \in A, a \neq b: (b,a) \in R \land a \in A_k \Rightarrow \exists j < k: b \in A_j \]
		
		Иными словами, если работа $b$ должна предшествовать работе $a$ (т.е. $(b,a) \in R$), 
		то они не могут находиться в одном слое расписания, причем $b$ обязательно находится 
		в слое с меньшим номером, чем $a$.
	\end{definitionbox}
	
	\myremark{Простыми словами:} Расписание — это разбиение множества работ на временные разбиения (слои) такие, что:
	\begin{itemize}
		\item Внутри одной части разбиения нет отношенией/зависимостей (все работы можно выполнять параллельно);
		\item Но есть между элементами разных частей, и они заданы так, что первый элемент всегда в более ранней части разбиения (т.е если работа $b$ должна предшествовать работе $a$, то $b$ находится в более ранней части разбиения, чем $a$).
	\end{itemize}
	
	\begin{examplebox}{4.2: Примеры расписаний}
		Для некоторого множества работ расписаниями могут быть:
		\[ \{\{1, 3\}, \{2, 7, 4\}, \{a, 9\}, \{12\}\} \]
		\[ \{\{1\}, \{3\}, \{2, 7\}, \{4\}, \{a, 9\}\} \]
		
		Здесь каждый блок $\{\dots\}$ соответствует одному временному слою (такту), 
		в котором работы выполняются параллельно.
	\end{examplebox}
	
	\begin{definitionbox}{4.3: Длина и высота расписания}
		\begin{itemize}
			\item \textbf{Длиной расписания} $L(F)$ называется количество слоев (тактов) в расписании $F$.
			\item \textbf{Высотой} частично упорядоченного множества $(A,R)$ называется максимальная длина цепи (см. определение ниже).
		\end{itemize}
	\end{definitionbox}
		
	\subsubsection{Функция глубины}
	
	\begin{definitionbox}{4.6: Глубина элемента}
		Для элемента $a \in A$ определим \textbf{глубину} $d(a)$ как длину максимальной цепи, заканчивающейся в элементе $a$:
		\[ d(a) = \max\{ |X| : X \text{ — цепь в } A, \text{ наибольший элемент } X = a \} \]
	\end{definitionbox}
	
	\begin{lemmabox}{4.7: Свойства глубины}
		\begin{enumerate}
			\item Для любого элемента $a \in A$: $1 \le d(a) \le h$, где $h$ — высота ЧУМ.
			\item Если $(b,a) \in R$, то $d(b) < d(a)$.
			\item В критическом пути $x_1, x_2, \dots, x_h$ выполняется $d(x_i) = i$ для всех $i$.
		\end{enumerate}
	\end{lemmabox}
	
	\myproof
	Рассмотрим критический путь $Z$, заканчивающийся в $z_n$. На этом пути есть предыдущий элемент $z_{n-1}$. Множество $Z \setminus \{z_n\}$ имеет те же свойства, что и $Z$. Для $Z$ наибольшим элементом является $z_n$, для $Z \setminus \{z_n\}$ наибольший элемент — $z_{n-1}$.
	
	Покажем, что $d(z_{n-1}) \neq d(z_n)$. Если $d(z_{n-1}) = d(z_n)$, то существует критический путь $Z'$, заканчивающийся в $z_{n-1}$. При этом $z_n \notin Z'$, т.к. $z_{n-1}$ — наибольший элемент $Z'$, $(z_{n-1}, z_n) \in R$ и $(z_n, z_{n-1}) \notin R$ (асимметричность). Тогда $|Z \setminus \{z_n\}| = d(z_n) - 1$, следовательно $d(z_{n-1}) = d(z_n) - 1$.
	
	Рассмотрим самый длинный критический путь $X$, $|X| = h$. $X$ заканчивается в своем наибольшем элементе $x_h$, $d(x_h) = h = \max_{a \in A} d(a)$. В $X$ все элементы строго линейно упорядочены, тогда $d(x_h) = h$, $d(x_{h-1}) = h-1$, $d(x_{h-2}) = h-2$ и т.д. Следовательно, $\forall i \in 1..h$ имеем $A_i \cap X \neq \emptyset$, где $A_i = \{a \in A : d(a) = i\}$. \qed
	
	\subsubsection{Теорема о кратчайшем расписании}
	
	\begin{theorembox}{4.8: О существовании кратчайшего расписания}
		Пусть $A$ — конечное строго частично упорядоченное множество относительно отношения $R$. 
		Рассмотрим множества $A_i = \{ x \in A : d(x) = i \}$ для $i = 1, 2, \dots, h$, где $h$ — высота $A$.
		Тогда семейство $\{A_i\}_{i=1}^{h}$ задает \textbf{кратчайшее} (наименьшей длины \ мощности) расписание на $A$.
	\end{theorembox}
	
	\myproof
	
	\textbf{Часть 1. Проверка, что $\{A_i\}$ — расписание.}
	
	\begin{enumerate}
		\item \textbf{Непустота:} $\forall i \in 1..h \; A_i \neq \emptyset$. Действительно, рассмотрим критический путь $X$ длины $h$. Как показано в лемме, $X \cap A_i \neq \emptyset$ для каждого $i$.
		
		\item \textbf{Дизъюнктность:} $\forall i \neq j \in 1..h \; A_i \cap A_j = \emptyset$, так как глубина элемента определяется единственным образом.
		
		\item \textbf{Полнота покрытия:} $\bigcup_{i=1}^{h} A_i = A$.
		\begin{itemize}
			\item $\bigcup A_i \subseteq A$ — очевидно.
			\item $\bigcup A_i \supseteq A$ — для любого $a \in A$ определена глубина $d(a) \in \mathbb{N}$, причем $d(a) \le h$, следовательно $a \in A_{d(a)}$.
		\end{itemize}
		Таким образом, $\{A_i\}_{i=1}^{h}$ является разбиением $A$.
		
		\item \textbf{Свойство расписания:} Пусть $a \neq b \in A$, $(b,a) \in R$ и $a \in A_k$. Покажем, что $\exists j < k: b \in A_j$.
		
		Предположим противное, что $j \ge k$, т.е. $d(b) \ge d(a) = k$. Рассмотрим критический путь $B$, заканчивающийся в $b$ (он существует по определению глубины). Тогда $d(b) \ge k$. Заметим, что $a \notin B$, иначе $(a,b) \in R$ (по транзитивности), что противоречило бы асимметричности.
		
		Рассмотрим множество $B \cup \{a\}$. Так как $\forall x \in B: (x,b) \in R$ и $(b,a) \in R$, то по транзитивности $(x,a) \in R$. Следовательно, $B \cup \{a\}$ строго линейно упорядочено и является цепью, заканчивающейся в $a$. Тогда $d(a) \ge |B \cup \{a\}| \ge k + 1$, что противоречит $d(a) = k$.
		
		Следовательно, $d(b) < d(a)$, т.е. $j < k$.
	\end{enumerate}
	
	\textbf{Часть 2. Доказательство кратчайшести.}
	
	Рассмотрим произвольное другое расписание $A'_1, A'_2, \dots, A'_s$ (где $s$ — его длина). Пусть $X$ — критический путь (самая длинная цепь) в $A$, $|X| = h$. Все элементы $X$ попарно сравнимы, поэтому в силу определения расписания они должны находиться в \textbf{различных} слоях расписания (если бы два элемента из $X$ попали в один слой, то между ними было бы отношение порядка, что запрещено). Следовательно, $s \ge h$.
	
	Таким образом, любое расписание имеет длину не меньше $h$, а построенное нами расписание $\{A_i\}_{i=1}^{h}$ имеет длину ровно $h$. Значит, оно является кратчайшим. \qed
	
	\begin{corollarybox}{4.9}
		Минимально возможное время выполнения всего комплекса работ (при неограниченном количестве исполнителей) равно длине критического пути.
	\end{corollarybox}
	
	\begin{center}
		\begin{tikzpicture}[scale=1.2]
			% Уровни (слои расписания)
			\foreach \y/\lev in {3/3,2/2,1/1,0/0} {
				\draw[gray, dashed] (-1,\y) -- (4,\y);
				\node[left] at (-1,\y) {Уровень \lev};
			}
			
			% Элементы
			\node[draw, circle, blue] (a) at (0,3) {$a$};
			\node[draw, circle, blue] (b) at (1,2) {$b$};
			\node[draw, circle, blue] (c) at (2,2) {$c$};
			\node[draw, circle, blue] (d) at (0,1) {$d$};
			\node[draw, circle, blue] (e) at (2,1) {$e$};
			\node[draw, circle, blue] (f) at (1,0) {$f$};
			
			\draw[->] (a) -- (b);
			\draw[->] (a) -- (c);
			\draw[->] (b) -- (d);
			\draw[->] (c) -- (e);
			\draw[->] (d) -- (f);
			\draw[->] (e) -- (f);
			
			\node at (1.5,-1) {Кратчайшее расписание (уровневое)};
		\end{tikzpicture}
		\captionof{figure}{Кратчайшее расписание: все элементы с глубиной $i$ попадают в $i$-й слой}
	\end{center}

	
	\begin{examplebox}{4.11: Пример построения кратчайшего расписания}
		Рассмотрим множество $A = \{1,2,3,4,5,6\}$ с порядком:
		$1 < 2$, $1 < 3$, $2 < 4$, $3 < 4$, $2 < 5$, $4 < 6$, $5 < 6$.
		
		Найдем глубины элементов:
		\begin{itemize}
			\item $d(1) = 1$ (нет предшественников)
			\item $d(2) = 2$ (путь 1-2)
			\item $d(3) = 2$ (путь 1-3)
			\item $d(4) = 3$ (путь 1-2-4 или 1-3-4)
			\item $d(5) = 3$ (путь 1-2-5)
			\item $d(6) = 4$ (путь 1-2-4-6 или 1-2-5-6)
		\end{itemize}
		
		Тогда кратчайшее расписание:
		\begin{align*}
			A_1 &= \{1\} \\
			A_2 &= \{2,3\} \\
			A_3 &= \{4,5\} \\
			A_4 &= \{6\}
		\end{align*}
		
		Длина расписания $= 4$, что совпадает с длиной критического пути $1 \to 2 \to 4 \to 6$ (или $1 \to 2 \to 5 \to 6$).
	\end{examplebox}
	
	
		\subsubsection{Цепи и критические пути}
	
	\begin{definitionbox}{4.4: Цепь (путь)}
		\textbf{Путем} (или \textbf{цепью}) в частично упорядоченном множестве называется любая последовательность элементов, в которой каждый последующий элемент больше предыдущего (в смысле порядка $R$):
		\[ a_1 R a_2 R a_3 R \dots R a_k \]
		
		\textbf{Полным путем} называется путь, который нельзя продолжить ни в начале, ни в конце (т.е. начинающийся с минимального и заканчивающийся максимальным элементом).
	\end{definitionbox}
	
	\begin{definitionbox}{4.5: Критический путь}
		\textbf{Критическим путем} называется наиболее продолжительный (самый длинный) полный путь в частично упорядоченном множестве. 
		
		Элементы, лежащие на критическом пути, называются \textbf{критическими}.
	\end{definitionbox}
	
	\begin{center}
		\begin{tikzpicture}[scale=1.2]
			% ЧУМ
			\node[draw, circle, blue] (a) at (0,3) {$a$};
			\node[draw, circle, blue] (b) at (-1,2) {$b$};
			\node[draw, circle, blue] (c) at (1,2) {$c$};
			\node[draw, circle, blue] (d) at (-1.5,1) {$d$};
			\node[draw, circle, blue] (e) at (0,1) {$e$};
			\node[draw, circle, blue] (f) at (1.5,1) {$f$};
			\node[draw, circle, blue] (g) at (0,0) {$g$};
			
			\draw[->] (a) -- (b);
			\draw[->] (a) -- (c);
			\draw[->] (b) -- (d);
			\draw[->] (b) -- (e);
			\draw[->] (c) -- (e);
			\draw[->] (c) -- (f);
			\draw[->] (d) -- (g);
			\draw[->] (e) -- (g);
			\draw[->] (f) -- (g);
			
			% Выделяем критический путь
			\draw[->, red, ultra thick] (a) -- (b);
			\draw[->, red, ultra thick] (b) -- (e);
			\draw[->, red, ultra thick] (e) -- (g);
			
			\node[red] at (0,-1) {Критический путь: $a \to b \to e \to g$, длина = 4};
		\end{tikzpicture}
		\captionof{figure}{Критический путь в частично упорядоченном множестве}
	\end{center}
	
	\myremark{Интуиция:} Высота множества — это минимально возможное время, за которое можно выполнить все задачи, даже если у нас бесконечное количество рабочих. Мы не можем закончить быстрее, чем выполним самую длинную цепочку зависимых задач (критический путь).
	
	\begin{definitionbox}{4.10: Сужение отношения}
		Пусть на множестве $A$ задано отношение $R$, и $\emptyset \neq X \subseteq A$. 
		\textbf{Сужением отношения $R$ на множество $X$} называется:
		\[ R(X) := R \cap (X \times X) \]
		т.е. все пары из $R$, оба элемента которых принадлежат $X$.
	\end{definitionbox}
	
	
	
	
	
	
	
	
	
	
	% ==================== РАЗДЕЛ 5: АНТИЦЕПИ ====================
	\newpage
	\section{Антицепи и разбиения}
	
	\subsection{Антицепи. Теорема о существовании разбиения на антицепи (5 вопрос)}
	
	\subsubsection{Определение антицепи}
	
	\begin{definitionbox}{5.1: Антицепь}
		Пусть $A$ — произвольное множество, строго частично упорядоченное отношением $R$. 
		Подмножество $X \subseteq A$ называется \textbf{антицепью}, если никакие два различных элемента из $X$ не сравнимы между собой:
		\[ \forall x, y \in X, x \neq y: (x, y) \notin R(X) \]
		где $R(X) = R \cap (X \times X)$ — сужение отношения $R$ на множество $X$.
	\end{definitionbox}
	
	\myremark{Напоминание:} $R(X)$ — сужение отношения $R$ на подмножество $X$, т.е. все пары из $R$, оба элемента которых принадлежат $X$.
	
	\begin{examplebox}{5.2: Примеры антицепей}
		Рассмотрим множество $A = \{1, 3, 5, 13\}$ с отношением порядка:
		\[ R = \{(1, 3), (3, 5), (1, 5)\} \]
		
		Антицепями в этом множестве являются:
		\[ X = \{1, 13\}, \{3, 13\}, \{5, 13\}, \{1\}, \{3\}, \{5\}, \{13\} \]
		
		Заметим, что $\{1, 3\}$ не является антицепью, так как $(1,3) \in R$.
	\end{examplebox}
	
	\begin{examplebox}{5.3: Более сложный пример}
		Рассмотрим множество $A = \{1, 2, 3, 4, 7, n, 9, 12\}$ с отношением:
		\begin{align*}
			R(A) = \{ &(1, 2), (2, n), (1, n), (n, 12), (1, 12), (2, 12), \\
			&(7, n), (7, 9), (9, 12), (7, 12), (3, 4)\}
		\end{align*}
		
		Примерами антицепей являются:
		\begin{itemize}
			\item $\{1, 3\}$ — элементы 1 и 3 не связаны отношением
			\item $\{1, 3, 7\}$ — все элементы попарно несравнимы
			\item $\{3, 7\}$, $\{1, 7\}$, $\{4, 7\}$ и т.д.
		\end{itemize}
		
		\myremark{Важно:} $\{1, 7, 12\}$ не является антицепью, так как $(1,12) \in R$ и $(7,12) \in R$ (по транзитивности).
		
		Для любого отдельно взятого элемента множество $\{a\}$ всегда является антицепью, так как один элемент ни с чем не сравним.
	\end{examplebox}
	
	\subsubsection{Высота множества и разбиение на антицепи}
	
	\begin{definitionbox}{5.4: Высота множества}
		\textbf{Высотой} частично упорядоченного множества называется максимальная глубина элемента в нём.
		
		Напоминание: \textbf{Глубина} элемента $a$ — длина критического пути (максимальной цепи), заканчивающегося в $a$.
	\end{definitionbox}
	
	\begin{theorembox}{5.5: О разбиении на антицепи}
		Всякое конечное частично упорядоченное множество высоты $h$ можно разбить на $h$ антицепей.
	\end{theorembox}
	
	
	\begin{remarkbox}{5.6: Связь с расписаниями}
		\textbf{Кратчайшее расписание} (построенное в теореме 4.8) в точности является таким разбиением на антицепи по уровням глубины.
		
		В примере 5.3 кратчайшее расписание (разбиение на антицепи) имеет вид:
		\[ \{\{1, 3\}, \{2, 4, 7\}, \{n, 9\}, \{12\}\} \]
		
		Это разбиение на 4 антицепи, что соответствует высоте множества (глубина элемента 12 равна 4).
	\end{remarkbox}
	
	\begin{center}
		\begin{tikzpicture}[>=stealth, scale=1.2]
			
			% ---------- СХЕМА ----------
			% Вершины
			\node[draw, circle] (1)  at (0,4)  {1};
			\node[draw, circle] (3)  at (3,4)  {3};
			
			\node[draw, circle] (2)  at (0,3)  {2};
			\node[draw, circle] (4)  at (3,3)  {4};
			
			\node[draw, circle] (7)  at (1.5,2) {7};
			
			\node[draw, circle] (n)  at (1.5,1) {$n$};
			\node[draw, circle] (9)  at (3.5,1) {9};
			
			\node[draw, circle] (12) at (2.5,0) {12};
			
			% Стрелки
			\draw[->] (1) -- (2);
			\draw[->] (3) -- (4);
			
			\draw[->] (2) -- (n);
			\draw[->] (7) -- (n);
			
			\draw[->] (4) -- (9);
			\draw[->] (7) -- (9);
			
			\draw[->] (n) -- (12);
			\draw[->] (9) -- (12);
			
			% ---------- ТЕКСТ СПРАВА ----------
			
			\node[align=left] at (9,4) {\textbf{Цепь} — стрелка или их последовательность};
			
			\node[align=left] at (9,3.2) {Критический путь для $n$ — $\{1,2,n\}$};
			
			\node[align=left] at (9,2.5) {Для $9$ — $\{7,9\}$};
			
			\node[align=left] at (9,1.8) {Глубина $n = 3$, глубина $9 = 2$, глубина $12 = 4$};
			
			\node[align=left] at (9,1) {Кратчайшее расписание = $\{\{1,3\}, \{2,7,4\}, \{n,9\}, \{12\}\}$};
			
		\end{tikzpicture}
	\end{center}
	
	{1}, {2}, {3}, {4}, {7}, {n}, {9}, {12}, {1, 3}, {1, 3, 7} , {3, 7}, {1, 7}, {4, 7} и т.д. – антицепи (но {1, 7, 12}, например, нет из-за транзитивности ЧУМ)
	
	Будет существовать антицепь для любого элемента, так как один элемент в множестве ни с чем не сравним.
	
	Высота множества = Глубине 12 = 4
	
	Кратчайшее разбиение A на антицепи = {{1, 3}, {2, 4, 7}, {n, 9}, {12}}  = кратчайшему расписанию - это не единственное разбиение на антицепи, которое можно сделать, и не единственная длина таких разбиений.
	
	\subsubsection{Лемма Дилуорса}
	
	\begin{theorembox}{5.7: Лемма Дилуорса}
		Во всяком конечном частично упорядоченном множестве $A$, $|A| = n$, для любого $t \in \{1, \ldots, n\}$ существует либо цепь длины $> t$, либо антицепь длины $\geq \dfrac{n}{t}$.\\
		
		(пояснение: Для всех чисел от 1 до длины A, есть цепи большие этого числа, либо антицепи большие противоположной длины)
	\end{theorembox}
	
	\myproof (От противного)
	Предположим, что максимальная цепь в $A$ имеет длину $\le t$. Тогда по теореме 5.5 множество $A$ можно разбить на $t$ антицепей. Пусть $l$ — длина самой длинной из этих антицепей, и предположим, что $l < \dfrac{n}{t}$.
	
	Тогда в $A$ не может быть больше, чем $l \times t$ элементов (так как всего $t$ антицепей, каждая размера не больше $l$). Следовательно:
	\[ |A| \le l \times t < \frac{n}{t} \times t = n \]
	Получили противоречие с тем, что $|A| = n$. \qed
	
	\begin{examplebox}{5.8: Пример применения леммы Дилуорса}
		Рассмотрим $A = \{1, 3, 5\}$ — строго линейно упорядоченное множество:
		\[ R = \{(1, 3), (3, 5), (1, 5)\} \]
		
		Для различных $t$ (здесь $n=3$):
		\begin{itemize}
			\item $t = 3$: 
			\begin{itemize}
				\item цепей длины $>3$ нет (максимальная цепь имеет длину 3);
				\item есть антицепь длины $\ge \frac{3}{3}=1$ (например, $\{1\}$, $\{3\}$ или $\{5\}$).
			\end{itemize}
			\item $t = 2$: 
			\begin{itemize}
				\item есть цепь длины $3 > 2$ (весь $A$: $1<3<5$);
				\item антицепей длины $\ge \frac{3}{2}=1.5$ не требуется, так как условие уже выполнено цепью.
			\end{itemize}
			\item $t = 1$: 
			\begin{itemize}
				\item есть цепь длины $3 > 1$;
				\item условие выполнено.
			\end{itemize}
		\end{itemize}
		
		Лемма Дилуорса утверждает: для любого $t$ найдётся либо цепь длиннее $t$, либо антицепь длиной не меньше $n/t$. В данном примере это выполняется для всех $t$.
	\end{examplebox}
	
	
	\begin{theorembox}{5.9: Теорема о существовании разбиения на антицепи}
		Кратчайшее расписание является разбиением на $h$ антицепей, где $h$ — высота множества.
	\end{theorembox}
	
	\myproof
	Следует из доказательства теоремы о существовании кратчайшего расписания. \qed
	
	
	
	
	
	
	% ==================== РАЗДЕЛ 6: ВОЗРАСТАЮЩИЕ ПОДПОСЛЕДОВАТЕЛЬНОСТИ ====================
	\newpage
	\section{Наибольшая возрастающая подпоследовательность}
	
	\subsection{Задача о построении наибольшей возрастающей подпоследовательности (6 вопрос)}
	
	\subsubsection{Основные определения}
	
	\begin{definitionbox}{6.1: Последовательность и подпоследовательность}
		\textbf{Последовательностью} называется упорядоченный набор элементов:
		\[ S = (s_1, s_2, \dots, s_n) \]
		
		\textbf{Подпоследовательностью} называется последовательность, полученная из исходной путем удаления некоторых элементов без изменения порядка оставшихся.
	\end{definitionbox}
	
	\begin{definitionbox}{6.2: Отношения на последовательности}
		На множестве элементов последовательности $S$ определим два отношения:
		\begin{itemize}
			\item $a <_S b$ — отношение порядка следования: $a$ предшествует $b$ в последовательности $S$ (т.е. индекс $a$ меньше индекса $b$).
			\item $a \prec b$ — отношение частичного порядка, задаваемое условием:
			\[ a \prec b \iff \begin{cases} 
				a <_S b & (\text{порядок следования}) \\
				a < b & (\text{числовое сравнение})
			\end{cases} \]
		\end{itemize}
		
		Иными словами, $a \prec b$ означает, что $a$ расположено раньше $b$ в последовательности и при этом $a$ меньше $b$ по значению.
	\end{definitionbox}
	
	\begin{examplebox}{6.3: Пример}
		Рассмотрим последовательность:
		\[ S = \{3, 2, 1, 8, 9, 4, 5, 6, 7\} \]
		
		\begin{itemize}
			\item $S_1 = \{2, 1, 4, 6\}$ — подпоследовательность (сохранен порядок индексов)
			\item $S_2 = \{1, 4, 5, 6\}$ — \textbf{возрастающая} подпоследовательность
			\item $S_3 = \{3, 2, 1\}$ — \textbf{убывающая} подпоследовательность
		\end{itemize}
	\end{examplebox}
	
	\begin{center}
		\begin{tikzpicture}[scale=0.9]
			
			% Оси
			\draw[->] (0,0) -- (10,0) node[right] {индекс};
			\draw[->] (0,0) -- (0,10) node[above] {значение};
			
			% Деления по X
			\foreach \x in {1,...,9}
			\node[below] at (\x,0) {\x};
			
			% Деления по Y
			\foreach \y in {1,...,9}
			\node[left] at (0,\y) {\y};
			
			% Все точки
			\foreach \x/\y in {1/3, 2/2, 3/1, 4/8, 5/9, 6/4, 7/5, 8/6, 9/7}
			\fill[blue] (\x,\y) circle (2.5pt);
			
			% НВП
			\draw[red, thick] (3,1) -- (6,4) -- (7,5) -- (8,6) -- (9,7);
			
			\foreach \x/\y in {3/1, 6/4, 7/5, 8/6, 9/7}
			\fill[red] (\x,\y) circle (3pt);
			
			\node at (5,-1.2) {НВП: индексы $(3,6,7,8,9)$};
			
		\end{tikzpicture}
		\captionof{figure}{Наибольшая возрастающая подпоследовательность}
	\end{center}
	
	\begin{definitionbox}{6.4: Цепи и антицепи в последовательности}
		Рассмотрим последовательность $S$ с отношением $\prec$ (как определено выше).
		\begin{itemize}
			\item \textbf{Цепь} в этом частичном порядке — это \textbf{возрастающая подпоследовательность}.
			\item \textbf{Антицепь} в этом частичном порядке — это \textbf{убывающая подпоследовательность}.
		\end{itemize}
	\end{definitionbox}
	
	\subsubsection{Алгоритм построения минимального разбиения на антицепи}
	
	\begin{algorithmbox}{6.5: Алгоритм разбиения на убывающие подпоследовательности}
		Проходя по последовательности слева направо, будем добавлять каждый следующий элемент в минимальный по номеру элемент разбиения, в который его можно добавить (т.е. в первый список, где последний элемент больше текущего).
	\end{algorithmbox}
	
	\begin{examplebox}{6.6: Пример работы алгоритма}
		Дана последовательность: $3, 5, 8, 9, 4, 6, 1, 2, 7, 10$.
		
		Пошаговое построение:
		\begin{align*}
			&3 \to \{3\} \\
			&3, 5 \to \{3\}, \{5\} \quad (5 > 3 \Rightarrow \text{новый список}) \\
			&3, 5, 8 \to \{3\}, \{5\}, \{8\} \\
			&3, 5, 8, 9 \to \{3\}, \{5\}, \{8\}, \{9\} \\
			&3, 5, 8, 9, 4 \to \{3\}, \{5, 4\}, \{8\}, \{9\} \quad (4 < 5 \Rightarrow \text{в список с }5) \\
			&3, 5, 8, 9, 4, 6 \to \{3\}, \{5, 4\}, \{8, 6\}, \{9\} \quad (6 < 8 \Rightarrow \text{в список с }8) \\
			&3, 5, 8, 9, 4, 6, 1 \to \{3, 1\}, \{5, 4\}, \{8, 6\}, \{9\} \\
			&3, 5, 8, 9, 4, 6, 1, 2 \to \{3, 1\}, \{5, 4, 2\}, \{8, 6\}, \{9\} \\
			&3, 5, 8, 9, 4, 6, 1, 2, 7 \to \{3, 1\}, \{5, 4, 2\}, \{8, 6\}, \{9, 7\} \\
			&3, 5, 8, 9, 4, 6, 1, 2, 7, 10 \to \{3, 1\}, \{5, 4, 2\}, \{8, 6\}, \{9, 7\}, \{10\}
		\end{align*}
		
		Итоговое разбиение на антицепи (убывающие подпоследовательности):
		\[ \{3, 1\}, \{5, 4, 2\}, \{8, 6\}, \{9, 7\}, \{10\} \]
	\end{examplebox}
	
	\subsubsection{Построение наибольшей возрастающей подпоследовательности}
	
	\begin{theorembox}{6.7: Построение НВП}
		Пусть $\Lambda = \{\Lambda_1, \Lambda_2, \dots, \Lambda_h\}$ — минимальное разбиение последовательности на антицепи (убывающие подпоследовательности), полученное алгоритмом 6.5. Тогда можно построить возрастающую подпоследовательность длины $h$ следующим образом:
		
		\begin{enumerate}
			\item Выберем произвольный элемент $A_h \in \Lambda_h$ (из последней антицепи).
			\item Для каждого $i = h-1, h-2, \dots, 1$ найдем элемент $A_i \in \Lambda_i$ такой, что:
			\begin{itemize}
				\item $A_i < A_{i+1}$ (по значению)
				\item $A_i$ предшествует $A_{i+1}$ в исходной последовательности
			\end{itemize}
			Такой элемент существует и называется $prev(A_{i+1})$.
		\end{enumerate}
		
		Полученная цепочка $A_1, A_2, \dots, A_h$ является \textbf{наибольшей возрастающей подпоследовательностью}.
	\end{theorembox}
	
	
	\myproof
	\begin{enumerate}
		\item \textbf{Существование $prev$:} Для любого $a \in \Lambda_i$ ($i > 1$) существует элемент $prev(a) \in \Lambda_{i-1}$, меньший $a$ и предшествующий ему. Если бы такого элемента не было, то по алгоритму $a$ попал бы в $\Lambda_{i-1}$.
		
		\item \textbf{Длина $h$ максимальна:} 
		\begin{itemize}
			\item В последовательности $S$ существует цепь (возрастающая подпоследовательность) длины $h$ (построенная выше).
			\item Следовательно, $S$ нельзя разбить менее чем на $h$ антицепей (по принципу Дирихле: если бы антицепей было меньше, то два сравнимых элемента из этой цепи попали бы в одну антицепь, что невозможно). \\
				
			*принцип Дирихле простыми словами – если в n клетках сидит m зайцев, m > n, то в хотя бы одной клетке сидит, по крайней мере, два зайца
			
			
			\item Значит, разбиение $\Lambda$ является \textbf{минимальным}.
			\item Так как $\Lambda$ минимально, не существует цепи длины $> h$ (иначе два элемента более длинной цепи попали бы в один элемент разбиения $\Lambda$).
		\end{itemize}
		
		Таким образом, построенная цепь имеет максимально возможную длину $h$.
	\end{enumerate}
	
	\begin{corollarybox}{6.8}
		Длина наибольшей возрастающей подпоследовательности равна минимальному количеству убывающих подпоследовательностей, на которые можно разбить исходную последовательность.
	\end{corollarybox}
	
	\begin{examplebox}{6.9: Пример построения НВП}
		Для последовательности из примера 6.6:
		\[ \Lambda = \{\{3, 1\}, \{5, 4, 2\}, \{8, 6\}, \{9, 7\}, \{10\}\} \]
		
		Строим цепь:
		\begin{align*}
			A_5 &= 10 \in \Lambda_5 \\
			prev(10) &= 7 \in \Lambda_4 \quad (7 < 10) \\
			prev(7) &= 6 \in \Lambda_3 \quad (6 < 7) \\
			prev(6) &= 4 \in \Lambda_2 \quad (4 < 6) \\
			prev(4) &= 1 \in \Lambda_1 \quad (1 < 4)
		\end{align*}
		
		Получаем наибольшую возрастающую подпоследовательность:
		\[ 1, 4, 6, 7, 10 \]
		
		Её длина $h = 5$, и это действительно максимально возможная длина.
	\end{examplebox}
	
	\begin{notebook}{6.10: Связь с теорией ЧУМ}
		Задача о наибольшей возрастающей подпоследовательности полностью эквивалентна задаче поиска максимальной цепи в частично упорядоченном множестве, где порядок задается двумя условиями: сохранение порядка индексов и возрастание значений.
	\end{notebook}
	
	\begin{center}
		\begin{tikzpicture}[scale=0.9]
			
			% Оси
			\draw[->] (0,0) -- (11,0) node[right] {индекс};
			\draw[->] (0,0) -- (0,11) node[above] {значение};
			
			% Подписи осей
			\foreach \x in {1,...,10}
			\node[below] at (\x,0) {\x};
			
			\foreach \y in {1,...,10}
			\node[left] at (0,\y) {\y};
			
			% Последовательность:
			% (3,5,8,9,4,6,1,2,7,10)
			
			\foreach \x/\y in {
				1/3, 2/5, 3/8, 4/9, 5/4,
				6/6, 7/1, 8/2, 9/7, 10/10}
			{
				\fill[blue] (\x,\y) circle (2.5pt);
			}
			
			% НВП: 1,4,6,7,10  (индексы 7,5,6,9,10)
			\draw[red, thick]
			(7,1) -- (5,4) -- (6,6) -- (9,7) -- (10,10);
			
			\foreach \x/\y in {7/1, 5/4, 6/6, 9/7, 10/10}
			\fill[red] (\x,\y) circle (3pt);
			
			\node at (5,-1.5)
			{НВП: $1 \rightarrow 4 \rightarrow 6 \rightarrow 7 \rightarrow 10$};
			
		\end{tikzpicture}
		\captionof{figure}{НВП как максимальная цепь в ЧУМ}
	\end{center}
	
	
		
		
		
		
		
		
		
	% ==================== РАЗДЕЛ 7: ТЕОРЕМА ДИЛУОРСА ====================
	\newpage
	\section{Теорема Дилуорса}
	
	\subsection{Теорема Дилуорса. Первое доказательство (7 вопрос)}
	
	\begin{theorembox}{7.1: Теорема Дилуорса}
		Пусть $A$ — конечное частично упорядоченное множество. Обозначим:
		\begin{itemize}
			\item $m$ — наименьшее число непересекающихся цепей, которыми может быть покрыто множество $A$ (минимальное разбиение на цепи);
			\item $M$ — наибольшая мощность антицепи в $A$ (размер максимальной антицепи).
		\end{itemize}
		Тогда $M = m$, т.е. \textbf{размер наибольшей антицепи равен размеру наименьшего разбиения на непересекающиеся цепи}.
	\end{theorembox}
	
	\begin{center}
		\begin{tikzpicture}[scale=1.2]
			% Пример частично упорядоченного множества
			\node[draw, circle, blue] (a) at (0,2) {$a$};
			\node[draw, circle, blue] (b) at (-1,1) {$b$};
			\node[draw, circle, blue] (c) at (1,1) {$c$};
			\node[draw, circle, blue] (d) at (0,0) {$d$};
			
			\draw[->] (a) -- (b);
			\draw[->] (a) -- (c);
			\draw[->] (b) -- (d);
			\draw[->] (c) -- (d);
			
			% Выделяем антицепь
			\fill[red, opacity=0.3] (-1.2,0.8) rectangle (1.2,1.2);
			\node[red] at (0,1.5) {Антицепь $\{b, c\}$ (мощность 2)};
			
			% Разбиение на цепи
			\node at (3,2) {Разбиение на цепи:};
			\node at (3,1.5) {$\{a, b, d\}$};
			\node at (3,1) {$\{c\}$};
			\node at (3,0.5) {Всего 2 цепи};
		\end{tikzpicture}
		\captionof{figure}{Иллюстрация теоремы Дилуорса: наибольшая антицепь $\{b,c\}$ имеет размер 2, и минимальное разбиение на цепи тоже состоит из 2 цепей}
	\end{center}
	
	\subsubsection{Предварительные замечания}
	
	\begin{remarkbox}{7.2: Интуитивное понимание}
		Рассмотрим пример множества, изображенного на рисунке. 
		\begin{itemize}
			\item \textbf{Непересекающиеся цепи:} например, $\{a, b, d\}$ и $\{c\}$ — всего 2 цепи. Меньше двух цепей сделать нельзя, так как элементы $b$ и $c$ несравнимы и не могут лежать в одной цепи.
			\item \textbf{Наибольшая антицепь:} $\{b, c\}$ — эти два элемента не состоят ни в каких отношениях. Мощность этой антицепи равна 2.
		\end{itemize}
		Таким образом, $M = m = 2$, что соответствует теореме.
	\end{remarkbox}
	
	\myproof
	
	Прежде чем перейти к формальному доказательству, установим очевидную оценку:
	
	\begin{lemmabox}{7.3: Простейшая оценка}
		Для любой антицепи $c$ и любого разбиения $\Lambda$ на цепи:
		\[ |c| \le |\Lambda| \]
		где $|\Lambda|$ — количество цепей в разбиении.
	\end{lemmabox}
	
	\myproof
	Если бы $|c| > |\Lambda|$, то по принципу Дирихле два элемента антицепи $c$ попали бы в одну цепь разбиения $\Lambda$, что невозможно, так как элементы одной цепи сравнимы, а элементы антицепи — нет. \qed
	
	Из этой леммы непосредственно следует, что для максимальной антицепи $c_{\text{max}}$ и минимального разбиения на цепи $\Lambda_{\text{min}}$:
	\[ |c_{\text{max}}| \le |\Lambda_{\text{min}}| \]
	
	Если существует антицепь $c_0$ и разбиение $\Lambda_0$ такие, что $|c_0| = |\Lambda_0|$, то очевидно $c_0$ — максимальная антицепь, а $\Lambda_0$ — минимальное разбиение. Наша цель — доказать, что такие $c_0$ и $\Lambda_0$ всегда существуют, т.е. что неравенство превращается в равенство.
	
	\subsubsection{Доказательство по индукции}
	
	Докажем теорему методом математической индукции по $|A|$.
	
	\textbf{База индукции:} 
	\begin{itemize}
		\item $A = \varnothing$: тривиально, $M = m = 0$.
		\item $|A| = 1$: единственный элемент образует и цепь, и антицепь, поэтому $M = m = 1$.
		\item $R = \varnothing$ (отношение пусто): все элементы попарно несравнимы, значит всё множество $A$ является антицепью, и одновременно каждая цепь состоит из одного элемента, поэтому минимальное разбиение на цепи — это разбиение на $|A|$ одноэлементных цепей. Таким образом, $M = |A| = m$.
	\end{itemize}
	
	\textbf{Шаг индукции:} Пусть $|A| \ge 2$, $R \neq \varnothing$. Выберем максимальный элемент в $A$ (относительно порядка $R$) и обозначим его $m$. По предположению индукции для множества $A \setminus \{m\}$ теорема верна.
	
	Следовательно, $A \setminus \{m\}$ можно разбить на цепи $C_1, C_2, \dots, C_k$, где $k$ — размер максимальной антицепи в $A \setminus \{m\}$. Обозначим через $k$ это число.
	
	\begin{center}
		\begin{tikzpicture}
			\draw[thick, blue] (0,0) rectangle (4,3);
			\node at (2,3.2) {$A$};
			\node[draw, circle, red] (m) at (3,2.5) {$m$};
			
			\draw[thick, dashed] (0,1.5) -- (4,1.5);
			\node at (2,1.7) {$A \setminus \{m\}$};
			
			\draw[thick, green!50!black] (0.5,0.5) -- (1.5,1.2) -- (2,0.8);
			\node at (0.5,0.2) {$C_1$};
			\draw[thick, orange] (2.5,0.3) -- (3.2,0.9) -- (3.8,0.5);
			\node at (3.5,0.2) {$C_k$};
		\end{tikzpicture}
		\captionof{figure}{Разбиение $A \setminus \{m\}$ на цепи $C_1,\dots,C_k$}
	\end{center}
	
	\textbf{Ключевое наблюдение:} Каждая антицепь размера $k$ в $A \setminus \{m\}$ пересекается с каждой цепью $C_i$ ровно по одному элементу. Действительно, если бы какая-то антицепь пересекалась с некоторой цепью более чем по одному элементу, то эти два элемента были бы сравнимы (как элементы одной цепи) и не могли бы одновременно принадлежать антицепи. А если бы антицепь не пересекалась с какой-то цепью, то общее количество элементов в антицепи было бы меньше $k$, так как всего цепей $k$.
	
	\begin{definitionbox}{7.4: Определение элементов $c_i$}
		Для каждого $i \in \{1,\dots,k\}$ определим $c_i$ как \textbf{максимальный элемент} в цепи $C_i$, который принадлежит хотя бы одной антицепи размера $k$ в множестве $A \setminus \{m\}$.
	\end{definitionbox}
	
	\begin{lemmabox}{7.5: \texorpdfstring{Множество $Y = \{c_1, \dots, c_k\}$ является антицепью}{Множество Y = {c_1, ..., c_k} является антицепью}}
	\end{lemmabox}
	
	\myproof
	Для каждого $c_i$ зафиксируем некоторую антицепь $Y_i$ размера $k$, содержащую $c_i$ (для разных $i$ антицепи $Y_i$ могут быть разными).
	
	Рассмотрим произвольные $i \neq j$. Так как $Y_i$ — антицепь размера $k$, она пересекается с каждой цепью $C_t$ ровно по одному элементу. В частности, $Y_i \cap C_j$ состоит из одного элемента. Обозначим этот элемент через $y_{ij}$:
	\[ Y_i \cap C_j = \{y_{ij}\} \]
	
	По определению $c_j$ как максимального элемента в $C_j$, принадлежащего некоторой антицепи размера $k$, имеем $(y_{ij}, c_j) \in R$ (возможно, $y_{ij} = c_j$, если $y_{ij}$ сам является максимальным).
	
	\begin{center}
		\begin{tikzpicture}
			\draw[thick, blue] (0,0) rectangle (3,2);
			\node at (1.5,2.2) {$C_i$};
			\draw[thick, blue] (4,0) rectangle (7,2);
			\node at (5.5,2.2) {$C_j$};
			
			\node[draw, circle, red] (ci) at (1.5,1.5) {$c_i$};
			\node[draw, circle, red] (cj) at (5.5,1.5) {$c_j$};
			\node[draw, circle, green] (y) at (5.5,0.8) {$y_{ij}$};
			
			\draw[->, dashed] (y) -- (cj);
			
			\draw[thick, purple, decorate, decoration={brace, mirror, amplitude=10pt}] (2.5,1) -- (4.5,1);
			\node at (3.5,0.5) {$Y_i$};
		\end{tikzpicture}
		\captionof{figure}{Иллюстрация: $y_{ij} \in C_j \cap Y_i$, $(y_{ij}, c_j) \in R$}
	\end{center}
	
	Предположим теперь, что $(c_j, c_i) \in R$. Тогда по транзитивности $(y_{ij}, c_i) \in R$ (так как $(y_{ij}, c_j) \in R$ и $(c_j, c_i) \in R$). Но $y_{ij} \in Y_i$ и $c_i \in Y_i$, а $Y_i$ — антицепь, поэтому никакие два её элемента не могут быть сравнимы. Получили противоречие. Следовательно, $(c_j, c_i) \notin R$.
	
	Аналогично, рассматривая антицепь $Y_j$, получаем $(c_i, c_j) \notin R$. Таким образом, никакие два элемента из $Y$ не сравнимы, значит $Y$ — антицепь. \qed
	
	Теперь рассмотрим два возможных случая относительно максимального элемента $m$.
	
	\subsubsection{Случай 1: Существует $i$ такой, что $(c_i, m) \in R$}
	
	В этом случае $m$ больше элемента $c_i$ (в смысле порядка). Построим множество:
	\[ X := \{x \in C_i : (x, c_i) \in R\} \cup \{m\} \]
	т.е. $X$ содержит все элементы цепи $C_i$, которые не больше $c_i$ (включая сам $c_i$), и добавленный $m$.
	
	\begin{center}
		\begin{tikzpicture}
			\draw[thick, blue] (0,0) rectangle (3,3);
			\node at (1.5,3.2) {$C_i$};
			
			\node[draw, circle, blue] (a) at (1.5,2.5) {$a_1$};
			\node[draw, circle, blue] (b) at (1.5,1.8) {$a_2$};
			\node[draw, circle, blue] (c) at (1.5,1.1) {$c_i$};
			
			\draw[->] (a) -- (b);
			\draw[->] (b) -- (c);
			
			\draw[thick, red, decorate, decoration={brace, amplitude=10pt}] (2,2.8) -- (2,0.8);
			\node at (2.5,1.8) {$X$};
			
			\node[draw, circle, red] (m) at (1.5,0.3) {$m$};
			\draw[->, red] (c) -- (m);
		\end{tikzpicture}
		\captionof{figure}{Построение множества $X$ как цепи}
	\end{center}
	
	По построению $X$ является цепью (все элементы $C_i$ сравнимы между собой, и $m$ больше $c_i$, а значит и всех элементов, меньших $c_i$).
	
	Рассмотрим множество $A \setminus X$. Утверждается, что в нём нет антицепи размера $k$. Действительно, любая антицепь размера $k$ в исходном $A \setminus \{m\}$ должна была иметь хотя бы один элемент в $C_i$. Но в $C_i \setminus X$ (т.е. в той части $C_i$, которая осталась после удаления $X$) нет элементов, которые входят в антицепи размера $k$, так как $c_i$ был максимальным таким элементом, а все меньшие мы удалили вместе с $X$.
	
	Однако в $A \setminus X$ есть антицепь размера $k-1$ — это $Y \setminus \{c_i\}$ (так как $c_i$ был удалён). По предположению индукции, $A \setminus X$ можно разбить на $k-1$ цепей. Добавляя к этому разбиению цепь $X$, получаем разбиение всего $A$ на $k$ цепей.
	
	Таким образом, в этом случае мы построили разбиение на $k$ цепей. Поскольку мы уже имели антицепь $Y$ размера $k$, то $M \ge k$ и $m \le k$, следовательно $M = m = k$.
	
	\subsubsection{Случай 2: Для всех $i$ выполняется $(c_i, m) \notin R$}
	
	Так как $m$ — максимальный элемент в $A$, то и $(m, c_i) \notin R$ (иначе $m$ был бы меньше $c_i$, что противоречит максимальности). Следовательно, $m$ не сравним ни с одним элементом из антицепи $Y$.
	
	\begin{center}
		\begin{tikzpicture}
			\node[draw, circle, blue] (c1) at (0,2) {$c_1$};
			\node[draw, circle, blue] (c2) at (1,2) {$c_2$};
			\node[draw, circle, blue] (ck) at (2,2) {$c_k$};
			
			\node[draw, circle, red] (m) at (1,0) {$m$};
			
			\draw[<->, dashed] (c1) -- (m);
			\draw[<->, dashed] (c2) -- (m);
			\draw[<->, dashed] (ck) -- (m);
			
			\node at (1,1) {несравнимы};
		\end{tikzpicture}
		\captionof{figure}{$m$ несравним с элементами $Y$}
	\end{center}
	
	Тогда $Y \cup \{m\}$ является антицепью размера $k+1$. Действительно, элементы $Y$ попарно несравнимы по доказанному, и $m$ не сравним ни с одним из них.
	
	Рассмотрим разбиение $\{C_1, \dots, C_k\}$ множества $A \setminus \{m\}$. Добавим к нему одноэлементную цепь $\{m\}$, получим разбиение $A$ на $k+1$ цепей.
	
	Покажем, что меньше чем $k+1$ цепей быть не может. Действительно, для любого разбиения $\Lambda$ множества $A$ на цепи, по лемме 7.3, $|\Lambda| \ge |Y \cup \{m\}| = k+1$. Следовательно, $k+1$ — минимально возможное количество цепей.
	
	Таким образом, в этом случае максимальная антицепь имеет размер $k+1$, и минимальное разбиение на цепи состоит из $k+1$ цепи, т.е. $M = m = k+1$.
	
	В обоих случаях мы получили равенство $M = m$. Индукционный переход завершён. \qed
	
	
	
	
	
	
	
	
	
	\begin{corollarybox}{7.6: Следствие}
		Теорема Дилуорса устанавливает двойственность между цепями и антицепями в частично упорядоченных множествах: минимальное количество цепей, необходимых для покрытия множества, равно максимальному размеру антицепи.
	\end{corollarybox}
	
	\begin{examplebox}{7.7: Применение теоремы Дилуорса}
		Рассмотрим множество прямоугольников, где $(w_1, h_1) \le (w_2, h_2)$ если $w_1 < w_2$ и $h_1 < h_2$ (один прямоугольник можно вложить в другой). 
		
		\begin{itemize}
			\item \textbf{Цепь} — последовательность прямоугольников, каждый из которых можно вложить в следующий.
			\item \textbf{Антицепь} — множество прямоугольников, ни один из которых нельзя вложить в другой.
		\end{itemize}
		
		Теорема Дилуорса утверждает, что минимальное количество цепей (последовательностей вложенных прямоугольников), на которые можно разбить все прямоугольники, равно максимальному размеру антицепи (множества попарно невложимых прямоугольников).
	\end{examplebox}
	
	\begin{notebook}{7.8: Историческая справка}
		Теорема названа в честь американского математика Роберта Дилуорса, который опубликовал её в 1950 году. Она является фундаментальным результатом теории частично упорядоченных множеств и имеет многочисленные применения в комбинаторике, теории графов и алгоритмах.
	\end{notebook}
	
	
	
	
	
	
	
	
	
	
	
	
	
	
	
	
	
	
	% ==================== РАЗДЕЛ 8: КОМБИНАТОРИКА ====================
	\newpage
	\section{Элементы комбинаторики}
	
	\subsection{Комбинаторное доказательство. Перестановки, размещения и сочетания. Факториальная система счисления (8 вопрос)}
	
	\subsubsection{Комбинаторное доказательство}
	
	\begin{definitionbox}{8.1: Комбинаторное доказательство}
		\textbf{Комбинаторное доказательство} — это способ нахождения мощности множества путем установления биекции с множеством известной мощности.\\
		
		Идея метода: чтобы доказать тождество $a = b$, мы придумываем задачу $X$, решением которой является и $a$, и $b$ (один и тот же ответ, получаемый разными способами). Преимущество данного способа — в почти полном отсутствии вычислений.
	\end{definitionbox}
	
	\begin{examplebox}{8.2: Пример комбинаторного доказательства}
		Докажем, что $C_n^k = C_n^{n-k}$.\\
		
		Придумываем задачу: $X$ — сколько существует способов выбрать без повторов $k$ объектов из $n$, когда порядок на выбранных объектах не важен.
		
		\begin{itemize}
			\item Левая часть $C_n^k$ — это решение задачи по определению числа сочетаний.
			\item Для правой части: выбрать $k$ объектов из $n$ — это то же самое, что выбрать $n-k$ объектов, которые мы \textbf{не} берем, и "выбросить" их; оставшиеся $k$ объектов образуют искомую выборку.
		\end{itemize}
		
		Обе части равенства являются решением одной и той же задачи, следовательно, они равны. Таким образом, тождество доказано.
	\end{examplebox}
	
	\begin{examplebox}{8.3: Задача о пончиках}
		Пусть у нас есть пончики 5 видов: с ванилью (В), шоколадом (Ш), карамелью (К), сгущенкой (С) и сахарной пудрой (П).\\
		
		\textbf{Задача 1:} Сколько способов купить \textbf{не более одного} пончика каждого вида?\\
		
		Установим биекцию между множеством способов купить пончики и множеством двоичных кортежей длины 5. Для кортежа $(0, 1, 0, 0, 1)$ скажем, что если на $i$-й позиции стоит 0, то мы купили пончик такого вида, а если 1 — то нет. Таким образом, каждому способу покупки соответствует уникальный двоичный кортеж, и наоборот.\\
		
		Количество двоичных кортежей длины 5 равно $2^5 = 32$. Следовательно, и способов купить пончики тоже 32.\\
		
		\textbf{Задача 2:} Хотим купить 12 пончиков, неважно какого вида. Сколько способов это сделать?\\
		
		Количество способов равно количеству двоичных кортежей длины 16, в которых ровно 4 единицы. \\
		
		Интерпретация: нули в кортеже обозначают пончики, причем нули, идущие до $i$-й единицы, отвечают за количество пончиков $i$-го вида. Например, кортеж
		\[ (0, 0, 1, 1, 0, 0, 0, 0, 1, 0, 0, 0, 1, 0, 0) \]
		соответствует: 2 ванильных, 0 шоколадных, 5 карамельных, 3 со сгущенкой, 2 с сахарной пудрой.\\
		
		Количество таких кортежей равно $C_{16}^4$, так как мы выбираем позиции для 4 единиц из 16 возможных.
	\end{examplebox}
	
	\subsubsection{Мощность множества всех подмножеств (комбинаторное доказательство)}
	
	\begin{theorembox}{8.4: Мощность булеана (комбинаторное доказательство)}
		Для любого конечного множества $A$ с $|A| = m$ верно $|2^A| = 2^m$.
	\end{theorembox}
	
	\myproof
	Пусть $A$ — произвольное конечное множество, $|A| = m$. Зафиксируем нумерацию элементов $A$, т.е. биекцию $\phi: A \to \{1, 2, \dots, m\}$.
	
	Рассмотрим множество $X = \{1, 2, \dots, m\}$. Для любого подмножества $S \subseteq X$ определим его \textbf{характеристический вектор} $\chi(S) \in \{0, 1\}^m$:
	\[ \chi(S) = (\chi(S)_1, \chi(S)_2, \dots, \chi(S)_m), \quad \text{где } \chi(S)_i = \begin{cases}
		1, & i \in S \\
		0, & i \notin S
	\end{cases} \]
	
	Отображение $\chi: 2^X \to \{0, 1\}^m$ является биекцией: каждому подмножеству соответствует уникальный двоичный кортеж, и каждый двоичный кортеж задает некоторое подмножество.
	
	Теперь построим биекцию $\psi: \{0, 1\}^m \to \{0, 1, \dots, 2^m-1\}$ следующим образом:
	\[ \psi(x_1, x_2, \dots, x_m) = \sum_{i=1}^m x_i \cdot 2^{m-i} \]
	т.е. интерпретируем двоичный кортеж как двоичную запись числа (старший бит — $x_1$).
	
	Это отображение также является биекцией, так как каждое число от $0$ до $2^m-1$ имеет единственное двоичное представление длины $m$ (возможно, с ведущими нулями).
	
	Таким образом, мы построили цепочку биекций:
	\[ A \xrightarrow{\phi} \{1,\dots,m\} \;\Rightarrow\; 2^A \xrightarrow{\text{биекция}} 2^{\{1,\dots,m\}} \xrightarrow{\chi} \{0,1\}^m \xrightarrow{\psi} \{0,\dots,2^m-1\} \]
	
	Следовательно, все эти множества равномощны, и
	\[ |2^A| = |\{0,1,\dots,2^m-1\}| = 2^m. \]
	\qed
	
	\subsubsection{Перестановки, размещения, сочетания}
	
	\begin{definitionbox}{8.5: Перестановка}
		\textbf{Перестановкой} множества $\{1, 2, \dots, n\}$ называется упорядоченный кортеж $\langle a_1, a_2, \dots, a_n \rangle$, удовлетворяющий условиям:
		\begin{enumerate}
			\item $\forall i \in 1:n \; a_i \in \{1, \dots, n\}$ (каждый элемент принадлежит множеству);
			\item $\forall i \neq j \; a_i \neq a_j$ (элементы не повторяются).
		\end{enumerate}
		
		\textbf{Количество перестановок} $P(n) = n!$ (читается: "эн факториал").
	\end{definitionbox}
	
	\begin{definitionbox}{8.6: Размещение}
		\textbf{Размещением} из $n$ элементов по $k$ ($k \le n$) называется упорядоченный набор из $k$ различных элементов, выбранных из $n$-элементного множества. Порядок элементов важен.\\
		
		\textbf{Количество размещений:}
		\[ A_n^k = \frac{n!}{(n-k)!} \]
		
		Перестановка является частным случаем размещения: $P(n) = A_n^n$.
	\end{definitionbox}
	
	\begin{examplebox}{8.7: Пример размещений}
		Пусть дано множество $\{1, 2, 3, 4, 5, 6\}$. Размещениями из 6 по 2 будут:
		\[ \langle 1,2 \rangle, \langle 1,3 \rangle, \langle 1,4 \rangle, \langle 1,5 \rangle, \langle 2,1 \rangle, \langle 2,3 \rangle, \dots, \langle 5,6 \rangle \]
		
		Всего их $A_6^2 = \frac{6!}{4!} = 6 \cdot 5 = 30$.
		
		Пример размещения из 6 по 4: $\langle 1, 3, 2, 5 \rangle$.
	\end{examplebox}
	
	\begin{definitionbox}{8.8: Сочетание}
		\textbf{Сочетанием} из $n$ элементов по $k$ ($k \le n$) называется набор из $k$ элементов, выбранных из $n$-элементного множества. \textbf{Порядок элементов не учитывается} — сочетания, отличающиеся только порядком следования элементов, считаются одинаковыми.\\
		
		\textbf{Количество сочетаний:}
		\[ C_n^k = \binom{n}{k} = \frac{n!}{k!(n-k)!} \]
	\end{definitionbox}
	
	\begin{remarkbox}{8.9: Связь между размещениями и сочетаниями}
		Размещение можно получить, сначала выбрав $k$ элементов (сочетание), а затем упорядочив их ($k!$ способов):
		\[ A_n^k = C_n^k \cdot k! \implies C_n^k = \frac{A_n^k}{k!} = \frac{n!}{k!(n-k)!} \]
	\end{remarkbox}
	
	\begin{center}
		\begin{tabular}{|l|l|l|}
			\hline
			\textbf{Понятие} & \textbf{Порядок важен?} & \textbf{Формула} \\ \hline
			Перестановка $P(n)$ & да & $n!$ \\ \hline
			Размещение $A_n^k$ & да & $\frac{n!}{(n-k)!}$ \\ \hline
			Сочетание $C_n^k$ & нет & $\frac{n!}{k!(n-k)!}$ \\ \hline
		\end{tabular}
		\captionof{table}{Основные комбинаторные объекты}
	\end{center}
	
	\subsubsection{Факториальная система счисления}
	
	\begin{definitionbox}{8.10: Факториальная система счисления}
		В \textbf{факториальной системе счисления} основаниями являются последовательности факториалов. Любое натуральное число $x$ представляется в виде:
		\[ x = \sum_{k=1}^n d_k \cdot k! \]
		где $0 \le d_k \le k$ для всех $k$, и $d_n \neq 0$ (кроме случая $x=0$).
		
		Запись числа в факториальной системе: $(d_n d_{n-1} \dots d_2 d_1)_{!}$.
	\end{definitionbox}
	
	\begin{examplebox}{8.11: Пример представления в факториальной системе}
		Число 28 в десятичной системе представляется в факториальной как $1020_{!}$, так как:
		\[ 1 \cdot 4! + 0 \cdot 3! + 2 \cdot 2! + 0 \cdot 1! = 1 \cdot 24 + 0 \cdot 6 + 2 \cdot 2 + 0 \cdot 1 = 24 + 4 = 28 \]
	\end{examplebox}
	
	\subsubsection{Применение факториальной системы: генерация перестановок}
	
	Факториальная система счисления используется для перебора всех возможных перестановок множества. Рассмотрим алгоритм:
	
	\begin{definitionbox}{8.12: Ключ перестановки (вектор инверсий)}
		Для перестановки $a = \langle a_1, a_2, \dots, a_n \rangle$ определим её \textbf{ключ} (вектор инверсий) $T(a) = (t_1, t_2, \dots, t_n)$, где $t_i$ — количество элементов перестановки, стоящих после $i$-го элемента, но при этом меньших его.\\
		
		Иными словами, $t_i$ — число инверсий, которые образует $i$-й элемент с последующими.
	\end{definitionbox}
	
	\begin{importantbox}{8.13: Свойства ключа перестановки}
		\begin{itemize}
			\item Всегда $0 \le t_i \le n-i$ (так как после $i$-го элемента может быть не более $n-i$ элементов).
			\item Кортеж $(t_1, t_2, \dots, t_n)$, рассматриваемый как число в факториальной системе счисления (с весами $(n-1)!, (n-2)!, \dots, 0!$), задает уникальный номер перестановки от $0$ до $n!-1$.
		\end{itemize}
	\end{importantbox}
	
	\begin{examplebox}{8.14: Пример построения ключа}
		Рассмотрим перестановку:
		\[ a = \langle 3, 8, 7, 6, 2, 4, 9, 5, 1 \rangle \]
		
		Вычисляем ключ:
		\begin{align*}
			t_1 &= \text{после 3: элементы } 8,7,6,2,4,9,5,1 \text{ — меньшие 3: } \{2,1\} \Rightarrow t_1 = 2 \\
			t_2 &= \text{после 8: } 7,6,2,4,9,5,1 \text{ — меньшие 8: } \{7,6,2,4,5,1\} \Rightarrow t_2 = 6 \\
			t_3 &= \text{после 7: } 6,2,4,9,5,1 \text{ — меньшие 7: } \{6,2,4,5,1\} \Rightarrow t_3 = 5 \\
			t_4 &= \text{после 6: } 2,4,9,5,1 \text{ — меньшие 6: } \{2,4,5,1\} \Rightarrow t_4 = 4 \\
			t_5 &= \text{после 2: } 4,9,5,1 \text{ — меньшие 2: } \{1\} \Rightarrow t_5 = 1 \\
			t_6 &= \text{после 4: } 9,5,1 \text{ — меньшие 4: } \{1\} \Rightarrow t_6 = 1 \\
			t_7 &= \text{после 9: } 5,1 \text{ — меньшие 9: } \{5,1\} \Rightarrow t_7 = 2 \\
			t_8 &= \text{после 5: } 1 \text{ — меньшие 5: } \{1\} \Rightarrow t_8 = 1 \\
			t_9 &= \text{после 1: } \emptyset \Rightarrow t_9 = 0
		\end{align*}
		
		Получаем ключ: $T(a) = (2, 6, 5, 4, 1, 1, 2, 1, 0)$.
		
		Если интерпретировать этот ключ как число в факториальной системе счисления (игнорируя последнюю цифру, которая всегда 0), получим $26541121_{!}$.
	\end{examplebox}
	
	\begin{algorithmbox}{8.15: Генерация следующей перестановки}
		Факториальная система счисления позволяет легко перебирать перестановки в лексикографическом порядке:
		
		\begin{enumerate}
			\item Представляем ключ текущей перестановки как число в факториальной системе.
			\item Прибавляем к этому числу 1 (выполняя сложение в факториальной системе).
			\item Полученное новое число преобразуем обратно в ключ.
			\item По новому ключу восстанавливаем перестановку.
		\end{enumerate}
	\end{algorithmbox}
	
	\begin{examplebox}{8.16: Пример генерации следующей перестановки}
		Для перестановки из примера 8.14 имеем ключ-число $26541121_{!}$.
		
		Прибавляем 1:
		\[ 26541121_{!} + 1 = 26541200_{!} \]
		
		Это соответствует новому ключу:
		\[ T'(a) = (2, 6, 5, 4, 1, 2, 0, 0, 0) \]
		
		По этому ключу можно восстановить новую перестановку:
		\[ \langle 3, 8, 7, 6, 2, 5, 1, 4, 9 \rangle \]
	\end{examplebox}
	
	\begin{remarkbox}{8.17: Преимущества факториальной системы}
		\begin{itemize}
			\item Позволяет эффективно генерировать перестановки в лексикографическом порядке.
			\item Дает возможность нумеровать все $n!$ перестановок числами от $0$ до $n!-1$.
			\item Используется в комбинаторных алгоритмах для перебора и оптимизации.
		\end{itemize}
	\end{remarkbox}
	
	\begin{center}
		\begin{tikzpicture}[scale=0.9]
			
			% Ось
			\draw[thick, blue] (0,0) -- (10,0);
			
			% Деления
			\foreach \x in {0,1,...,9} {
				\draw (\x,0.15) -- (\x,-0.15);
			}
			
			% Подписи крайних значений
			\node[below] at (0,-0.15) {$0$};
			\node[below] at (9,-0.15) {$n!-1$};
			
			% Подпись оси
			\node at (5,-0.8) {Номера перестановок в диапазоне $0 \dots n!-1$};
			
			% Стрелка перехода
			\draw[->, red, thick] (2,0.6) -- (4,0.6);
			\node[above] at (3,0.6) {+1 в факториальной системе};
			
			% Пример перестановок
			\node[draw, rounded corners, fill=green!10] 
			at (2,1.6) {\texttt{26541121}};
			
			\node[draw, rounded corners, fill=green!10] 
			at (4,1.6) {\texttt{26541200}};
			
		\end{tikzpicture}
		\captionof{figure}{Итерация по перестановкам через факториальную систему счисления}
	\end{center}
		
	
	
	
	
	
	
	
	
	
	
	
	
	
	
	
	
	
	
	% ==================== РАЗДЕЛ 9: ПЕРЕСТАНОВКИ И КОЭФФИЦИЕНТЫ ====================
	\newpage
	\section{Способы нумерации перестановок и перечисления}
	
	\subsection{Способы нумерации перестановок и перечисления. Биномиальные и мультиномиальные коэффициенты и их свойства (9 вопрос)}
	
	\subsubsection{Нумерация перестановок с помощью факториальной системы счисления}
	
	В предыдущем разделе мы познакомились с факториальной системой счисления. Теперь рассмотрим, как с её помощью можно нумеровать перестановки и переходить от одной перестановки к следующей.
	
	\begin{definitionbox}{9.1: Нумерация перестановок}
		Мы построили биекцию из множества всех перестановок $\{1,2,\dots,n\}$ через множество чисел, имеющих не более $n$ знаков в факториальной системе счисления, в множество $\{0, 1, \dots, n!-1\}$.
		
		Эта биекция устроена следующим образом:
		\begin{enumerate}
			\item Каждой перестановке $a$ сопоставляется её ключ (вектор инверсий) $T(a) = (t_1, t_2, \dots, t_n)$.
			\item Ключ интерпретируется как число в факториальной системе счисления (первые $n-1$ цифр, так как $t_n$ всегда равно 0).
			\item Это число и есть номер перестановки.
		\end{enumerate}
	\end{definitionbox}
	
	\begin{examplebox}{9.2: Пример нумерации}
		Рассмотрим перестановку:
		\[ a = \langle 3, 8, 7, 6, 2, 4, 9, 5, 1 \rangle \]
		
		Её ключ (вычисленный ранее): $T(a) = (2, 6, 5, 4, 1, 1, 2, 1, 0)$.
		
		В факториальной системе счисления этому ключу соответствует число:
		\[ 26541121_{!} \]
		
		Это и есть номер данной перестановки среди всех $9! = 362880$ перестановок.
	\end{examplebox}
	
	\subsubsection{Переход к следующей перестановке}
	
	\begin{importantbox}{9.3: Алгоритм получения следующей перестановки}
		Чтобы найти перестановку, следующую за данной в лексикографическом порядке:
		
		\begin{enumerate}
			\item Представляем ключ текущей перестановки как число в факториальной системе.
			\item Прибавляем к этому числу $1$ (выполняя сложение в факториальной системе).
			\item Полученное новое число преобразуем обратно в ключ $T(b)$.
			\item По новому ключу восстанавливаем перестановку $b$.
		\end{enumerate}
	\end{importantbox}
	
	\begin{examplebox}{9.4: Пример перехода к следующей перестановке}
		Для перестановки из примера 9.2 имеем число $26541121_{!}$.
		
		Прибавляем 1 в факториальной системе:
		\[ 26541121_{!} + 1 = 26541200_{!} \]
		
		Это соответствует новому ключу:
		\[ T(b) = (2, 6, 5, 4, 1, 2, 0, 0, 0) \]
		
		По этому ключу восстанавливаем новую перестановку:
		\[ b = \langle 3, 8, 7, 6, 2, 5, 1, 4, 9 \rangle \]
	\end{examplebox}
	
	\begin{center}
		\begin{tikzpicture}[scale=0.9]
			\draw[thick, blue] (0,0) rectangle (13,2);
			
			% Верхняя строка
			\node at (2,1.4) {$\langle 3,8,7,6,2,4,9,5,1 \rangle$};
			\node at (5.5,1.4) {$\xrightarrow{\text{ключ}}$};
			\node at (8,1.4) {$26541121_{!}$};
			\node at (10,1.4) {$\xrightarrow{+1}$};
			\node at (11.5,1.4) {$26541200_{!}$};
			
			% Восстановление
			\node at (6.5,-0.4) {Восстановление из $T=(2,6,5,4,1,2,0,0,0)$};
			
			\draw[thick, blue] (0,-2.2) rectangle (13,-0.8);
			
			\node at (2,-1.5) {$L=[1,\dots,9]$};
			
			\node at (5.5,-1.5) {$\xrightarrow{\text{индекс }T_i}$};
			
			\node at (10,-1.5) {$\langle 3,8,7,6,2,5,1,4,9 \rangle$};
		\end{tikzpicture}
		\captionof{figure}{Переход к следующей перестановке}
	\end{center}
	
	\begin{theorembox}{9.5: Структура изменения перестановки}
		При переходе к следующей перестановке изменяется так называемый "убывающий хвост". Если $l$ — длина убывающего хвоста (максимальный суффикс, в котором элементы идут в убывающем порядке), и $T(a)_{n-l} = k$, то:
		
		\begin{itemize}
			\item Последние $l+1$ элементов ключа $T(b)$ будут иметь вид $(k+1, 0, 0, \dots, 0)$.
			\item Первые $n-l-1$ элементов перестановки остаются неизменными.
			\item На позиции $n-l$ в перестановке $b$ будет стоять минимальный элемент $x$ из множества $\{a_j : j \ge n-l\}$, такой что $x > k$.
			\item Оставшиеся $l$ элементов из множества $\{a_j : j \ge n-l\} \setminus \{x\}$ располагаются в конце перестановки $b$ в порядке возрастания.
		\end{itemize}
	\end{theorembox}
	
	\myremark{Пояснение:} Берем перестановку, строим ключ $T$ по принципу "сколько элементов меньше данного стоит правее от него". Переводим в факториальную систему — это номер текущей перестановки. Чтобы найти перестановку на $n$ позиций вперед/назад, прибавляем/вычитаем это число в факториальной системе и получаем ключ, по которому восстанавливаем нужную перестановку.
	
	\subsubsection{Биномиальные коэффициенты}
	
	\begin{definitionbox}{9.6: Биномиальные коэффициенты}
		\textbf{Биномиальные коэффициенты} — это числа $C_n^k = \binom{n}{k}$, выражающие количество способов выбрать $k$ элементов из $n$ без учёта порядка.\\
		
		Название связано с тем, что они появляются в формуле бинома Ньютона:
		\[ (x + y)^n = \sum_{k=0}^n \binom{n}{k} x^{n-k} y^k \]
	\end{definitionbox}
	
	\begin{propertybox}{9.7: Основные свойства биномиальных коэффициентов}
		
		\textbf{Свойство 1 (Симметричность):}
		\[ \binom{n}{k} = \binom{n}{n-k} \]
		
		\myproof
		Выбрав подмножество размера $k$ из $n$-элементного множества, мы однозначно определяем его дополнение — подмножество размера $n-k$. Следовательно, количество способов выбрать $k$ элементов равно количеству способов выбрать $n-k$ элементов. \qed \\
		
		\textbf{Свойство 2 (Рекуррентное соотношение):}
		\[ \binom{n}{k} = \binom{n-1}{k-1} + \binom{n-1}{k} \]
		
		\myproof
		Зафиксируем некоторый элемент $a$ исходного множества из $n$ элементов. Рассмотрим все подмножества размера $k$:
		\begin{itemize}
			\item Подмножества, содержащие $a$: нужно выбрать оставшиеся $k-1$ элементов из $n-1$ оставшихся — $\binom{n-1}{k-1}$ способов.
			\item Подмножества, не содержащие $a$: нужно выбрать все $k$ элементов из $n-1$ оставшихся — $\binom{n-1}{k}$ способов.
		\end{itemize}
		Эти два класса не пересекаются и покрывают все подмножества, поэтому общее количество равно сумме. \qed \\
		
		\textbf{Аналитическое доказательство свойства 2:}
		\[ \binom{n+1}{k} = \binom{n}{k} + \binom{n}{k-1} = \frac{n!}{k!(n-k)!} + \frac{n!}{(k-1)!(n-k+1)!} = \]
		\[ = \frac{n! \cdot (n-k+1)}{k!(n-k+1)!} + \frac{n! \cdot k}{k!(n-k+1)!} = \frac{n!(n-k+1+k)}{k!(n-k+1)!} = \]
		\[ = \frac{n!(n+1)}{k!((n+1)-k)!} = \frac{(n+1)!}{k!((n+1)-k)!} = \binom{n+1}{k} \]
	\end{propertybox}
	
	\begin{examplebox}{9.8: Пример комбинаторного тождества}
		Рассмотрим тождество:
		\[ \binom{3n}{n} = \sum_{r=0}^n \binom{n}{r} \binom{2n}{n-r} \]
		
		\textbf{Комбинаторная интерпретация:}
		Представим, что у нас есть колода из $3n$ карт: $n$ карт заклинаний и $2n$ карт существ. Нужно выбрать $n$ карт для составления колоды.
		
		Зафиксируем $r$ — количество заклинаний в колоде ($r$ может быть от $0$ до $n$). Тогда:
		\begin{itemize}
			\item Выбрать $r$ заклинаний из $n$ можно $\binom{n}{r}$ способами.
			\item Выбрать оставшиеся $n-r$ карт из $2n$ существ можно $\binom{2n}{n-r}$ способами.
		\end{itemize}
		По правилу произведения для фиксированного $r$ получаем $\binom{n}{r} \binom{2n}{n-r}$ способов. Суммируя по всем возможным $r$, получаем общее количество способов выбрать $n$ карт из $3n$, которое равно $\binom{3n}{n}$.
	\end{examplebox}
	
	\begin{remarkbox}{9.9: Альтернативная формулировка}
		Ту же идею можно проиллюстрировать на более простом примере:
		
		\textit{Задача:} Есть 30 шаров: 20 белых и 10 черных. Сколько способов выбрать 10 шаров любого цвета?
		
		Решение: фиксируем $r$ — количество черных шаров в выборке ($r = 0..10$). Тогда способов выбрать $r$ черных из 10 — $C_{10}^r$, а оставшиеся $10-r$ белых из 20 — $C_{20}^{10-r}$. Суммируем по $r$:
		\[ \sum_{r=0}^{10} C_{10}^r C_{20}^{10-r} = C_{30}^{10} \]
	\end{remarkbox}
	
	\subsubsection{Мультиномиальные коэффициенты}
	
	\begin{definitionbox}{9.10: Мультиномиальные (полиномиальные) коэффициенты}
		\textbf{Мультиномиальный коэффициент} — это обобщение биномиального коэффициента на случай разбиения множества на несколько групп. Он обозначается как:
		\[ \binom{n}{n_1, n_2, \dots, n_k} = \frac{n!}{n_1! \, n_2! \, \cdots \, n_k!} \]
		где $n_1 + n_2 + \dots + n_k = n$.
		
		Этот коэффициент выражает количество способов разбить $n$-элементное множество на $k$ групп, содержащих $n_1, n_2, \dots, n_k$ элементов соответственно (порядок групп важен).
	\end{definitionbox}
	
	\begin{examplebox}{9.11: Пример мультиномиального коэффициента}
		Сколькими способами можно разбить 10 студентов на 3 группы по 5, 3 и 2 человека?
		
		Ответ: $\binom{10}{5, 3, 2} = \frac{10!}{5! \cdot 3! \cdot 2!} = \frac{3628800}{120 \cdot 6 \cdot 2} = \frac{3628800}{1440} = 2520$ способов.
	\end{examplebox}
	
	\begin{theorembox}{9.12: Мультиномиальная теорема}
		Мультиномиальные коэффициенты появляются в разложении степени суммы нескольких слагаемых:
		\[ (x_1 + x_2 + \dots + x_k)^n = \sum_{\substack{n_1+\dots+n_k=n \\ n_i \ge 0}} \binom{n}{n_1, n_2, \dots, n_k} x_1^{n_1} x_2^{n_2} \cdots x_k^{n_k} \]
		
		При $k=2$ получается бином Ньютона.
	\end{theorembox}
	
	\myremark{Замечание:} В презентациях и на лекциях мультиномиальные коэффициенты подробно не рассматривались, поэтому для экзамена достаточно знать их определение и связь с биномиальными коэффициентами. Если потребуются дополнительные детали, можно сослаться на отсутствие этого материала в официальных слайдах.
	
	\subsubsection{Таблица свойств биномиальных коэффициентов}
	
	\begin{center}
		\begin{tabular}{|l|l|l|}
			\hline
			\textbf{Свойство} & \textbf{Формула} & \textbf{Комбинаторный смысл} \\ \hline
			Симметричность & $\binom{n}{k} = \binom{n}{n-k}$ & Дополнение подмножества \\ \hline
			Рекуррентность & $\binom{n}{k} = \binom{n-1}{k-1} + \binom{n-1}{k}$ & Выбор с фиксированным элементом \\ \hline
			Граничные значения & $\binom{n}{0} = \binom{n}{n} = 1$ & Пустое множество и всё множество \\ \hline
			Сумма по $k$ & $\sum_{k=0}^n \binom{n}{k} = 2^n$ & Все подмножества \\ \hline
			Сумма с чередованием & $\sum_{k=0}^n (-1)^k \binom{n}{k} = 0$ (при $n>0$) & Принцип включения-исключения \\ \hline
		\end{tabular}
		\captionof{table}{Основные свойства биномиальных коэффициентов}
	\end{center}
	
	\begin{notebook}{9.13: Дополнительные источники}
		\begin{itemize}
			\item Подробнее о перестановках и факториальной системе: \textit{Романовский В.Ф. "Дискретная математика"}, стр. 14–20.
			\item О биномиальных и мультиномиальных коэффициентах: там же, стр. 53–60.
			\item Онлайн-ресурсы: ссылки на материалы по комбинаторике и теории перестановок.
		\end{itemize}
	\end{notebook}
		
	\begin{center}
		\begin{tikzpicture}
			
			\begin{axis}[
				width=12cm,
				height=8cm,
				xlabel={$k$},
				ylabel={$\binom{6}{k}$},
				xtick={0,1,2,3,4,5,6},
				ytick={0,5,10,15,20},
				ymin=0,
				ymax=22,
				grid=both
				]
				
				\addplot[
				only marks,
				mark=*,
				red
				] coordinates {
					(0,1)
					(1,6)
					(2,15)
					(3,20)
					(4,15)
					(5,6)
					(6,1)
				};
				
			\end{axis}
			
		\end{tikzpicture}
		\captionof{figure}{Биномиальные коэффициенты $\binom{6}{k}$}
	\end{center}
		
	
	
	
	
	
	
	
	
	
	% ==================== РАЗДЕЛ 10: ФОРМУЛА ВКЛЮЧЕНИЙ-ИСКЛЮЧЕНИЙ ====================
	\newpage
	\section{Формула включений-исключений}
	
	\subsection{Формула включений-исключений (10 вопрос)}
	
	\begin{definitionbox}{10.1: Формула включений-исключений}
		\textbf{Формула включений-исключений} — это формула, позволяющая определить мощность объединения конечного числа конечных множеств, которые в общем случае могут пересекаться друг с другом.
		
		Для любых конечных множеств $A_1, A_2, \dots, A_n$ справедливо равенство:
		\[ |A_1 \cup A_2 \cup \dots \cup A_n| = \sum_{i} |A_i| - \sum_{i < j} |A_i \cap A_j| + \sum_{i < j < k} |A_i \cap A_j \cap A_k| - \dots + (-1)^{n+1} |A_1 \cap A_2 \cap \dots \cap A_n| \]
		
		Или в компактной записи:
		\[ \left|\bigcup_{i=1}^n A_i\right| = \sum_{j=1}^n (-1)^{j+1} \sum_{1 \le i_1 < i_2 < \dots < i_j \le n} |A_{i_1} \cap A_{i_2} \cap \dots \cap A_{i_j}| \]
	\end{definitionbox}
	
	\myremark{Суть формулы:}
	\begin{itemize}
		\item Сначала суммируем мощности всех множеств ($n$ слагаемых);
		\item Затем вычитаем мощности всех попарных пересечений ($C_n^2$ слагаемых);
		\item Затем прибавляем мощности всех тройных пересечений ($C_n^3$ слагаемых);
		\item И так далее, чередуя знаки: \textbf{плюс} при нечетном количестве множеств в пересечении, \textbf{минус} при четном;
		\item Последнее слагаемое — мощность пересечения всех $n$ множеств со знаком $(-1)^{n+1}$.
	\end{itemize}
	
	\subsubsection{Пример для двух множеств}
	
	\begin{examplebox}{10.2: Два множества}
		Для двух множеств формула принимает простой вид:
		\[ |A \cup B| = |A| + |B| - |A \cap B| \]
		
		\textbf{Пояснение:} Когда мы складываем мощности множеств $A$ и $B$, их пересекающиеся элементы учитываются два раза. Поэтому для получения правильной мощности объединения необходимо вычесть мощность пересечения.\\
		
		\textbf{Пример:} 
		\begin{align*}
			A &= \{1, 2, 3, 4, 5\} \quad |A| = 5 \\
			B &= \{4, 5, 6, 7, 8\} \quad |B| = 5 \\
			A \cap B &= \{4, 5\} \quad |A \cap B| = 2 \\
			A \cup B &= \{1, 2, 3, 4, 5, 6, 7, 8\} \quad |A \cup B| = 8
		\end{align*}
		Проверка: $5 + 5 - 2 = 8$ — верно.
	\end{examplebox}
	
	\subsubsection{Вспомогательная лемма}
	
	Для доказательства формулы нам потребуется следующая лемма.
	
	\begin{lemmabox}{10.3: Лемма о сумме биномиальных коэффициентов}
		Для любого натурального $m$ справедливо равенство:
		\[ C_1^m - C_2^m + C_3^m - \dots + (-1)^{m+1} C_m^m = 1 \]
		т.е.
		\[ \sum_{k=1}^m (-1)^{k+1} \binom{m}{k} = 1 \]
	\end{lemmabox}
	
	\myproof
	Воспользуемся биномом Ньютона:
	\[ 0 = (1 - 1)^m = \sum_{k=0}^m \binom{m}{k} \cdot 1^{m-k} \cdot (-1)^k = 1 + \sum_{k=1}^m (-1)^k \binom{m}{k} \]
	Перенесем 1 в левую часть:
	\[ -1 = \sum_{k=1}^m (-1)^k \binom{m}{k} \]
	Умножим обе части на $-1$:
	\[ 1 = \sum_{k=1}^m (-1)^{k+1} \binom{m}{k} = C_1^m - C_2^m + C_3^m - \dots + (-1)^{m+1} C_m^m \]
	\qed
	
	\subsubsection{Доказательство формулы включений-исключений}
	
	\myproof \textbf{(Методом подсчёта вклада элемента)}
	
	Возьмём произвольный элемент $x \in A_1 \cup A_2 \cup \dots \cup A_n$. Покажем, что в правой части формулы (3) этот элемент учитывается ровно один раз.\\
	
	Предположим, что $x$ принадлежит пересечению ровно $m$ наших множеств. Не ограничивая общности, можно считать, что $x$ принадлежит множествам $A_1, A_2, \dots, A_m$ и не принадлежит множествам $A_{m+1}, \dots, A_n$ (перенумеровав множества при необходимости).\\
	
	Теперь посчитаем, сколько раз элемент $x$ появляется в каждом слагаемом правой части:
	
	\begin{itemize}
		\item В первой сумме $\sum_i |A_i|$ элемент $x$ посчитан $m = C_1^m$ раз (в слагаемых $|A_1|, |A_2|, \dots, |A_m|$).
		\item Во второй сумме $\sum_{i<j} |A_i \cap A_j|$ элемент $x$ посчитан $C_2^m$ раз — это количество попарных пересечений множеств из $\{A_1, \dots, A_m\}$, в которые входит $x$.
		\item В третьей сумме $\sum_{i<j<k} |A_i \cap A_j \cap A_k|$ элемент $x$ посчитан $C_3^m$ раз.
		\item ...
		\item В $m$-й сумме $\sum_{i_1 < i_2 < \dots < i_m} |A_{i_1} \cap A_{i_2} \cap \dots \cap A_{i_m}|$ элемент $x$ посчитан $1 = C_m^m$ раз (он входит только в пересечение $A_1 \cap A_2 \cap \dots \cap A_m$).
		\item Суммы, содержащие $m+1$ и более пересечений, не учитывают элемент $x$, поскольку $x$ не входит в пересечение более чем $m$ множеств.
	\end{itemize}
	
	Знаки слагаемых чередуются: первая сумма со знаком $+$, вторая со знаком $-$, третья со знаком $+$ и т.д. Таким образом, общий вклад элемента $x$ в правую часть равен:
	\[ C_1^m - C_2^m + C_3^m - \dots + (-1)^{m+1} C_m^m \]
	
	Согласно лемме 10.3, это выражение равно $1$. Следовательно, каждый элемент объединения учитывается ровно один раз, что и доказывает формулу. \qed
	
	\subsubsection{Доказательство по индукции}
	
	Приведём также альтернативное доказательство методом математической индукции.
	
	\myproof \textbf{(По индукции)}
	
	\textbf{База индукции:} $n = 1$. Формула принимает вид $|A_1| = |A_1|$ — очевидно верно.
	
	\textbf{Шаг индукции:} Предположим, что формула верна для $n = m$. Докажем её для $n = m+1$.
	
	Рассмотрим объединение $m+1$ множества:
	\[ |A_1 \cup \dots \cup A_m \cup A_{m+1}| = |(A_1 \cup \dots \cup A_m) \cup A_{m+1}| \]
	
	По формуле для двух множеств:
	\[ = |A_1 \cup \dots \cup A_m| + |A_{m+1}| - |(A_1 \cup \dots \cup A_m) \cap A_{m+1}| \]
	
	Используя дистрибутивность пересечения относительно объединения:
	\[ (A_1 \cup \dots \cup A_m) \cap A_{m+1} = (A_1 \cap A_{m+1}) \cup (A_2 \cap A_{m+1}) \cup \dots \cup (A_m \cap A_{m+1}) \]
	
	Теперь применим предположение индукции к $|A_1 \cup \dots \cup A_m|$ и к $|(A_1 \cap A_{m+1}) \cup \dots \cup (A_m \cap A_{m+1})|$:
	
	\[ |A_1 \cup \dots \cup A_m| = \sum_{j=1}^m (-1)^{j+1} \sum_{1 \le i_1 < \dots < i_j \le m} |A_{i_1} \cap \dots \cap A_{i_j}| \]
	
	\[ |(A_1 \cap A_{m+1}) \cup \dots \cup (A_m \cap A_{m+1})| = \sum_{j=1}^m (-1)^{j+1} \sum_{1 \le i_1 < \dots < i_j \le m} |A_{i_1} \cap \dots \cap A_{i_j} \cap A_{m+1}| \]
	
	Подставляем эти выражения:
	\begin{align*}
		|A_1 \cup \dots \cup A_{m+1}| &= \left[ \sum_{j=1}^m (-1)^{j+1} \sum_{i_1<\dots<i_j} |A_{i_1} \cap \dots \cap A_{i_j}| \right] + |A_{m+1}| \\
		&\quad - \left[ \sum_{j=1}^m (-1)^{j+1} \sum_{i_1<\dots<i_j} |A_{i_1} \cap \dots \cap A_{i_j} \cap A_{m+1}| \right]
	\end{align*}
	
	Заметим, что второе слагаемое $|A_{m+1}|$ можно рассматривать как сумму по $j=1$ с множеством $\{A_{m+1}\}$. Перегруппируем слагаемые по количеству множеств в пересечении:
	
	\begin{itemize}
		\item Для пересечений, не содержащих $A_{m+1}$, они остаются как есть с индексами от $1$ до $j$.
		\item Для пересечений, содержащих $A_{m+1}$, они появляются из последней суммы, причём знак меняется: если в последней сумме был знак $(-1)^{j+1}$, то после вычитания знак становится $(-1)^{j+2}$.
	\end{itemize}
	
	В результате получаем сумму по всем непустым подмножествам $\{i_1, \dots, i_j\}$ множества $\{1, 2, \dots, m+1\}$:
	\begin{itemize}
		\item Если подмножество не содержит $m+1$, то вклад равен $(-1)^{j+1} |A_{i_1} \cap \dots \cap A_{i_j}|$.
		\item Если подмножество содержит $m+1$, то $m+1$ — максимальный индекс, и вклад равен $(-1)^{j+2} |A_{i_1} \cap \dots \cap A_{i_{j-1}} \cap A_{m+1}| = (-1)^{(j+1)+1} |A_{i_1} \cap \dots \cap A_{i_{j-1}} \cap A_{m+1}|$.
	\end{itemize}
	
	Объединяя, получаем требуемую формулу для $m+1$ множеств. \qed
	
	\subsubsection{Примеры применения}
	
	\begin{examplebox}{10.4: Три множества}
		Для трёх множеств формула принимает вид:
		\[ |A \cup B \cup C| = |A| + |B| + |C| - |A \cap B| - |A \cap C| - |B \cap C| + |A \cap B \cap C| \]
		
		\textbf{Пример:} 
		\begin{align*}
			A &= \{1, 2, 3, 4\} \\
			B &= \{3, 4, 5, 6\} \\
			C &= \{4, 5, 6, 7\}
		\end{align*}
		
		Вычислим необходимые мощности:
		\begin{align*}
			|A| &= 4, |B| = 4, |C| = 4 \\
			|A \cap B| &= |\{3,4\}| = 2 \\
			|A \cap C| &= |\{4\}| = 1 \\
			|B \cap C| &= |\{4,5,6\}| = 3 \\
			|A \cap B \cap C| &= |\{4\}| = 1
		\end{align*}
		
		Тогда:
		\[ |A \cup B \cup C| = 4 + 4 + 4 - 2 - 1 - 3 + 1 = 7 \]
		
		Проверим прямым подсчётом:
		\[ A \cup B \cup C = \{1, 2, 3, 4, 5, 6, 7\} \quad |A \cup B \cup C| = 7 \]
	\end{examplebox}
	
	\begin{examplebox}{10.5: Задача о беспорядках}
		Сколько существует перестановок чисел $1, 2, \dots, n$, в которых ни один элемент не стоит на своём месте (т.е. $\sigma(i) \neq i$ для всех $i$)? Такие перестановки называются \textbf{беспорядками}.
		
		Обозначим через $A_i$ множество перестановок, в которых элемент $i$ стоит на своём месте. Тогда:
		\begin{itemize}
			\item $|A_i| = (n-1)!$ (фиксируем $i$ на месте, остальные переставляем произвольно)
			\item $|A_i \cap A_j| = (n-2)!$ (фиксируем два элемента)
			\item $|A_{i_1} \cap \dots \cap A_{i_k}| = (n-k)!$
		\end{itemize}
		
		Искомое число беспорядков $!n$ (субфакториал) равно общему числу перестановок $n!$ минус те, в которых хотя бы один элемент на своём месте. По формуле включений-исключений:
		\begin{align*}
			!n &= n! - \sum_i |A_i| + \sum_{i<j} |A_i \cap A_j| - \sum_{i<j<k} |A_i \cap A_j \cap A_k| + \dots + (-1)^n |A_1 \cap \dots \cap A_n| \\
			&= n! - C_n^1 (n-1)! + C_n^2 (n-2)! - C_n^3 (n-3)! + \dots + (-1)^n C_n^n 0! \\
			&= n! \sum_{k=0}^n \frac{(-1)^k}{k!}
		\end{align*}
		
		Например, для $n=4$:
		\[ !4 = 4! \left(1 - 1 + \frac{1}{2} - \frac{1}{6} + \frac{1}{24}\right) = 24 \cdot \frac{9}{24} = 9 \]
	\end{examplebox}
	
	\begin{center}
		\begin{tikzpicture}
			\draw[thick, blue] (0,0) circle (2cm);
			\node at (0,1) {$A$};
			\draw[thick, red] (2,0) circle (2cm);
			\node at (2,1) {$B$};
			\draw[thick, green!50!black] (1,-1.5) circle (2cm);
			\node at (1,-2.5) {$C$};
			
			\fill[yellow, opacity=0.3] (0,0) circle (2cm);
			\fill[yellow, opacity=0.3] (2,0) circle (2cm);
			\fill[yellow, opacity=0.3] (1,-1.5) circle (2cm);
			
			\node at (1,0.5) {$A \cap B \cap C$};
		\end{tikzpicture}
		\captionof{figure}{Иллюстрация пересечения трёх множеств}
	\end{center}
	
	\begin{notebook}{10.6: Обобщение}
		Формула включений-исключений обобщается на случай, когда вместо мощности множеств используется мера (например, вероятность). В теории вероятностей она позволяет вычислять вероятность объединения событий:
		\[ P\left(\bigcup_{i=1}^n A_i\right) = \sum_{j=1}^n (-1)^{j+1} \sum_{1 \le i_1 < \dots < i_j \le n} P(A_{i_1} \cap \dots \cap A_{i_j}) \]
	\end{notebook}
	
	
	
	
	
	
	
	
	
	
	
	
	
	
	% ==================== РАЗДЕЛ 11: ЧИСЛА БЕЛЛА И СТИРЛИНГА ====================
	\newpage
	\section{Числа Белла и Стирлинга}
	
	\subsection{Числа Белла и Стирлинга (11 вопрос)}
	
	\subsubsection{Числа Белла}
	
	\begin{definitionbox}{11.1: Числа Белла}
		\textbf{Числом Белла} \(B(n)\) называется количество всех возможных разбиений конечного множества мощности \(n\).\\
		
		Иными словами, если \(|A| = n\), то \(B(n)\) — это число способов представить \(A\) в виде объединения непустых попарно непересекающихся подмножеств (порядок подмножеств не важен).
	\end{definitionbox}
	
	\begin{examplebox}{11.2: Пример для n=3}
		Рассмотрим множество \(A = \{1,2,3\}\). Все его разбиения:
		\begin{itemize}
			\item \(\{\{1,2,3\}\}\) — одно подмножество
			\item \(\{\{1\}, \{2,3\}\}\), \(\{\{2\}, \{1,3\}\}\), \(\{\{3\}, \{1,2\}\}\) — разбиения на два подмножества
			\item \(\{\{1\}, \{2\}, \{3\}\}\) — разбиение на три подмножества
		\end{itemize}
		Всего \(B(3) = 5\) разбиений.
	\end{examplebox}
	
	\subsubsection{Разбиения на заданное число подмножеств}
	
	\begin{definitionbox}{11.3: Разбиение мощности \(k\)}
		Пусть \(A\) — произвольное конечное множество, \(|A| = n\). Требуется подсчитать количество его разбиений мощности \(k\), т.е. таких наборов множеств \(\{X_1, X_2, \dots, X_k\}\), что:
		\begin{itemize}
			\item \(\forall i \in 1..k : X_i \neq \emptyset\) (все подмножества непусты);
			\item \(\forall i \neq j : X_i \cap X_j = \emptyset\) (они попарно дизъюнктны);
			\item \(\bigcup_{i=1}^k X_i = A\) (в объединении дают \(A\)).
		\end{itemize}
		Порядок подмножеств не учитывается (т.е. \(\{X_1, X_2\}\) и \(\{X_2, X_1\}\) — одно и то же разбиение).
	\end{definitionbox}
	
	\begin{examplebox}{11.4: Случай}
		Для \(k = 2\) нужно выбрать элементы, которые пойдут в первое множество разбиения, а все остальные пойдут во второе. Количество способов выбрать непустое собственное подмножество \(A\) (т.е. не совпадающее ни с \(\emptyset\), ни с \(A\)) равно \(2^n - 2\). 
		
		Так как мы не упорядочиваем элементы разбиения (пары \(\{X_1, X_2\}\) и \(\{X_2, X_1\}\) не различаем), полученное число нужно разделить на \(2\):
		\[ \frac{2^n - 2}{2} = 2^{n-1} - 1 \]
		
		Таким образом, количество разбиений множества из \(n\) элементов на 2 непустых подмножества равно \(2^{n-1} - 1\).
	\end{examplebox}
	
	\subsubsection{Рекуррентная формула для чисел Белла}
	
	\begin{theorembox}{11.5: Рекуррентная формула для чисел Белла}
		Для чисел Белла справедливо рекуррентное соотношение:
		\[ B(n+1) = \sum_{k=0}^{n} \binom{n}{k} B(k) \]
		где \(B(0) = 1\) (пустое множество имеет единственное разбиение — пустое).
	\end{theorembox}
	
	\myproof
	Рассмотрим множество \(A = \{a_1, a_2, \dots, a_n, a_{n+1}\}\) и произвольное его разбиение. Зафиксируем элемент \(a_{n+1}\). Пусть в разбиении он попадает в некоторый блок \(X_i\), причем \(|X_i| = j\), где \(1 \le j \le n+1\). Тогда \(|A \setminus X_i| = n+1-j\).
	
	\begin{itemize}
		\item Выбрать \(j-1\) элемент из оставшихся \(n\) элементов для блока, содержащего \(a_{n+1}\), можно \(\binom{n}{j-1}\) способами.
		\item Оставшееся множество \(A \setminus X_i\) мощности \(n+1-j\) можно разбить на произвольное число блоков \(B(n+1-j)\) способами.
	\end{itemize}
	
	Суммируя по всем возможным размерам блока, содержащего \(a_{n+1}\), получаем:
	\[ B(n+1) = \sum_{j=1}^{n+1} \binom{n}{j-1} B(n+1-j) \]
	
	Сделаем замену индекса: пусть \(k = n+1-j\), тогда \(j-1 = n-k\), и \(k\) пробегает значения от \(0\) до \(n\):
	\[ B(n+1) = \sum_{k=0}^{n} \binom{n}{n-k} B(k) = \sum_{k=0}^{n} \binom{n}{k} B(k) \]
	\qed
	
	\begin{examplebox}{11.6: Вычисление первых чисел Белла}
		\begin{align*}
			B(0) &= 1 \\
			B(1) &= \binom{0}{0} B(0) = 1 \cdot 1 = 1 \\
			B(2) &= \binom{1}{0} B(0) + \binom{1}{1} B(1) = 1 \cdot 1 + 1 \cdot 1 = 2 \\
			B(3) &= \binom{2}{0} B(0) + \binom{2}{1} B(1) + \binom{2}{2} B(2) = 1 \cdot 1 + 2 \cdot 1 + 1 \cdot 2 = 5 \\
			B(4) &= \binom{3}{0} B(0) + \binom{3}{1} B(1) + \binom{3}{2} B(2) + \binom{3}{3} B(3) \\
			&= 1 \cdot 1 + 3 \cdot 1 + 3 \cdot 2 + 1 \cdot 5 = 1 + 3 + 6 + 5 = 15
		\end{align*}
	\end{examplebox}
	
	\subsubsection{Числа Стирлинга второго рода}
	
	\begin{definitionbox}{11.7: Числа Стирлинга второго рода}
		\textbf{Числом Стирлинга второго рода} \(S(n, k)\) (обозначается также \(\genfrac{\{}{\}}{0pt}{}{n}{k}\)) называется количество разбиений множества мощности \(n\) ровно на \(k\) непустых подмножеств (блоков).\\
		
		Таким образом, числа Белла выражаются через числа Стирлинга:
		\[ B(n) = \sum_{k=0}^{n} S(n, k) \]
		где \(S(0,0) = 1\), а при \(k > n\) или \(k < 0\) полагают \(S(n, k) = 0\).
	\end{definitionbox}
	
	\myremark{ Обычно, говоря о числах Стирлинга без уточнения, имеют в виду числа второго рода, так как они встречаются чаще. \\
	
	Обозначение: $\{{}_n^k\}$\\
	
	Вспомним уже известное нам свойство: $C_n^k=C_{n-1}^k+C_{n-1}^{k-1},$ $C_n^0=1$\\
	
	Сформулируем аналогичное для чисел Стирлинга.}
	
	\subsubsection{Рекуррентная формула для чисел Стирлинга}
	
	\begin{theorembox}{11.8: Рекуррентная формула для \(S{(n,k)}\)}
		Для чисел Стирлинга второго рода справедливо рекуррентное соотношение:
		\[ S(n, k) = S(n-1, k-1) + k \cdot S(n-1, k) \]
		с начальными условиями:
		\begin{itemize}
			\item \(S(0,0) = 1\)
			\item \(S(n,0) = 0\) при \(n > 0\)
			\item \(S(0,k) = 0\) при \(k > 0\)
			\item \(S(n,k) = 0\) при \(k > n\)
		\end{itemize}
	\end{theorembox}
	
	\myproof
	Рассмотрим множество \(A = \{a_1, a_2, \dots, a_n\}\) и его разбиения на \(k\) блоков. Зафиксируем элемент \(a_n\). Возможны два случая:
	
	\begin{enumerate}
		\item \textbf{\(\{a_n\}\) образует отдельный блок.} Тогда оставшиеся \(n-1\) элементов нужно разбить на \(k-1\) блоков. Число таких разбиений равно \(S(n-1, k-1)\).
		
		\item \textbf{\(a_n\) находится в блоке вместе с другими элементами.} Тогда сначала строим разбиение множества \(A \setminus \{a_n\}\) на \(k\) блоков — это можно сделать \(S(n-1, k)\) способами. После этого элемент \(a_n\) нужно добавить в один из \(k\) уже существующих блоков — это можно сделать \(k\) способами. Всего получаем \(k \cdot S(n-1, k)\) способов.
	\end{enumerate}
	
	Эти два случая не пересекаются и покрывают все возможные разбиения, следовательно:
	\[ S(n, k) = S(n-1, k-1) + k \cdot S(n-1, k) \]
	\qed
	
	\begin{propertybox}{11.9: Свойства чисел Стирлинга}
		\begin{enumerate}
			\item \(S(n, 0) = 0\) для \(n > 0\)
			\item \(S(0, 0) = 1\)
			\item \(S(n, k) = 0\) при \(k > n\)
			\item \(S(n, 2) = 2^{n-1} - 1\) (разбиение на два подмножества)
			\item \(S(n, n-1) = \binom{n}{2}\) (разбиение состоит из одного блока мощности 2 и \(n-2\) блоков мощности 1; количество способов выбрать два элемента, которые попадут в один блок)
			\item \(S(n, n) = 1\) (единственное разбиение — все элементы по отдельности)
			\item \(S(n, 1) = 1\) (единственное разбиение — всё множество целиком)
		\end{enumerate}
	\end{propertybox}
	
	\begin{examplebox}{11.10: Пример вычисления \(S{(n,k)}\)}
		\begin{align*}
			S(3,1) &= 1 \\
			S(3,2) &= S(2,1) + 2 \cdot S(2,2) = 1 + 2 \cdot 1 = 3 \\
			S(3,3) &= S(2,2) + 3 \cdot S(2,3) = 1 + 3 \cdot 0 = 1 \\
			S(4,2) &= S(3,1) + 2 \cdot S(3,2) = 1 + 2 \cdot 3 = 7 = 2^{3} - 1 \\
			S(4,3) &= S(3,2) + 3 \cdot S(3,3) = 3 + 3 \cdot 1 = 6 = \binom{4}{2}
		\end{align*}
	\end{examplebox}
	
	\subsubsection{Явная формула для чисел Стирлинга второго рода}
	
	\begin{theorembox}{11.11: Явная формула для \(S{(n,k)}\)}
		Для любых \(n \ge 0\) и \(k \ge 1\) справедливо:
		\[ S(n, k) = \frac{1}{k!} \sum_{j=0}^{k} (-1)^j \binom{k}{j} (k-j)^n \]
	\end{theorembox}
	
	\myproof
	Пусть \(A\) — множество из \(n\) элементов. Число \(S(n,k)\) равно числу разбиений \(A\) на \(k\) непустых \textbf{неупорядоченных} подмножеств.\\
	
	Каждому разбиению \(A\) на \(k\) блоков можно сопоставить \textbf{сюръективное} отображение \(f: A \to \{1, 2, \dots, k\}\), где элемент \(a \in A\) отображается в номер блока, в который он входит. При этом разным нумерациям блоков соответствуют разные отображения. Таким образом, каждому разбиению соответствует ровно \(k!\) различных сюръекций (все возможные способы перенумеровать блоки).\\
	
	Следовательно,
	\[ S(n, k) = \frac{1}{k!} \cdot L \]
	где \(L\) — число сюръективных отображений из \(A\) в \(\{1, 2, \dots, k\}\).\\
	
	Найдём \(L\) методом включений-исключений.
	
	\begin{enumerate}
		\item Общее число всех отображений \(f: A \to \{1, \dots, k\}\) равно \(k^n\).
		\item Для каждого \(i \in \{1, \dots, k\}\) обозначим через \(P_i\) множество отображений, не принимающих значение \(i\) (т.е. в образе которых нет числа \(i\)). Тогда \(|P_i| = (k-1)^n\).
		\item Для набора индексов \(\{i_1, \dots, i_j\}\) пересечение \(P_{i_1} \cap \dots \cap P_{i_j}\) состоит из отображений, не принимающих значения \(i_1, \dots, i_j\), т.е. отображений в множество из \(k-j\) элементов. Таких отображений \((k-j)^n\).
	\end{enumerate}
	
	По формуле включений-исключений, число \textbf{несюръективных} отображений равно:
	\[ \sum_{j=1}^{k} (-1)^{j+1} \sum_{1 \le i_1 < \dots < i_j \le k} |P_{i_1} \cap \dots \cap P_{i_j}| = \sum_{j=1}^{k} (-1)^{j+1} \binom{k}{j} (k-j)^n \]
	
	Тогда число сюръективных отображений:
	\begin{align*}
		L &= k^n - \sum_{j=1}^{k} (-1)^{j+1} \binom{k}{j} (k-j)^n \\
		&= k^n + \sum_{j=1}^{k} (-1)^{j} \binom{k}{j} (k-j)^n \\
		&= \sum_{j=0}^{k} (-1)^{j} \binom{k}{j} (k-j)^n
	\end{align*}
	
	Делим на \(k!\) и получаем искомую формулу:
	\[ S(n, k) = \frac{1}{k!} \sum_{j=0}^{k} (-1)^j \binom{k}{j} (k-j)^n \]
	\qed
	
	\begin{corollarybox}{11.12: Проверка для (n=0)}
		При \(n = 0\) формула даёт:
		\[ S(0, k) = \frac{1}{k!} \sum_{j=0}^{k} (-1)^j \binom{k}{j} (k-j)^0 = \frac{1}{k!} \sum_{j=0}^{k} (-1)^j \binom{k}{j} = \frac{1}{k!} (1 - 1)^k = 0 \]
		для \(k > 0\), что согласуется с определением (пустое множество нельзя разбить на \(k>0\) непустых частей). При \(k=0\) сумма даёт \(1/0! \cdot (-1)^0 \binom{0}{0} \cdot 0^0\), что интерпретируется как \(1\) (соглашение \(0^0 = 1\) в комбинаторике).
	\end{corollarybox}
	
	\begin{examplebox}{11.13: Пример применения явной формулы}
		Вычислим \(S(4,2)\) по явной формуле:
		\begin{align*}
			S(4,2) &= \frac{1}{2!} \sum_{j=0}^{2} (-1)^j \binom{2}{j} (2-j)^4 \\
			&= \frac{1}{2} \left[ (-1)^0 \binom{2}{0} \cdot 2^4 + (-1)^1 \binom{2}{1} \cdot 1^4 + (-1)^2 \binom{2}{2} \cdot 0^4 \right] \\
			&= \frac{1}{2} \left[ 1 \cdot 1 \cdot 16 - 2 \cdot 1 \cdot 1 + 1 \cdot 1 \cdot 0 \right] = \frac{1}{2} (16 - 2) = 7
		\end{align*}
		Что совпадает с \(2^{3} - 1 = 7\).
	\end{examplebox}
	
	\begin{center}
		\begin{tabular}{|c|cccccc|}
			\hline
			\(n \setminus k\) & 0 & 1 & 2 & 3 & 4 & 5 \\ \hline
			0 & 1 & 0 & 0 & 0 & 0 & 0 \\
			1 & 0 & 1 & 0 & 0 & 0 & 0 \\
			2 & 0 & 1 & 1 & 0 & 0 & 0 \\
			3 & 0 & 1 & 3 & 1 & 0 & 0 \\
			4 & 0 & 1 & 7 & 6 & 1 & 0 \\
			5 & 0 & 1 & 15 & 25 & 10 & 1 \\ \hline
		\end{tabular}
		\captionof{table}{Треугольник чисел Стирлинга второго рода \(S(n,k)\)}
	\end{center}
	
	\myremark{11.14: Связь с числами Белла}
		Числа Белла получаются суммированием чисел Стирлинга по строкам:\\ \\
		B(0)=1;\\
		B(1)=1; \\
		B(2)=1+1=2;\\ 
		B(3)=1+3+1=5;\\
		B(4)=1+7+6+1=15;\\
		B(5)=1+15+25+10+1=52
	
	
	
	
	
	
	
	
	
	
	% ==================== РАЗДЕЛ 12: ПРОИЗВОДЯЩИЕ ФУНКЦИИ ====================
	\newpage
	\section{Производящие функции}
	
	\subsection{Производящие функции (12 вопрос)}
	
	\subsubsection{Определение производящей функции}
	
	\begin{definitionbox}{12.1: Производящая функция}
		Для последовательности чисел \(\{a_n\}\) рассмотрим формальную сумму \(\sum a_n t^n\), где \(t \in \mathbb{C}\).
		
		\begin{itemize}
			\item Если последовательность \(\{a_n\}\) \textbf{конечна}, то эта сумма всегда определяет функцию (многочлен)
			\[ F(t) = \sum a_n t^n \]
			которая называется \textbf{производящей функцией} для \(\{a_n\}\).
			
			\item Если последовательность \(\{a_n\}\) \textbf{бесконечна}, то получается степенной ряд
			\[ \sum_{n=0}^\infty a_n t^n \]
			который может сходиться в некоторой области \(D \subseteq \mathbb{C}\) к некоторой функции \(F(t)\). Эта функция также называется \textbf{производящей функцией} для \(\{a_n\}\).
				
			\item Если задана производящая функция \(F(t)\) для некоторой последовательности \(\{a_n\}\), то говорят, что последовательность \(\{a_n\}\) \textbf{полностью определена}.
		\end{itemize}	
	\end{definitionbox}
	
	\begin{examplebox}{12.2: Пример конечной последовательности}
		Дана конечная последовательность \(a_k = C_n^k\), где \(0 \leq k \leq n\). Найдём её производящую функцию \(F(t)\).
		
		\textbf{Решение:}
		\[ F_n(t) = \sum_{k=0}^n C_n^k t^k = (1 + t)^n \]
		— это формула бинома Ньютона.
	\end{examplebox}
	
	\subsubsection{Элементарные преобразования производящих функций}
	
	\begin{definitionbox}{12.3: Сумма производящих функций}
		Суммой двух производящих функций
		\[ A(s) = a_0 + a_1 s + a_2 s^2 + \dots \]
		\[ B(s) = b_0 + b_1 s + b_2 s^2 + \dots \]
		называется производящая функция
		\[ A(s) + B(s) = (a_0 + b_0) + (a_1 + b_1)s + (a_2 + b_2)s^2 + \dots \]
	\end{definitionbox}
	
	\begin{definitionbox}{12.4: Произведение производящих функций}
		Произведением двух производящих функций \(A\) и \(B\) называется производящая функция
		\[ A(s)B(s) = a_0b_0 + (a_0b_1 + a_1b_0)s + (a_0b_2 + a_1b_1 + a_2b_0)s^2 + \dots \]
		где коэффициент при \(s^n\) равен \(\sum_{k=0}^n a_k b_{n-k}\).
	\end{definitionbox}
	
	\myremark{Свойства:} Операции сложения и умножения производящих функций коммутативны и ассоциативны (легко доказывается по определению).
	
	\begin{definitionbox}{12.5: Подстановка производящих функций}
		Пусть даны две производящие функции
		\[ A(s) = a_0 + a_1 s + a_2 s^2 + \dots \]
		\[ B(t) = b_1 t + b_2 t^2 + \dots \]
		причём \(B(0) = 0\). Тогда \textbf{подстановкой} \(B\) в \(A\) называется производящая функция
		\[ A(B(t)) = a_0 + a_1 B(t) + a_2 (B(t))^2 + \dots = \]
		\[ = a_0 + a_1 b_1 t + (a_1 b_2 + a_2 b_1^2) t^2 + \dots \]
	\end{definitionbox}
	
	\begin{examplebox}{12.6: Пример подстановки}
		Пусть \(B(t) = -t\). Это соответствует последовательности \(0, -1, 0, 0, \dots\).
		
		Тогда
		\[ A(B(t)) = A(-t) = a_0 - a_1 t + a_2 t^2 - a_3 t^3 + \dots \]
	\end{examplebox}
	
	\subsubsection{Элементарные производящие функции}
	
	\begin{remarkbox}{12.7: Стандартные разложения}
		Всякий раз записывать производящие функции в виде ряда неудобно. Поэтому для некоторых часто встречающихся функций используется сокращенная запись (ряды Тейлора):
		
		\begin{enumerate}
			\item \textbf{Биномиальное разложение:}
			\[ (1 + s)^\alpha = 1 + \frac{\alpha}{1!}s + \frac{\alpha(\alpha-1)}{2!}s^2 + \frac{\alpha(\alpha-1)(\alpha-2)}{3!}s^3 + \dots, \quad \alpha \in \mathbb{C} \]
			
			\item \textbf{Экспонента:}
			\[ e^s = \exp(s) = 1 + \frac{1}{1!}s + \frac{1}{2!}s^2 + \frac{1}{3!}s^3 + \dots \]
			
			\item \textbf{Логарифм:}
			\[ \ln\left(\frac{1}{1-s}\right) = s + \frac{1}{2}s^2 + \frac{1}{3}s^3 + \dots \]
			
			\item \textbf{Синус:}
			\[ \sin s = s - \frac{1}{3!}s^3 + \frac{1}{5!}s^5 - \dots \]
			
			\item \textbf{Косинус:}
			\[ \cos s = 1 - \frac{1}{2!}s^2 + \frac{1}{4!}s^4 - \dots \]
		\end{enumerate}
	\end{remarkbox}
	
	\subsubsection{Дифференцирование и интегрирование}
	
	\begin{definitionbox}{12.8: Производная производящей функции}
		Пусть \(A(s) = a_0 + a_1 s + a_2 s^2 + \dots\) — производящая функция. Её \textbf{производной} называется функция
		\[ A'(s) = a_1 + 2a_2 s + 3a_3 s^2 + \dots + n a_n s^{n-1} + \dots \]
	\end{definitionbox}
	
	\begin{definitionbox}{12.9: Интеграл производящей функции}
		\textbf{Интегралом} производящей функции \(A(s)\) называется функция
		\[ \int A(s)   = a_0 s + a_1 \frac{s^2}{2} + a_2 \frac{s^3}{3} + \dots + a_n \frac{s^{n+1}}{n+1} + \dots \]
		
		Эта формула соответствует значению интеграла с переменным верхним пределом:
		\[ \int A(s) = \int_0^s A(\xi) \, d\xi \]
	\end{definitionbox}
	
	\begin{propertybox}{12.10: Свойства дифференцирования и интегрирования}
		\begin{itemize}
			\item \(\left( \int A(s) \right)' = A(s)\)
			\item \(\int A'(s) = A(s) - a_0\)
		\end{itemize}
	\end{propertybox}
	
	\begin{examplebox}{12.11: Вычисление производящей функции}
		Вычислить производящую функцию
		\[ f(s) = \frac{1}{1\cdot2} + \frac{1}{2\cdot3} s + \frac{1}{3\cdot4} s^2 + \dots + \frac{1}{(n+1)(n+2)} s^n + \dots \]
		
		\textbf{Решение:} Умножим функцию \(f\) на \(s^2\) и продифференцируем. Получаем элементарную производящую функцию:
		\[ (s^2 f(s))' = s + \frac{1}{2} s^2 + \frac{1}{3} s^3 + \dots = \ln\left(\frac{1}{1-s}\right) \]
		
		Откуда
		\[ f(s) = s^{-2} \int \ln\left(\frac{1}{1-s}\right) ds = s^{-2}\left((s-1)\ln\left(\frac{1}{1-s}\right) + s\right) \]
	\end{examplebox}
	
	\subsubsection{Геометрическая прогрессия}
	
	\begin{definitionbox}{12.12: Производящая функция постоянной последовательности}
		Простейшая последовательность — это постоянная последовательность \(1, 1, 1, \dots\). Её производящая функция имеет вид
		\[ G(s) = 1 + s + s^2 + s^3 + \dots \]
		
		Умножим обе части равенства на \(s\):
		\[ sG(s) = s + s^2 + s^3 + s^4 + \dots = G(s) - 1 \]
		
		Откуда
		\[ G(s) = \frac{1}{1-s} \]
	\end{definitionbox}
	
	\begin{corollarybox}{12.13: Обобщение на геометрическую прогрессию}
		Для произвольной геометрической прогрессии вида \(a, ar, ar^2, ar^3, \dots\):
		\[ G_{a,r}(s) = a + ar s + ar^2 s^2 + ar^3 s^3 + \dots = a(1 + (rs) + (rs)^2 + (rs)^3 + \dots) \]
		
		Аналогично получаем:
		\[ r \cdot G_{a,r}(s) = G_{a,r}(s) - a \quad\Rightarrow\quad G_{a,r}(s) = \frac{a}{1-rs} \]
	\end{corollarybox}
	
	\subsubsection{Сочетания с повторениями}
	
	\begin{examplebox}{12.14: Производящая функция для \(C^k_{n+k-1}\)}
		Рассмотрим последовательность \(a_k = C^k_{n+k-1}\). Для частного случая \(n = 3\):
		\[ a_k = C^k_{k+2} = 1, 3, 6, 10, 15, 21, \dots \]
		
		Производящая функция:
		\[ f(t) = 1 + 3t + 6t^2 + 10t^3 + 15t^4 + 21t^5 + \dots \]
		
		Для получения производящей функции необходимо перемножить \(n = 3\) бесконечных сумм:
		\[ 1 + t + t^2 + t^3 + \dots = \frac{1}{1-t} \]
		
		Следовательно,
		\[ f(t) = \frac{1}{(1-t)^3} \]
		
		В общем случае:
		\[ \sum_{k=0}^\infty C^k_{n+k-1} t^k = \frac{1}{(1-t)^n} \]
	\end{examplebox}
	
	\subsubsection{Задачи с лекции}
	
	\begin{examplebox}{12.15: Разбиения числа}
		Доказать, что \(\displaystyle \prod_{n=1}^\infty (1 - x^n)\) является производящей функцией для разности количества разбиений числа \(n\) в чётное и нечётное число различных слагаемых.\\
		
		\textbf{Пояснение:} Каждая скобка \((1 - x^k)\) даёт право выбора: если мы не берём слагаемое \(k\), то выбираем \(1\) в скобке, а если берём — то выбираем \(x^k\). При этом такие \(x^k\) для каждой скобки различны, и суммарно должна получаться заданная степень \(n\).\\
		
		Так как перед \(x^k\) в скобке стоит минус, каждый раз при его выборе знак изменяется. При чётном количестве "изъятий" каждое разбиение добавляет \(1\) к коэффициенту, а при нечётном — вычитает \(1\). Таким образом, итоговый коэффициент равен разности числа чётных и нечётных разбиений.
	\end{examplebox}
	
	\begin{examplebox}{12.16: Задача о гирях}
		Какие грузы можно взвесить гирями в \(2^0, 2^1, 2^2, \dots, 2^m\) грамм, и сколькими способами?\\
		
		\textbf{Решение:} Производящая функция для этой задачи:
		\[ \alpha(x) = (1 + x)(1 + x^2)(1 + x^4)(1 + x^8)\cdots(1 + x^{2^m})\cdots \]
		
		После раскрытия скобок и приведения подобных членов получаем:
		\[ \alpha(x) = 1 + A_1 x + A_2 x^2 + A_3 x^3 + \dots \]
		
		Коэффициенты \(A_k\) равны числу различных представлений числа \(k\) в виде суммы некоторых из чисел \(1, 2, 2^2, 2^3, \dots\).\\
		
		Домножим обе части на \((1 - x)\) и воспользуемся тождествами:
		\begin{align*}
			(1 - x)(1 + x) &= (1 - x^2) \\
			(1 - x^2)(1 + x^2) &= (1 - x^4) \\
			(1 - x^4)(1 + x^4) &= (1 - x^8) \\
			&\vdots
		\end{align*}
		
		Получим:
		\[ (1 - x)\alpha(x) = 1 \quad\Rightarrow\quad \alpha(x) = \frac{1}{1 - x} = 1 + x + x^2 + x^3 + \dots \]
		
		Таким образом, \(A_k = 1\) для всех \(k\), т.е. всякий груз в целое число грамм можно взвесить указанными гирями, и притом единственным способом.
	\end{examplebox}
	
	\subsubsection{Числа Шрёдера}
	
	\begin{definitionbox}{12.17: Числа Шрёдера}
		\textbf{Числа Шрёдера} \(S_n\) в комбинаторике описывают количество путей из левого нижнего угла квадратной решётки \(n \times n\) в противоположный по диагонали угол, используя ходы:
		\begin{itemize}
			\item вверх \((0, 1)\)
			\item вправо \((1, 0)\)
			\item вверх-вправо \((1, 1)\)
		\end{itemize}
		при этом пути не поднимаются выше диагонали квадратной решётки.
	\end{definitionbox}
	
	
	
	\subsubsection{Рекуррентное соотношение для чисел Шрёдера}
	
	\begin{theorembox}{12.19: Рекуррентная формула для чисел Шрёдера}
		Для чисел Шрёдера справедливо рекуррентное соотношение:
		\[ S_n = S_{n-1} + \sum_{k=0}^{n-1} S_k S_{n-1-k}, \quad n \ge 1 \]
		с начальным условием \(S_0 = 1\).
	\end{theorembox}
	
	\myproof
	Рассмотрим пути из \((0,0)\) в \((n,n)\), не пересекающие диагональ. Первый шаг не может быть вверх (иначе пересечём диагональ). Возможны два случая:
	
	\begin{enumerate}
		\item \textbf{Первый шаг по диагонали} \((1,1)\). Тогда мы попадаем в точку \((1,1)\), и остаётся квадрат \((n-1) \times (n-1)\). Число таких путей равно \(S_{n-1}\).
		
		\item \textbf{Первый шаг вправо} \((1,0)\). Тогда на некотором шаге мы впервые выйдем на диагональ. Пусть это произойдёт в точке \((k,k)\). От начала до этой точки мы шли, не поднимаясь выше диагонали, и первый шаг был вправо — число таких путей равно \(S_{n-1-k}\). От точки \((k,k)\) до \((n,n)\) — \(S_k\) путей. Суммируем по всем \(k\) от \(0\) до \(n-1\).
	\end{enumerate}
	
	Складывая оба случая, получаем искомую формулу. \qed
	
	
		
	\begin{center}
		\begin{tikzpicture}[scale=0.6]
			
			% Параметр n
			\def\n{6}
			
			% Сетка
			\draw[step=1cm, gray!30, very thin] (0,0) grid (\n,\n);
			
			% Диагональ
			\draw[thick, red] (0,0) -- (\n,\n);
			
			% Оси
			\draw[->, thick] (0,0) -- (\n+0.5,0);
			\draw[->, thick] (0,0) -- (0,\n+0.5);
			
			\node at (\n/2,-0.7) {вправо};
			\node[rotate=90] at (-0.7,\n/2) {вверх};
			
			\node at (\n/2,\n+0.7) {Диагональ};
			
		\end{tikzpicture}
		\captionof{figure}{Квадратная решётка и ограничение диагональю}
	\end{center}
	
	
	\begin{center}
		\begin{tikzpicture}[scale=0.6]
			
			\def\n{6}
			
			% Сетка
			\draw[step=1cm, gray!30, very thin] (0,0) grid (\n,\n);
			
			% Диагональ
			\draw[thick, red] (0,0) -- (\n,\n);
			
			% Путь
			\draw[very thick, blue]
			(0,0)
			-- (1,0)
			-- (2,1)
			-- (3,1)
			-- (4,2)
			-- (5,3)
			-- (6,6);
			
			% Точки начала и конца
			\fill (0,0) circle (3pt);
			\fill (\n,\n) circle (3pt);
			
			\node at (0,-0.4) {Старт};
			\node at (\n,\n+0.4) {Финиш};
			
		\end{tikzpicture}
		\captionof{figure}{Пример допустимого пути}
	\end{center}
	
	
	\subsubsection{Производящая функция для чисел Шрёдера}
	
	\begin{theorembox}{12.20: Производящая функция чисел Шрёдера}
		Производящая функция для чисел Шрёдера \(G(x) = \sum_{n=0}^\infty S_n x^n\) имеет вид:
		\[ G(x) = \frac{1 - \sqrt{1 - 6x + x^2}}{2x} \]
	\end{theorembox}
	
	\myproof
	Запишем рекуррентное соотношение для \(n \ge 1\):
	\[ S_n = S_{n-1} + \sum_{k=0}^{n-1} S_k S_{n-1-k} \]
	
	Умножим на \(x^n\) и просуммируем по \(n \ge 1\):
	\[ \sum_{n=1}^\infty S_n x^n = \sum_{n=1}^\infty S_{n-1} x^n + \sum_{n=1}^\infty \sum_{k=0}^{n-1} S_k S_{n-1-k} x^n \]
	
	Заметим, что:
	\[ \sum_{n=1}^\infty S_{n-1} x^n = x \sum_{n=1}^\infty S_{n-1} x^{n-1} = x G(x) \]
	и
	\[ \sum_{n=1}^\infty \sum_{k=0}^{n-1} S_k S_{n-1-k} x^n = x \sum_{n=1}^\infty \sum_{k=0}^{n-1} (S_k x^k)(S_{n-1-k} x^{n-1-k}) = x (G(x))^2 \]
	
	Таким образом,
	\[ G(x) - S_0 = x G(x) + x (G(x))^2 \]
	Учитывая \(S_0 = 1\), получаем:
	\[ G(x) - 1 = x G(x) + x G^2(x) \]
	\[ x G^2(x) + (x - 1) G(x) + 1 = 0 \]
	
	Решая квадратное уравнение относительно \(G(x)\) и выбирая ветвь с \(G(0) = 1\), получаем:
	\[ G(x) = \frac{1 - \sqrt{1 - 6x + x^2}}{2x} \]
	\qed
	
	\begin{center}
		\begin{tabular}{|c|c|c|c|c|c|c|c|}
			\hline
			\(n\) & 0 & 1 & 2 & 3 & 4 & 5 & 6 \\ \hline
			\(S_n\) & 1 & 2 & 6 & 22 & 90 & 394 & 1806 \\ \hline
		\end{tabular}
		\captionof{table}{Первые числа Шрёдера}
	\end{center}
	
	
	\begin{remarkbox}{12.18: Эквивалентные формулировки}
		Числа Шрёдера также равны:
		\begin{itemize}
			\item Количеству способов разрезания прямоугольника на \(n + 1\) меньших прямоугольников с помощью \(n\) разрезов, проводимых через заданные точки внутри прямоугольника (никакие две из которых не лежат на одной прямой, параллельной сторонам), при этом каждый разрез проходит через одну из точек и делит только один прямоугольник на два.
			
			\item Количеству путей из точки \((0, 0)\) в \((2n, 0)\), использующих шаги вправо-вверх \((1, 1)\), вправо-вниз \((1, -1)\) или двойные шаги вправо \((2, 0)\), которые не опускаются ниже оси \(x\).
		\end{itemize}
	\end{remarkbox}
	
	
	\begin{center}
		\begin{tikzpicture}[scale=0.8]
			
			% Прямоугольник
			\draw[thick] (0,0) rectangle (8,5);
			
			% Заданные точки
			\fill (2,1) circle (3pt);
			\fill (5,3) circle (3pt);
			\fill (6,1.5) circle (3pt);
			\fill (3,4) circle (3pt);
			
			% Вертикальные разрезы
			\draw[very thick, blue] (2,0) -- (2,5);
			\draw[very thick, blue] (5,0) -- (5,5);
			
			% Горизонтальный разрез
			\draw[very thick, red] (0,3) -- (8,3);
			
			\node at (4,-0.7) {Разбиение на $n+1$ прямоугольников};
			
		\end{tikzpicture}
		\captionof{figure}{Разрезание прямоугольника (эквивалентная формулировка)}
	\end{center}
	
	
	\begin{center}
		\begin{tikzpicture}[scale=0.8]
			
			% Оси
			\draw[->, thick] (0,0) -- (10,0);
			\draw[thick] (0,0) -- (0,5);
			
			% Ось x
			\node at (5,-0.6) {$x$};
			
			% Запрещённая область
			\fill[red!10] (0,0) -- (10,0) -- (10,5) -- (0,5) -- cycle;
			
			% Разрешённый путь
			\draw[very thick, blue]
			(0,0)
			-- (1,1)
			-- (3,1)
			-- (4,0)
			-- (6,0)
			-- (8,2)
			-- (10,0);
			
			% Подписи
			\node at (10,0) [below right] {$(2n,0)$};
			\node at (0,0) [below left] {$(0,0)$};
			
			% Легенда шагов
			\node at (6,4.5) {шаги: $(1,1)$, $(1,-1)$, $(2,0)$};
			
		\end{tikzpicture}
		\captionof{figure}{Путь, не опускающийся ниже оси $x$}
	\end{center}
	
	
	
	
	
	
	
	
	
	
	
	
	
	
	
	% ==================== РАЗДЕЛ 13: ТЕОРИЯ ВЕРОЯТНОСТЕЙ ====================
	\newpage
	\section{Основы теории вероятностей}
	
	\subsection{Парадокс Монти Холла. Эксперимент, исход, событие, вероятность. Условная вероятность (13 вопрос)}
	
	\subsubsection{Вероятностное пространство}
	
	\begin{definitionbox}{13.1: Вероятностное пространство}
		Пусть \(S\) — конечное множество, \(|S| = n\). Пусть задана функция \(f: S \to [0,1]\), такая что:
		\begin{itemize}
			\item \(\forall \omega \in S\) существует единственное значение \(f(\omega) \in [0,1]\);
			\item \(\sum_{\omega \in S} f(\omega) = 1\).
		\end{itemize}
		
		Для любого подмножества \(A \subseteq S\) определим:
		\[ \Pr(A) := \sum_{\omega \in A} f(\omega) \]
		
		В частности:
		\begin{itemize}
			\item \(\Pr(\emptyset) = 0\);
			\item \(\Pr(S) = 1\);
			\item \(\Pr(\{\omega\}) = f(\omega)\).
		\end{itemize}
		
		Таким образом, необходимость в исходной функции \(f\) пропадает — нам достаточно иметь \(\Pr\).
	\end{definitionbox}
	
	\begin{definitionbox}{13.2: Терминология}
		\begin{itemize}
			\item \((S, \Pr)\) называется \textbf{вероятностным пространством};
			\item \(S\) — \textbf{пространством элементарных событий};
			\item \(\omega \in S\) — \textbf{элементарным событием} (исходом);
			\item \(A \subseteq S\) — \textbf{событием};
			\item \(\Pr(A)\) — \textbf{вероятностью} события \(A\).
		\end{itemize}
	\end{definitionbox}
	
	\begin{definitionbox}{13.3: Несовместные события}
		События \(A, B \subseteq S\) называются \textbf{несовместными}, если
		\[ \Pr(A \cap B) = 0 \]
	\end{definitionbox}
	
	\begin{notebook}{13.4: Обозначения}
		Для краткости используется запись:
		\[ \Pr\{P(x)\} := \Pr(\{\omega \in S : P(\omega)\}) \]
		— вероятность множества таких элементарных исходов, для которых выполняется условие \(P\).\\
		
		Условие может быть записано в произвольном формате, например:
		\[ \Pr(\text{Сборная России по футболу выиграет чемпионат мира}) \]
	\end{notebook}
	
	\subsubsection{Свойства вероятности}
	
	\begin{propertybox}{13.5: Основные свойства вероятности}
		
		\textbf{1. Формула включений и исключений для двух событий:}
		\[ \forall A, B \subseteq S: \Pr(A \cup B) = \Pr(A) + \Pr(B) - \Pr(A \cap B) \]
		
		\myproof
		\[ \Pr(A \cup B) = \sum_{\omega \in A \cup B} f(\omega) = \sum_{\omega \in A} f(\omega) + \sum_{\omega \in B} f(\omega) - \sum_{\omega \in A \cap B} f(\omega) \]
		\qed
		
		\textbf{2. Вероятность дополнения:}
		\[ \forall A \subseteq S: \Pr(A) + \Pr(\overline{A}) = 1, \quad \text{где } \overline{A} = S \setminus A \]
		
		\myproof
		Следует из определения вероятности и того факта, что \(A\) и \(\overline{A}\) образуют разбиение \(S\):
		\[ \Pr(A) + \Pr(\overline{A}) = \sum_{\omega \in A} f(\omega) + \sum_{\omega \in \overline{A}} f(\omega) = \sum_{\omega \in S} f(\omega) = 1 \]
		\qed
		
		\textbf{3. Полуаддитивность:}
		\[ \forall A, B \subseteq S: \Pr(A \cup B) \le \Pr(A) + \Pr(B) \]
		
		\myproof
		Очевидно из свойства 1 и неотрицательности вероятности \(\Pr(A \cap B) \ge 0\). \qed \\
		
		\textbf{4. Разложение вероятности:}
		\[ \forall A, B \subseteq S: \Pr(A) = \Pr(A \setminus B) + \Pr(A \cap B) \]
		
		\myproof
		Множества \(A \setminus B\) и \(A \cap B\) не пересекаются и в объединении дают \(A\). \qed
	\end{propertybox}
	
	
	\subsubsection{Парадокс Монти Холла}
	
	\begin{examplebox}{13.8: Парадокс Монти Холла}
		\textbf{Неформальная формулировка:} На телешоу ведущий предлагает игроку выбрать одну из трёх дверей. Известно, что за одной из дверей находится автомобиль, а за двумя другими — по козе. После того как игрок сделал выбор, ведущий открывает одну из двух оставшихся дверей (причём обязательно ту, за которой коза — открыть дверь с автомобилем он не может) и предлагает игроку изменить свой выбор. \\
		
		\textbf{Вопрос:} Стоит ли игроку менять выбор? \\
		
		\textbf{Предположения:}
		\begin{itemize}
			\item Организаторы выбирают дверь с автомобилем наугад;
			\item Игрок выбирает дверь наугад;
			\item Если у ведущего есть выбор (т.е. игрок изначально выбрал дверь с автомобилем), он выбирает одну из двух оставшихся дверей наугад;
			\item Никто не нарушает правила игры.
		\end{itemize}
	\end{examplebox}
	
	\subsubsection{Дерево вариантов}
	
	Для решения задачи построим дерево всех возможных исходов эксперимента.
	
	\begin{center}
		\begin{tikzpicture}[scale=0.9, transform shape]
			% Уровень 1: выбор организаторов
			\node[circle, draw, blue] (root) at (0,0) {Старт};
			\node[circle, draw, blue] (A1) at (-4,-1.5) {\(A_1\)};
			\node[circle, draw, blue] (A2) at (0,-1.5) {\(A_2\)};
			\node[circle, draw, blue] (A3) at (4,-1.5) {\(A_3\)};
			
			\draw[->] (root) -- (A1) node[midway, left] {\(\frac{1}{3}\)};
			\draw[->] (root) -- (A2) node[midway, above] {\(\frac{1}{3}\)};
			\draw[->] (root) -- (A3) node[midway, right] {\(\frac{1}{3}\)};
			
			% Уровень 2: выбор игрока (исправлено позиционирование)
			\node[circle, draw, blue] (A1P1) at (-6,-3) {\(P_1\)};
			\node[circle, draw, blue] (A1P2) at (-4,-3) {\(P_2\)};
			\node[circle, draw, blue] (A1P3) at (-2,-3) {\(P_3\)};
			
			\node[circle, draw, blue] (A2P1) at (-1,-3) {\(P_1\)};
			\node[circle, draw, blue] (A2P2) at (0,-3) {\(P_2\)};
			\node[circle, draw, blue] (A2P3) at (1,-3) {\(P_3\)};
			
			\node[circle, draw, blue] (A3P1) at (2,-3) {\(P_1\)};
			\node[circle, draw, blue] (A3P2) at (4,-3) {\(P_2\)};
			\node[circle, draw, blue] (A3P3) at (6,-3) {\(P_3\)};
			
			\draw[->] (A1) -- (A1P1) node[midway, left] {\(\frac{1}{3}\)};
			\draw[->] (A1) -- (A1P2) node[midway, above] {\(\frac{1}{3}\)};
			\draw[->] (A1) -- (A1P3) node[midway, right] {\(\frac{1}{3}\)};
			
			\draw[->] (A2) -- (A2P1) node[midway, left] {\(\frac{1}{3}\)};
			\draw[->] (A2) -- (A2P2) node[midway, above] {\(\frac{1}{3}\)};
			\draw[->] (A2) -- (A2P3) node[midway, right] {\(\frac{1}{3}\)};
			
			\draw[->] (A3) -- (A3P1) node[midway, left] {\(\frac{1}{3}\)};
			\draw[->] (A3) -- (A3P2) node[midway, above] {\(\frac{1}{3}\)};
			\draw[->] (A3) -- (A3P3) node[midway, right] {\(\frac{1}{3}\)};
			
			% Пояснения
			\node at (-8, -1) {\(A_i\) — автомобиль за дверью \(i\)};
			\node at (-8, -2.5) {\(P_j\) — игрок выбрал дверь \(j\)};
		\end{tikzpicture}
	\end{center}
	
	Далее ведущий выбирает, какую дверь открыть. Заметим, что выбор у ведущего есть только если игрок исходно выбрал дверь, за которой находится приз. В таком случае мы считаем, что ведущий делает выбор случайным образом (вероятность каждого из двух выборов — \(\frac{1}{2}\)).
	
	\begin{center}
		\begin{tikzpicture}[scale=0.8, transform shape]
			% Дерево для случая A1
			\node[circle, draw, blue] (A1) at (0,0) {\(A_1\)};
			
			\node[circle, draw, blue] (A1P1) at (-4,-1.5) {\(P_1\)};
			\node[circle, draw, blue] (A1P2) at (0,-1.5) {\(P_2\)};
			\node[circle, draw, blue] (A1P3) at (4,-1.5) {\(P_3\)};
			
			\draw[->] (A1) -- (A1P1) node[midway, left] {\(\frac{1}{3}\)};
			\draw[->] (A1) -- (A1P2) node[midway, above] {\(\frac{1}{3}\)};
			\draw[->] (A1) -- (A1P3) node[midway, right] {\(\frac{1}{3}\)};
			
			% Выбор ведущего для A1P1 (игрок угадал)
			\node[circle, draw, blue] (A1P1H2) at (-6,-3) {\(H_2\)};
			\node[circle, draw, blue] (A1P1H3) at (-2,-3) {\(H_3\)};
			\draw[->] (A1P1) -- (A1P1H2) node[midway, left] {\(\frac{1}{2}\)};
			\draw[->] (A1P1) -- (A1P1H3) node[midway, right] {\(\frac{1}{2}\)};
			\node[draw, fill=red!20] at (-6,-4) {Проигрыш};
			\node[draw, fill=red!20] at (-2,-4) {Проигрыш};
			
			% Выбор ведущего для A1P2 (игрок не угадал)
			\node[circle, draw, blue] (A1P2H3) at (0,-3) {\(H_3\)};
			\draw[->] (A1P2) -- (A1P2H3) node[midway, right] {\(1\)};
			\node[draw, fill=green!20] at (0,-4) {Выигрыш};
			
			% Выбор ведущего для A1P3 (игрок не угадал)
			\node[circle, draw, blue] (A1P3H2) at (4,-3) {\(H_2\)};
			\draw[->] (A1P3) -- (A1P3H2) node[midway, right] {\(1\)};
			\node[draw, fill=green!20] at (4,-4) {Выигрыш};
			
			% Пояснения
			\node at (0,-5) {\(H_k\) — ведущий открыл дверь \(k\)};
			\node at (0,-5.5) {При смене двери};
		\end{tikzpicture}
	\end{center}
	
	Аналогичные деревья строятся для случаев \(A_2\) и \(A_3\).
	
	Вероятность каждого исхода можно посчитать как произведение вероятностей на пути из корня дерева в лист, соответствующий данному исходу.
	
	Например:
	\[ \Pr(\{(A_1, P_1, H_2)\}) = \frac{1}{3} \cdot \frac{1}{3} \cdot \frac{1}{2} = \frac{1}{18} \]
	
	\begin{theorembox}{13.9: Решение парадокса Монти Холла}
		Вероятность выигрыша при смене выбора равна \(\frac{2}{3}\), а вероятность выигрыша при сохранении первоначального выбора равна \(\frac{1}{3}\). Следовательно, игроку выгодно менять свой выбор.
	\end{theorembox}
	
	\myproof
	Подсчитаем вероятность выигрыша при стратегии «всегда менять выбор».
	
	Рассмотрим все случаи, когда игрок выигрывает, поменяв выбор:
	\begin{itemize}
		\item Игрок изначально выбрал дверь с козой (вероятность \(\frac{2}{3}\)), и после открытия ведущим другой двери с козой, оставшаяся дверь будет с автомобилем.
	\end{itemize}
	
	Таким образом, вероятность выигрыша при смене равна \(\frac{2}{3}\).
	
	Альтернативно, можно посчитать по дереву:
	\[ \Pr(\text{выигрыш при смене}) = 6 \cdot \frac{1}{18} + 6 \cdot \frac{1}{9} = \frac{2}{3} \]
	(где \(\frac{1}{9} = \frac{1}{3} \cdot \frac{1}{3} \cdot 1\) — случаи, когда игрок изначально ошибся).
	
	Вероятность выигрыша при сохранении выбора равна, соответственно, \(\frac{1}{3}\). \qed
	
	\begin{notebook}{13.10: Интуитивное объяснение}
		Парадокс заключается в том, что многие люди ошибочно полагают, что после открытия одной двери вероятности для двух оставшихся дверей становятся равными (\(\frac{1}{2}\) и \(\frac{1}{2}\)). Однако это не так, потому что ведущий всегда открывает дверь с козой, что даёт игроку дополнительную информацию.
		
		Ключевой момент: вероятность того, что первоначальный выбор был правильным, остаётся \(\frac{1}{3}\), а вероятность того, что автомобиль за одной из двух других дверей, была \(\frac{2}{3}\). После того как ведущий открывает дверь с козой, вся эта вероятность \(\frac{2}{3}\) переходит на оставшуюся закрытую дверь.
	\end{notebook}
	
	
	\subsubsection{Условная вероятность}
	
	\begin{definitionbox}{13.6: Условная вероятность}
		Пусть \((S, \Pr)\) — вероятностное пространство, \(A, B \subseteq S\), и \(\Pr(B) \neq 0\).
		
		\textbf{Условной вероятностью} события \(A\) при условии, что событие \(B\) произошло, называется величина:
		\[ \Pr(A \mid B) = \frac{\Pr(A \cap B)}{\Pr(B)} \]
	\end{definitionbox}
	
	\myremark{Смысл:} Условная вероятность показывает, какова вероятность события \(A\) с учётом того, что нам уже известно о наступлении события \(B\).
	
	\begin{propertybox}{13.7: Свойства условной вероятности}
		\begin{itemize}
			\item \(\Pr(A \cap B) = \Pr(A \mid B) \cdot \Pr(B) = \Pr(B \mid A) \cdot \Pr(A)\) (формула умножения вероятностей);
			\item Если события \(A\) и \(B\) независимы, то \(\Pr(A \mid B) = \Pr(A)\);
			\item \(\Pr(\Omega \mid B) = 1\);
			\item \(\Pr(\emptyset \mid B) = 0\);
			\item Для любого разбиения \(\bigcup_{i=1}^n B_i = S\) с \(\Pr(B_i) > 0\) справедлива формула полной вероятности:
			\[ \Pr(A) = \sum_{i=1}^n \Pr(A \mid B_i) \cdot \Pr(B_i) \]
		\end{itemize}
	\end{propertybox}
	
	
	
	
	
	
	
	
	
	
	% ==================== РАЗДЕЛ 14: ЗАКОН ПОЛНОЙ ВЕРОЯТНОСТИ И ФОРМУЛА БАЙЕСА ====================
	\newpage
	\section{Закон полной вероятности и формула Байеса}
	
	\subsection{Закон полной вероятности. Формула Байеса. Дискретная случайная величина, арифметические операции над случайными величинами (14 вопрос)}
	
	\subsubsection{Закон полной вероятности}
	
	\begin{definitionbox}{14.1: Полная система событий}
		\textbf{Полной системой событий} называется не более чем счётное множество событий \(B_1, B_2, \dots, B_n\), удовлетворяющее условиям:
		\begin{itemize}
			\item все события попарно несовместны: \(B_i \cap B_j = \emptyset\) при \(i \neq j\);
			\item их объединение образует всё пространство элементарных событий: \(\bigcup_{i=1}^n B_i = \Omega\);
			\item \(\Pr(B_i) > 0\) для всех \(i\).
		\end{itemize}
	\end{definitionbox}
	
	\begin{theorembox}{14.2: Закон полной вероятности}
		Пусть событие \(A \subseteq \Omega\) может произойти только вместе с одним из событий \(B_1, B_2, \dots, B_n\), образующих полную систему событий. Тогда вероятность события \(A\) равна сумме произведений вероятностей гипотез на условные вероятности события, вычисленные соответственно при каждой из гипотез:
		
		\[ \Pr(A) = \sum_{i=1}^n \Pr(A \mid B_i) \cdot \Pr(B_i) \]
	\end{theorembox}
	
	\myproof
	Так как события \(B_1, B_2, \dots, B_n\) образуют полную систему событий, то по определению событие \(A\) можно представить следующим образом:
	\[ A = A \cap \Omega = A \cap \left( \bigcup_{i=1}^n B_i \right) = \bigcup_{i=1}^n (A \cap B_i) \]
	
	События \(B_1, B_2, \dots, B_n\) попарно несовместны, значит, события \((A \cap B_i)\) тоже несовместны. Тогда, воспользовавшись определением условной вероятности, получаем:
	\[ \Pr(A) = \Pr\left( \bigcup_{i=1}^n (A \cap B_i) \right) = \sum_{i=1}^n \Pr(A \cap B_i) = \sum_{i=1}^n \Pr(A \mid B_i) \cdot \Pr(B_i) \]
	\qed
	
	\begin{examplebox}{14.3: Пример использования закона полной вероятности}
		\textbf{Задача:} Имеются 3 одинаковые урны с шарами. В первой из них находится 3 белых и 4 чёрных шара, во второй — 2 белых и 5 чёрных, а в третьей — 10 чёрных шаров. Из случайно выбранной урны наудачу вынут шар. С какой вероятностью он окажется белым?\\
		
		\textbf{Решение:} Будем считать события \(B_1, B_2, B_3\) выбором урны с соответствующим номером, а событие \(A\) — выбором белого шара.\\
		
		По условию задачи все события выбора урны равновероятны, значит:
		\[ \Pr(B_1) = \Pr(B_2) = \Pr(B_3) = \frac{1}{3} \]
		
		Теперь найдём вероятность события \(A\) при выборе каждой урны:
		\[ \Pr(A \mid B_1) = \frac{3}{7}, \quad \Pr(A \mid B_2) = \frac{2}{7}, \quad \Pr(A \mid B_3) = 0 \]
		
		В результате по формуле полной вероятности получаем:
		\[ \Pr(A) = \frac{1}{3} \cdot \frac{3}{7} + \frac{1}{3} \cdot \frac{2}{7} + \frac{1}{3} \cdot 0 = \frac{5}{21} \approx 0.238 \]
	\end{examplebox}
	
	\subsubsection{Формула Байеса}
	
	\begin{definitionbox}{14.4: Формула Байеса}
		\textbf{Формула Байеса} позволяет переоценить вероятности гипотез \(B_i\) после того, как стало известно, что событие \(A\) произошло:
		
		\[ \Pr(B_i \mid A) = \frac{\Pr(A \mid B_i) \cdot \Pr(B_i)}{\sum_{j=1}^n \Pr(A \mid B_j) \cdot \Pr(B_j)} \]
		
		где:
		\begin{itemize}
			\item \(\Pr(A)\) — вероятность события \(A\);
			\item \(\Pr(A \mid B)\) — вероятность события \(A\) при наступлении события \(B\);
			\item \(\Pr(B \mid A)\) — вероятность наступления события \(B\) при условии, что событие \(A\) произошло;
			\item \(\Pr(B)\) — априорная вероятность наступления события \(B\).
		\end{itemize}
	\end{definitionbox}
	
	\myproof
	Из определения условной вероятности следует:
	\[ \Pr(B_i \cap A) = \Pr(A \cap B_i) = \Pr(A \mid B_i) \cdot \Pr(B_i) \]
	и
	\[ \Pr(B_i \mid A) = \frac{\Pr(B_i \cap A)}{\Pr(A)} \]
	
	Подставляя первое равенство и используя закон полной вероятности для \(\Pr(A)\), получаем:
	\[ \Pr(B_i \mid A) = \frac{\Pr(A \mid B_i) \cdot \Pr(B_i)}{\sum_{j=1}^n \Pr(A \mid B_j) \cdot \Pr(B_j)} \]
	\qed
	
	\begin{examplebox}{14.5: Пример использования формулы Байеса}
		\textbf{Задача:} Пусть вероятность заразиться гриппом равна 0.01. После проведения анализа вероятность того, что положительный результат действительно указывает на грипп, равна 0.9, а вероятность ложного положительного результата при другой болезни равна 0.001. Найти вероятность того, что при положительном анализе у пациента действительно грипп. \\
		
		\textbf{Решение:} Определим события:
		\begin{itemize}
			\item \(A\) — анализ положительный;
			\item \(B_1\) — пациент болен гриппом;
			\item \(B_2\) — пациент болен другой болезнью (или здоров).
		\end{itemize}
		
		Априорные вероятности:
		\[ \Pr(B_1) = 0.01, \quad \Pr(B_2) = 0.99 \]
		
		Условные вероятности:
		\[ \Pr(A \mid B_1) = 0.9, \quad \Pr(A \mid B_2) = 0.001 \]
		
		По формуле Байеса:
		\[ \Pr(B_1 \mid A) = \frac{0.9 \cdot 0.01}{0.9 \cdot 0.01 + 0.001 \cdot 0.99} = \frac{0.009}{0.009 + 0.00099} = \frac{0.009}{0.00999} = \frac{900}{999} = \frac{100}{111} \approx 0.9 \]
		
		Таким образом, даже при положительном анализе вероятность того, что пациент действительно болен гриппом, составляет около 90\% (а не 100\%), что объясняется редкостью заболевания и возможностью ложных срабатываний.
	\end{examplebox}
	
	\subsubsection{Дискретные случайные величины}
	
	\begin{definitionbox}{14.6: Дискретная случайная величина}
		Пусть \((S, \Pr)\) — вероятностное пространство, \(|S| < \infty\). Функция \(\xi : S \to \mathbb{R}\) называется \textbf{дискретной случайной величиной} (ДСВ). \\
		
		В более общем случае, дискретная случайная величина — это случайная величина, множество значений которой конечно или счётно.
	\end{definitionbox}
	
	\begin{definitionbox}{14.7: Вероятность значений ДСВ}
		Для дискретной случайной величины \(\xi\) определяют:
		\begin{itemize}
			\item \(\Pr\{\xi = a\} = \Pr\{\omega \in S : \xi(\omega) = a\}\) — вероятность того, что \(\xi\) примет значение \(a\);
			\item \(\Pr\{\xi \leq a\} = \Pr\{\omega \in S : \xi(\omega) \leq a\}\) — вероятность того, что \(\xi\) не превосходит \(a\).
		\end{itemize}
		
		То есть это суммарная вероятность всех элементарных исходов \(\omega \in S\), для которых значение \(\xi(\omega)\) удовлетворяет соответствующему условию.
	\end{definitionbox}
	
	\subsubsection{Способы задания дискретных случайных величин}
	
	\begin{definitionbox}{14.8: Способы задания ДСВ}
		Существуют два основных способа задания дискретной случайной величины:
		
		\textbf{1. Аналитический способ:}
		\[ \Pr\{\xi = a\} = p \]
		
		\textbf{2. Табличный способ (закон распределения ДСВ):}
		\[ \xi = \begin{pmatrix}
			a_1 & a_2 & \cdots & a_n \\
			p_1 & p_2 & \cdots & p_n
		\end{pmatrix} \]
		где \(a_i\) — возможные значения случайной величины, а \(p_i = \Pr\{\xi = a_i\}\) — соответствующие вероятности.
		
		При этом выполняется условие нормировки:
		\[ \sum_{i=1}^n p_i = 1 \]
	\end{definitionbox}
	
	\begin{examplebox}{14.9: Примеры дискретных случайных величин}
		
		\textbf{Пример 1:} Для вероятностного пространства, описывающего три броска монеты, можно ввести дискретную случайную величину \(\phi\), отражающую количество выпавших орлов:
		\[ \Pr\{\phi = 2\} = \Pr(\{OOP, OPO, POO\}) = \frac{3}{8} \]
		
		\textbf{Пример 2:} С помощью дискретной случайной величины \(\xi\) можно выразить число, выпавшее при броске шестигранного игрального кубика:
		\[ \xi = \begin{pmatrix}
			1 & 2 & 3 & 4 & 5 & 6 \\
			\frac{1}{6} & \frac{1}{6} & \frac{1}{6} & \frac{1}{6} & \frac{1}{6} & \frac{1}{6}
		\end{pmatrix} \]
		
		\textbf{Пример 3:} Несимметричный кубик:
		\[ \xi = \begin{pmatrix}
			1 & 2 & 3 & 4 & 5 & 6 \\
			\frac{1}{3} & \frac{1}{6} & \frac{1}{6} & \frac{1}{6} & \frac{1}{6} & \frac{1}{6}
		\end{pmatrix} \]
	\end{examplebox}
	
	\subsubsection{Арифметические операции над дискретными случайными величинами}
	
	\begin{propertybox}{14.10: Арифметические операции над ДСВ}
		Пусть \(\xi\) — дискретная случайная величина, \(c \in \mathbb{R}\) — константа. Определим новые случайные величины:
		
		\textbf{1. Сдвиг на константу:} \(\eta(\omega) = \xi(\omega) + c\)
		\[ \Pr\{\eta = a\} = \Pr\{\xi + c = a\} = \Pr\{\xi = a - c\} \]
		
		\textbf{2. Умножение на константу:} \(\eta(\omega) = \xi(\omega) \cdot c\), при \(c \neq 0\)
		\[ \Pr\{\eta = a\} = \Pr\{\xi \cdot c = a\} = \Pr\left\{\xi = \frac{a}{c}\right\} \]
		
		\textbf{3. Возведение в квадрат:} \(\eta(\omega) = (\xi(\omega))^2\)
		\[ \Pr\{\eta = a\} = \Pr\{\xi^2 = a\} = \begin{cases}
			\Pr\{\xi = \sqrt{a}\} + \Pr\{\xi = -\sqrt{a}\}, & \text{если } a > 0 \text{ и } \sqrt{a} \in \text{область значений } \xi \\
			\Pr\{\xi = 0\}, & \text{если } a = 0 \\
			0, & \text{иначе}
		\end{cases} \]
	\end{propertybox}
	
	\begin{examplebox}{14.11: Примеры операций над ДСВ}
		Пусть \(\xi\) задана таблично:
		\[ \xi = \begin{pmatrix}
			1 & 2 & 3 \\
			0.3 & 0.5 & 0.2
		\end{pmatrix} \]
		
		Тогда:
		\begin{itemize}
			\item \(\eta = \xi + 2\):
			\[ \eta = \begin{pmatrix}
				3 & 4 & 5 \\
				0.3 & 0.5 & 0.2
			\end{pmatrix} \]
			
			\item \(\eta = 2\xi\):
			\[ \eta = \begin{pmatrix}
				2 & 4 & 6 \\
				0.3 & 0.5 & 0.2
			\end{pmatrix} \]
			
			\item \(\eta = \xi^2\):
			\[ \eta = \begin{pmatrix}
				1 & 4 & 9 \\
				0.3 & 0.5 & 0.2
			\end{pmatrix} \]
		\end{itemize}
	\end{examplebox}
	
	\begin{notebook}{14.12: Замечание}
		Для двух случайных величин \(\xi\) и \(\eta\) можно определить их сумму \(\xi + \eta\), произведение \(\xi \cdot \eta\) и другие операции. Однако в общем случае распределение суммы не является просто комбинацией распределений слагаемых, если величины не являются независимыми.
	\end{notebook}
	
	
	
	
	
	% ==================== РАЗДЕЛ 15: МАТЕМАТИЧЕСКОЕ ОЖИДАНИЕ И ДИСПЕРСИЯ ====================
	\newpage
	\section{Математическое ожидание и дисперсия}
	
	\subsection{Математическое ожидание и дисперсия дискретной случайной величины, их свойства. Схема Бернулли (15 вопрос)}
	
	\subsubsection{Математическое ожидание}
	
	\begin{definitionbox}{15.1: Математическое ожидание ДСВ}
		\textbf{Математическим ожиданием} дискретной случайной величины \(\xi\) на вероятностном пространстве \((S, \Pr)\) называется среднее значение случайной величины, равное сумме произведений всех возможных значений случайной величины на их вероятности:
		
		\[ E\xi = \sum_{\omega \in S} \xi(\omega) \cdot \Pr(\{\omega\}) \]
		
		Говоря простым языком, это \textit{среднеожидаемое значение} при многократном повторении испытаний.
	\end{definitionbox}
	
	\begin{remarkbox}{15.2: Альтернативная формула}
		Математическое ожидание можно также выразить через закон распределения случайной величины:
		
		\[ E\xi = \sum_{a \in \text{Im}(\xi)} a \cdot \Pr\{\xi = a\} \]
		
		где \(\Pr\{\xi = a\} = \Pr\{\omega \in S : \xi(\omega) = a\}\).
		
		\textbf{Вывод:}
		\begin{align*}
			E\xi &= \sum_{\omega \in S} \xi(\omega) \cdot \Pr(\{\omega\}) \\
			&= \sum_{a \in \text{Im}(\xi)} \sum_{\omega: \xi(\omega)=a} a \cdot \Pr(\{\omega\}) \\
			&= \sum_{a \in \text{Im}(\xi)} a \cdot \sum_{\omega: \xi(\omega)=a} \Pr(\{\omega\}) \\
			&= \sum_{a \in \text{Im}(\xi)} a \cdot \Pr\{\xi = a\}
		\end{align*}
	\end{remarkbox}
	
	\begin{center}
		\begin{tikzpicture}
			\draw[thick, blue] (0,0) -- (8,0);
			\foreach \x in {0,1,2,3} {
				\draw[thick, red] (\x*2,0) -- (\x*2, {(\x==0?1:\x==1?3:\x==2?3:1)/4});
				\draw (\x*2,0) node[below] {\x};
			}
			\draw[thick, green, <->] (0,0.5) -- (6,0.5);
			\node at (3,0.8) {Математическое ожидание};
			\node at (3,-1) {Иллюстрация: мат. ожидание как точка равновесия};
		\end{tikzpicture}
		\captionof{figure}{Математическое ожидание как центр масс распределения вероятностей}
	\end{center}
	
	\subsubsection{Свойства математического ожидания}
	
	\begin{propertybox}{15.3: Свойства математического ожидания}
		
		\textbf{1. Математическое ожидание постоянной величины:}
		\[ E[a] = a, \quad \text{где } a \in \mathbb{R} \text{ — константа} \]
		
		\textbf{2. Линейность математического ожидания:}
		\[ E[aX + bY] = aE[X] + bE[Y] \]
		
		\textbf{3. Математическое ожидание произведения независимых случайных величин:}
		\[ E[XY] = E[X] \cdot E[Y] \]
		
		\textbf{4. Сдвиг на константу:} для \(\eta(\omega) = \xi(\omega) + c\)
		\[ E\eta = c + E\xi \]
		
		\textbf{5. Умножение на константу:} для \(\eta(\omega) = \xi(\omega) \cdot c\), \(c \neq 0\)
		\[ E\eta = c \cdot E\xi \]
		
		\textbf{6. Квадрат случайной величины:} для \(\eta(\omega) = (\xi(\omega))^2\)
		\[ E\eta = \sum_{\omega \in S} \xi^2(\omega) \cdot \Pr(\{\omega\}) \]
		
		\textbf{7. Математическое ожидание суммы:} для \(\eta = \xi + \psi\)
		\[ E(\xi + \psi) = \sum_{\omega \in S} (\xi(\omega) + \psi(\omega)) \cdot \Pr(\{\omega\}) = E\xi + E\psi \]
	\end{propertybox}
	
	\subsubsection{Дисперсия}
	
	\begin{definitionbox}{15.4: Дисперсия ДСВ}
		\textbf{Дисперсией} дискретной случайной величины \(\xi\) называется математическое ожидание квадрата разности значения случайной величины и её математического ожидания:
		
		\[ D\xi = E(\xi - E\xi)^2 \]
		
		Дисперсия показывает, насколько в среднем значения сгруппированы около математического ожидания. Если дисперсия маленькая — значения сравнительно близки друг к другу, если большая — далеки друг от друга.
	\end{definitionbox}
	
	\begin{theorembox}{15.5: Формула для вычисления дисперсии}
		\[ D\xi = E\xi^2 - (E\xi)^2 \]
	\end{theorembox}
	
	\myproof
	\begin{align*}
		D\xi &= E(\xi - E\xi)^2 = E(\xi^2 - 2 \cdot E\xi \cdot \xi + (E\xi)^2) \\
		&= E\xi^2 - E(2 \cdot E\xi \cdot \xi) + E((E\xi)^2) \\
		&= E\xi^2 - 2 \cdot E\xi \cdot E\xi + (E\xi)^2 \\
		&= E\xi^2 - (E\xi)^2
	\end{align*}
	\qed
	
	\begin{examplebox}{15.6: Пример вычисления дисперсии}
		Сравним дисперсию двух законов распределения:
		
		\[ X: \begin{pmatrix} 1 & 2 \\ 0.5 & 0.5 \end{pmatrix} \quad\text{и}\quad Y: \begin{pmatrix} -10 & 10 \\ 0.5 & 0.5 \end{pmatrix} \]
		
		Вычислим математические ожидания:
		\[ EX = 1 \cdot 0.5 + 2 \cdot 0.5 = 1.5 \]
		\[ EY = (-10) \cdot 0.5 + 10 \cdot 0.5 = 0 \]
		
		Вычислим вторые моменты:
		\[ EX^2 = 1^2 \cdot 0.5 + 2^2 \cdot 0.5 = 0.5 + 2 = 2.5 \]
		\[ EY^2 = (-10)^2 \cdot 0.5 + 10^2 \cdot 0.5 = 100 \]
		
		Теперь дисперсии:
		\[ DX = EX^2 - (EX)^2 = 2.5 - 1.5^2 = 2.5 - 2.25 = 0.25 \]
		\[ DY = EY^2 - (EY)^2 = 100 - 0^2 = 100 \]
		
		Как видим, дисперсия \(Y\) значительно больше, что соответствует разбросу значений.
	\end{examplebox}
	
	\subsubsection{Свойства дисперсии}
	
	\begin{propertybox}{15.7: Свойства дисперсии}
		
		\textbf{1. Неотрицательность:}
		\[ D\xi \ge 0 \]
		
		\textbf{2. Конечность:}
		Если дисперсия случайной величины конечна, то конечно и её математическое ожидание.
		
		\textbf{3. Дисперсия константы:}
		\[ D[a] = 0 \]
		
		\textbf{4. Вынесение константы:}
		\[ D[c \cdot X] = c^2 \cdot D[X] \]
		
		\textbf{5. Сдвиг на константу:}
		\[ D[X + c] = D[X] \]
		
		\textbf{6. Дисперсия суммы двух случайных величин:}
		\[ D[X + Y] = D[X] + D[Y] + 2 \cdot \text{cov}(X, Y) \]
		где \(\text{cov}(X, Y)\) — ковариация (корреляционный момент).
		
		\textbf{7. Дисперсия суммы независимых случайных величин:}
		Если \(X\) и \(Y\) независимы, то \(\text{cov}(X, Y) = 0\) и
		\[ D[X + Y] = D[X] + D[Y] \]
	\end{propertybox}
	
	\subsubsection{Схема Бернулли}
	
	\begin{definitionbox}{15.8: Схема Бернулли}
		\textbf{Схема Бернулли} — это последовательность независимых испытаний, в каждом из которых возможны лишь два исхода:
		\begin{itemize}
			\item "успех" с вероятностью \(p\);
			\item "неудача" с вероятностью \(q = 1 - p\).
		\end{itemize}
	\end{definitionbox}
	
	\begin{definitionbox}{15.9: Вероятностное пространство в схеме Бернулли}
		Пространство элементарных событий для \(n\) испытаний:
		\[ \Omega = \{ (a_1, a_2, \dots, a_n) \mid a_i \in \{0, 1\}, i = 1, \dots, n \} \]
		где \(a_i = 1\) соответствует успеху в \(i\)-м испытании, \(a_i = 0\) — неудаче.
		
		Мощность пространства: \(|\Omega| = 2^n\).
		
		Вероятность любого элементарного события \(\omega = (a_1, \dots, a_n) \in \Omega\):
		\[ \Pr(\omega) = p^{\sum_{i=1}^n a_i} \cdot (1-p)^{n - \sum_{i=1}^n a_i} \]
		
		Если в исходе \(\omega\) наблюдается \(k\) успехов, то:
		\[ \Pr(\omega) = p^k (1-p)^{n-k} \]
	\end{definitionbox}
	
	\begin{theorembox}{15.10: Формула Бернулли}
		Пусть событие \(A_k = \{ \omega \in \Omega \mid \sum_{i=1}^n a_i = k \}\) — ровно \(k\) успехов в \(n\) испытаниях. Тогда:
		
		\[ \Pr(A_k) = \sum_{\omega \in A_k} \Pr(\omega) = C_n^k \cdot p^k \cdot (1-p)^{n-k} \]
	\end{theorembox}
	
	\myproof
	Количество способов выбрать \(k\) позиций для успехов из \(n\) равно \(C_n^k\). Для каждого такого набора позиций вероятность равна \(p^k (1-p)^{n-k}\). Так как эти события несовместны, вероятность их объединения равна сумме вероятностей. \qed
	
	\begin{examplebox}{15.11: Задача на схему Бернулли}
		\textbf{Задача:} Вероятность попадания в цель при каждом выстреле из лука равна \(\frac{1}{3}\). Производится 6 выстрелов. 
		\begin{itemize}
			\item Какова вероятность ровно двух попаданий?
			\item Какова вероятность не менее двух попаданий?
		\end{itemize}
		
		\textbf{Решение:} Обозначим \(A = \{\text{попадание при одном выстреле}\}\), \(p = \Pr(A) = \frac{1}{3}\), \(q = 1 - p = \frac{2}{3}\). Число выстрелов \(n = 6\).
		
		Естественно предположить, что выстрелы независимы друг от друга.
		
		\textbf{1. Вероятность ровно двух попаданий} находим по формуле Бернулли:
		\[ \Pr_6(2) = C_6^2 \cdot \left(\frac{1}{3}\right)^2 \cdot \left(\frac{2}{3}\right)^4 = 15 \cdot \frac{1}{9} \cdot \frac{16}{81} = 15 \cdot \frac{16}{729} = \frac{240}{729} = \frac{80}{243} \approx 0.329 \]
		
		\textbf{2. Вероятность не менее двух попаданий:}
		\[ \Pr_6(\ge 2) = 1 - \Pr_6(0) - \Pr_6(1) \]
		\[ \Pr_6(0) = C_6^0 \cdot \left(\frac{1}{3}\right)^0 \cdot \left(\frac{2}{3}\right)^6 = 1 \cdot 1 \cdot \frac{64}{729} = \frac{64}{729} \]
		\[ \Pr_6(1) = C_6^1 \cdot \left(\frac{1}{3}\right)^1 \cdot \left(\frac{2}{3}\right)^5 = 6 \cdot \frac{1}{3} \cdot \frac{32}{243} = 6 \cdot \frac{32}{729} = \frac{192}{729} \]
		\[ \Pr_6(\ge 2) = 1 - \frac{64}{729} - \frac{192}{729} = 1 - \frac{256}{729} = \frac{473}{729} \approx 0.649 \]
	\end{examplebox}
	
	--- доп ---
	
	\begin{definitionbox}{15.12: Наивероятнейшее число успехов}
		Число \(k_0\) называется \textbf{наивероятнейшим числом успехов} в схеме Бернулли, если \(\Pr(A_{k_0}) \ge \Pr(A_k)\) для всех \(k\). Оно определяется из неравенства:
		\[ np - q \le k_0 \le np + p \]
		
		где \(k_0\) — целое число в этом интервале.
	\end{definitionbox}
	
	\begin{examplebox}{15.13: Наивероятнейшее число успехов}
		Для задачи с луком:
		\[ np - q = 6 \cdot \frac{1}{3} - \frac{2}{3} = 2 - \frac{2}{3} = \frac{4}{3} \approx 1.33 \]
		\[ np + p = 2 + \frac{1}{3} = \frac{7}{3} \approx 2.33 \]
		
		Наивероятнейшее число попаданий лежит в пределах от \(1\frac{1}{3}\) до \(2\frac{1}{3}\), следовательно, оно равно \(2\).
	\end{examplebox}
	
	\subsubsection{Математическое ожидание и дисперсия в схеме Бернулли}
	
	\begin{theorembox}{15.14: Моменты биномиального распределения}
		Пусть \(\xi\) — число успехов в \(n\) испытаниях Бернулли с вероятностью успеха \(p\). Тогда:
		
		\[ E\xi = np \]
		\[ D\xi = npq = np(1-p) \]
	\end{theorembox}
	
	\myproof
	Представим \(\xi\) как сумму индикаторов успеха в каждом испытании:
	\[ \xi = \xi_1 + \xi_2 + \dots + \xi_n \]
	где \(\xi_i\) — индикатор успеха в \(i\)-м испытании (\(1\) — успех, \(0\) — неудача).
	
	Для каждого \(\xi_i\):
	\[ E\xi_i = 1 \cdot p + 0 \cdot q = p \]
	\[ E\xi_i^2 = 1^2 \cdot p + 0^2 \cdot q = p \]
	\[ D\xi_i = E\xi_i^2 - (E\xi_i)^2 = p - p^2 = p(1-p) = pq \]
	
	В силу линейности математического ожидания:
	\[ E\xi = \sum_{i=1}^n E\xi_i = np \]
	
	В силу независимости испытаний (и, следовательно, независимости \(\xi_i\)):
	\[ D\xi = \sum_{i=1}^n D\xi_i = npq \]
	\qed
	
	\begin{center}
		\begin{tabular}{|c|c|c|c|}
			\hline
			\(n\) & \(k\) & \(p\) & \(\Pr = C_n^k p^k (1-p)^{n-k}\) \\ \hline
			6 & 0 & 1/3 & 64/729 \\
			6 & 1 & 1/3 & 192/729 \\
			6 & 2 & 1/3 & 240/729 \\
			6 & 3 & 1/3 & 160/729 \\
			6 & 4 & 1/3 & 60/729 \\
			6 & 5 & 1/3 & 12/729 \\
			6 & 6 & 1/3 & 1/729 \\ \hline
		\end{tabular}
		\captionof{table}{Распределение числа попаданий при 6 выстрелах с вероятностью попадания 1/3}
	\end{center}
	
	
	
	
	
	% ==================== РАЗДЕЛ 16: ДАТЧИК СЛУЧАЙНЫХ ЧИСЕЛ ====================
	\newpage
	\section{Датчик случайных чисел}
	
	\subsection{Датчик случайных чисел, псевдослучайные числа. Моделирование дискретного распределения (табличный метод) (16 вопрос)}
	
	\subsubsection{Датчики случайных чисел}
	
	\begin{definitionbox}{16.1: Модель дискретной случайной величины}
		Многократно повторяемый эксперимент, исходы которого имеют заданный набор вероятностей, соответствующий набору вероятностей дискретной случайной величины, называют \textbf{моделью этой ДСВ}. То есть заранее известно, какая вероятность будет у каждого исхода.
	\end{definitionbox}
	
	\begin{examplebox}{16.2: Пример модели ДСВ}
		\textbf{Хороший пример:} Моделью \(\xi\) с распределением
		\[ \Pr\{\xi = 0\} = \frac{1}{2}, \quad \Pr\{\xi = 1\} = \frac{1}{2} \]
		можно считать бросок монетки (если считать, что стороны монетки одинаковые, монетка никогда не падает на ребро и бросают её наугад).
	\end{examplebox}
	
	\begin{definitionbox}{16.3: Дискретный датчик случайных чисел}
		Многократно повторяемый эксперимент, результатом которого является (псевдо)случайно выбранное число из некоторого конечного множества, называют \textbf{дискретным генератором (датчиком) случайных чисел}.
	\end{definitionbox}
	
	\begin{definitionbox}{16.4: Непрерывный датчик случайных чисел}
		Многократно повторяемый эксперимент, результатом которого является (псевдо)случайно выбранное число из некоторого промежутка числовой прямой, называют \textbf{непрерывным генератором (датчиком) случайных чисел}.
	\end{definitionbox}
	
	\subsubsection{Псевдослучайные числа}
	
	\begin{importantbox}{16.5: Псевдослучайные числа}
		То, что числа называются \textbf{псевдослучайными}, а не \textbf{случайными}, означает, что они получены алгоритмом, который не является полностью случайным, и который теоретически или практически можно повторить с той же последовательностью.
		
		На практике генераторы псевдослучайных чисел используются повсеместно в компьютерном моделировании, так как получить истинно случайные числа сложно и дорого.
	\end{importantbox}
	
	\begin{examplebox}{16.6: Пример дискретного датчика}
		Назовём \(\alpha\) дискретный генератор случайных чисел, для которого
		\[ \Pr\{\alpha = k\} = \frac{1}{N} \quad \forall k \in \{1, 2, \dots, N\}. \]
		
		Такой датчик моделирует равномерное дискретное распределение на множестве из \(N\) элементов.
	\end{examplebox}
	
	\begin{examplebox}{16.7: Пример непрерывного датчика}
		Назовём \(\alpha_0\) непрерывный генератор случайных чисел, который с равной вероятностью попадает в каждую точку отрезка \([0, 1]\).
		
		Поскольку отрезок \([0, 1]\) содержит несчётное количество точек, вероятность
		\[ \Pr\{\alpha_0 = x\} = \frac{1}{\infty} = 0 \quad \forall x \in [0, 1]. \]
		
		Однако вероятность того, что точка попадёт в некоторый промежуток, не нулевая:
		\[ \Pr\{\alpha_0 \in [a, a+\delta]\} = \frac{|[a, a+\delta]|}{|[0,1]|} = \frac{\delta}{1} = \delta. \]
		
		В частности,
		\[ \Pr\{\alpha_0 \in [0, 1]\} = 1. \]
	\end{examplebox}
	
	\begin{remarkbox}{16.8: Замечание о мере нуля}
		Поскольку для конкретной точки вероятность попадания туда \(\alpha_0\) равна нулю, можно смело выкидывать из множества конечное (и даже счётное) количество точек, и вероятность попадания в него останется неизменной. Следовательно, вместо \([a, b]\) можно рассматривать \((a, b)\) или \([a, b)\). Последнее особенно удобно, поскольку полуинтервалы хорошо стыкуются.
	\end{remarkbox}
	
	\begin{notebook}{16.9: Наглядное пояснение разницы между дискретным и непрерывным датчиками}
		\textbf{Дискретный датчик} имеет конечное число возможных исходов. Например, сегодня может пойти дождь, выпасть снег, или то и другое вместе, или ничего не произойдёт. Имеем всего 4 случая, и мы знаем вероятность каждого из них. \\
		
		\textbf{Непрерывный датчик} — это попытка угадать, когда пойдёт дождь. У нас есть промежуток с 00:00 до 23:59. Так как время может измеряться секундами, миллисекундами, микросекундами и т.д., дать точный ответ становится практически невозможно (время можно делить бесконечно). Поэтому проще сказать, что дождь пойдёт в период с 12:00 до 14:00. Тогда вероятность равна \(\frac{2}{24} = \frac{1}{12}\).
	\end{notebook}
	
	\subsubsection{Табличный метод моделирования дискретной случайной величины}
	
	\begin{definitionbox}{16.10: Постановка задачи}
		Пусть \((S, \Pr)\) — вероятностное пространство, \(\xi\) — заданная на нём дискретная случайная величина.
		
		При этом \(S = \{0, 1, \dots, n-1\}\), и для всех \(i \in \{0, 1, \dots, n-1\}\) заданы вероятности \(\Pr(\{i\}) = p_i\), причём
		\[ \sum_{i=0}^{n-1} p_i = 1. \]
		
		\textbf{Примечание:} Для ДСВ с произвольным набором значений можно их пронумеровать и свести задачу к данной.
	\end{definitionbox}
	
	\begin{theorembox}{16.11: Табличный метод моделирования ДСВ}
		Возьмём отрезок \([0, 1]\) и разделим его на полуинтервалы вида \([x_i, x_{i+1})\), где
		\[ x_0 = 0, \quad \forall i \in \{0, \dots, n-1\}: x_{i+1} = x_i + p_i. \]
		
		Для такого разбиения отрезка \([0, 1]\) вероятность попадания непрерывного датчика \(\alpha_0\) в полуинтервал \([x_i, x_{i+1})\) равна:
		\[ \Pr\{\alpha_0 \in [x_i, x_{i+1})\} = p_i. \]
		
		Таким образом, эксперимент "в полуинтервал с каким номером попадёт случайное число, выданное \(\alpha_0\)?" моделирует случайную величину \(\xi\).
	\end{theorembox}
	
	\begin{center}
		\begin{tikzpicture}[scale=1.2]
			\draw[thick, blue] (0,0) -- (10,0);
			
			\foreach \x/\label in {0/0, 2/p_0, 4/p_1, 6/p_2, 8/p_3, 10/1} {
				\draw (\x,0.1) -- (\x,-0.1);
				\node[below] at (\x,-0.1) {\(\label\)};
			}
			
			\node[above] at (5,0.2) {Разбиение отрезка \([0,1]\) на полуинтервалы};
			
			\foreach \x/\label in {0/p_0, 2/p_1, 4/p_2, 6/p_3, 8/p_4} {
				\draw[decorate, decoration={brace, mirror, amplitude=10pt}] (\x,-0.5) -- ({\x+2},-0.5);
				\node[below] at ({\x+1},-1) {\(\label\)};
			}
		\end{tikzpicture}
		\captionof{figure}{Табличный метод: разбиение отрезка \([0,1]\) на полуинтервалы длиной \(p_i\)}
	\end{center}
	
	\begin{examplebox}{16.12: Пример табличного метода}
		Пусть требуется смоделировать дискретную случайную величину с распределением:
		\[ p_0 = 0.2, \quad p_1 = 0.3, \quad p_2 = 0.1, \quad p_3 = 0.4. \]
		
		Тогда строим разбиение:
		\begin{align*}
			x_0 &= 0 \\
			x_1 &= x_0 + p_0 = 0.2 \\
			x_2 &= x_1 + p_1 = 0.5 \\
			x_3 &= x_2 + p_2 = 0.6 \\
			x_4 &= x_3 + p_3 = 1.0
		\end{align*}
		
		Получаем полуинтервалы:
		\[ [0, 0.2), \quad [0.2, 0.5), \quad [0.5, 0.6), \quad [0.6, 1) \]
		
		Генерируем случайное число \(r \in [0,1)\) от непрерывного датчика. В зависимости от того, в какой полуинтервал попало \(r\), выбираем соответствующее значение случайной величины.
	\end{examplebox}
	
	\begin{importantbox}{16.13: Важное замечание о граничных точках}
		При попадании в точку \(1\) значение \(\alpha_0\) оказывается вне пределов всех полуинтервалов. С одной стороны, вероятность такого события равна \(0\), следовательно на распределение вероятностей это не повлияет. С другой стороны, если это всё же произошло, мы можем просто объявить эксперимент неудачным и провести его заново.
	\end{importantbox}
	
	\begin{remarkbox}{16.14: Недостатки табличного метода}
		\textbf{Основным недостатком} данного метода является большое количество сравнений вещественных чисел, поскольку эта операция является не очень точной и достаточно долгой. Для каждого сгенерированного числа приходится последовательно проверять, в какой полуинтервал оно попало, что в худшем случае требует \(O(n)\) операций сравнения.
		
		Существуют более эффективные методы моделирования дискретных распределений (например, метод Алиаса или метод двоичного поиска), которые минимизируют количество таких операций.
	\end{remarkbox}
	
	\begin{theorembox}{16.15: Алгоритм табличного метода}
		\textbf{Вход:} Массив вероятностей \(p_0, p_1, \dots, p_{n-1}\) с суммой \(1\).
		
		\textbf{Выход:} Случайное число \(i \in \{0, 1, \dots, n-1\}\) с распределением \(\Pr\{i\} = p_i\).
		
		\begin{enumerate}
			\item Вычислить накопленные суммы: \(s_0 = 0\), \(s_{i+1} = s_i + p_i\) для \(i = 0, \dots, n-1\).
			\item Сгенерировать \(r \sim U[0,1)\) (равномерное на \([0,1)\)).
			\item Найти наименьшее \(i\) такое, что \(r < s_{i+1}\).
			\item Вернуть \(i\).
		\end{enumerate}
	\end{theorembox}
	
	\begin{examplebox}{16.16: Программная реализация (псевдокод)}
		\begin{verbatim}
			function discrete_sample(probabilities):
			n = length(probabilities)
			cumulative = [0] * (n+1)
			for i from 0 to n-1:
			cumulative[i+1] = cumulative[i] + probabilities[i]
			
			r = uniform_random(0, 1)  # случайное число от 0 до 1
			
			for i from 0 to n-1:
			if r < cumulative[i+1]:
			return i
		\end{verbatim}
	\end{examplebox}
	
	\begin{notebook}{16.17: Заключение}
		Табличный метод является простейшим способом моделирования дискретной случайной величины на основе равномерно распределённого непрерывного датчика. Несмотря на свою простоту и наглядность, он имеет линейную сложность по количеству сравнений, что может быть неэффективно при большом количестве значений или при необходимости многократной генерации.
		
		В таких случаях применяют более совершенные методы, например:
		\begin{itemize}
			\item метод двоичного поиска по массиву накопленных вероятностей (\(O(\log n)\));
			\item метод Алиаса (alias method) — \(O(1)\) после предварительной обработки.
		\end{itemize}
	\end{notebook}
	
	
	
	
	
	
	% ==================== РАЗДЕЛ 17: МЕТОД УОЛКЕРА ====================
	\newpage
	\section{Метод Уолкера}
	
	\subsection{Метод Уолкера моделирования дискретного распределения (17 вопрос)}
	
	\subsubsection{Введение и мотивация}
	
	\begin{definitionbox}{17.1: Напоминание основных понятий}
		\textbf{Модель ДСВ} — многократно повторяемый эксперимент, исходы которого имеют заданный набор вероятностей, соответствующий набору вероятностей дискретной случайной величины. \\
		
		\textbf{Пример:} Модель \(\xi\) с распределением \(\Pr\{\xi = 0\} = 1/2, \Pr\{\xi = 1\} = 1/2\) — это один бросок монетки.\\
		
		\textbf{Непрерывный генератор случайных чисел} — многократно повторяемый эксперимент, результаты которого — (псевдо)случайно выбранные числа из некоторого промежутка числовой прямой.\\
		
		\textbf{Пример:} \(\alpha_0\) — непрерывный генератор случайных чисел на отрезке \([0, 1]\). Отрезок \([0, 1]\) содержит несчётное количество точек, поэтому:
		\[ \Pr\{\alpha_0 = x\} = \frac{1}{\infty} = 0 \quad \forall x \in [0, 1] \]
		Однако вероятность попадания в интервал:
		\[ \Pr\{\alpha_0 \in [a, a+\delta]\} = \frac{\delta}{1} = \delta \]
		
		\textbf{Бонус:} Можем выкидывать счётное число точек, и \([a, b] \iff (a, b) \iff [a, b)\) с точки зрения вероятности.
	\end{definitionbox}
	
	\begin{importantbox}{17.2: Постановка задачи}
		Задача метода Уолкера та же, что и у табличного метода из вопроса 16: имея непрерывный генератор \(\alpha_0\) (равномерно распределённый на \([0,1]\)), нужно получать значения дискретной случайной величины с заданным распределением вероятностей \(\{p_0, p_1, \dots, p_{n-1}\}\). \\
		
		\textbf{Отличие от табличного метода:}
		\begin{itemize}
			\item Табличный метод даёт соответствие за \(O(n)\) операций (линейный поиск по массиву накопленных вероятностей).
			\item Метод Уолкера можно реализовать так, чтобы он выдавал соответствие за \(O(1)\) (константное время), потратив перед этим \(O(n)\) времени на построение модели.
		\end{itemize}
		
		\(O(1)\) — это существенно быстрее, особенно при большом количестве генераций!
	\end{importantbox}
	
	\subsubsection{Идея метода Уолкера}
	
	\begin{definitionbox}{17.3: Базовое разбиение}
		Разобьём интервал \([0, 1]\) на \(n\) равных полуинтервалов:
		\[ \left[ \frac{i}{n}, \frac{i+1}{n} \right), \quad i \in \{0, 1, \dots, n-1\} \]
		
		Каждый такой полуинтервал имеет длину \(\frac{1}{n}\).
	\end{definitionbox}
	
	\begin{center}
		\begin{tikzpicture}[scale=1.2]
			\draw[thick, blue] (0,0) -- (8,0);
			\foreach \x/\lab in {0/0,2/1/4,4/2/4,6/3/4,8/1} {
				\draw (\x,0.1) -- (\x,-0.1);
			}
			\node[below] at (0,-0.1) {0};
			\node[below] at (2,-0.1) {\(\frac{1}{4}\)};
			\node[below] at (4,-0.1) {\(\frac{2}{4}\)};
			\node[below] at (6,-0.1) {\(\frac{3}{4}\)};
			\node[below] at (8,-0.1) {1};
			
			\draw[decorate, decoration={brace, mirror, amplitude=10pt}] (0,-0.5) -- (2,-0.5);
			\node[below] at (1,-1) {\(i=0\)};
			\draw[decorate, decoration={brace, mirror, amplitude=10pt}] (2,-0.5) -- (4,-0.5);
			\node[below] at (3,-1) {\(i=1\)};
			\draw[decorate, decoration={brace, mirror, amplitude=10pt}] (4,-0.5) -- (6,-0.5);
			\node[below] at (5,-1) {\(i=2\)};
			\draw[decorate, decoration={brace, mirror, amplitude=10pt}] (6,-0.5) -- (8,-0.5);
			\node[below] at (7,-1) {\(i=3\)};
		\end{tikzpicture}
		\captionof{figure}{Разбиение \([0,1]\) на \(n=4\) равных полуинтервала}
	\end{center}
	
	\begin{remarkbox}{17.4: Ключевая идея}
		В каждом базовом полуинтервале мы размещаем "кусочки" вероятностей различных исходов ДСВ. В каждом полуинтервале есть:
		\begin{itemize}
			\item \textbf{База} — исход, чья вероятность целиком или частично занимает этот полуинтервал;
			\item \textbf{Донор} — другой исход, часть вероятности которого "досыпается" в этот полуинтервал, чтобы заполнить его до длины \(\frac{1}{n}\);
			\item \textbf{Перегородка} — граница внутри полуинтервала, разделяющая область базы и донора.
		\end{itemize}
	\end{remarkbox}
	
	\subsubsection{Алгоритм построения таблицы Уолкера}
	
	\begin{algorithmbox}{17.5: Построение таблицы Уолкера}
		\textbf{Вход:} Вероятности \(p_0, p_1, \dots, p_{n-1}\) (сумма = 1).
		
		\textbf{Выход:} Два массива длины \(n\):
		\begin{itemize}
			\item \(base[i]\) — индекс ДСВ, который является базой в \(i\)-м полуинтервале;
			\item \(prob[i]\) — положение перегородки (вероятность, соответствующая базе в этом полуинтервале, нормированная к длине полуинтервала).
		\end{itemize}
		
		\textbf{Алгоритм:}
		\begin{enumerate}
			\item Инициализируем два списка: \(less\) (индексы с вероятностью \(< \frac{1}{n}\)) и \(greater\) (индексы с вероятностью \(> \frac{1}{n}\)).
			\item Пока есть элементы в обоих списках:
			\begin{enumerate}
				\item Берём \(l\) из \(less\) и \(g\) из \(greater\).
				\item Устанавливаем \(base[l] = l\), \(prob[l] = n \cdot p_l\).
				\item Уменьшаем вероятность \(g\): \(p_g = p_g - (\frac{1}{n} - p_l)\).
				\item Если новая \(p_g < \frac{1}{n}\), перемещаем \(g\) в \(less\), иначе оставляем в \(greater\).
			\end{enumerate}
			\item Если остались элементы только в \(greater\) (с вероятностью ровно \(\frac{1}{n}\)), для них \(base[i] = i\), \(prob[i] = 1\).
		\end{enumerate}
	\end{algorithmbox}
	
	\begin{examplebox}{17.6: Пример построения}
		Рассмотрим распределение с \(n=4\) и вероятностями:
		\[ p_0 = 0.18, \quad p_1 = 0.48, \quad p_2 = 0.31, \quad p_3 = 0.03 \]
		
		Шаг 0: \(\frac{1}{n} = 0.25\). Разделяем:
		\begin{itemize}
			\item \(less\): \(\{3\}\) (0.03), \(\{0\}\) (0.18) — позже
			\item \(greater\): \(\{1\}\) (0.48), \(\{2\}\) (0.31)
		\end{itemize}
		
		\textbf{Шаг 1:}
		\begin{itemize}
			\item Берём \(l = 3\) (0.03), \(g = 1\) (0.48)
			\item \(base[3] = 3\), \(prob[3] = n \cdot p_3 = 4 \cdot 0.03 = 0.12\)
			\item \(p_1 = 0.48 - (0.25 - 0.03) = 0.48 - 0.22 = 0.26\)
			\item 0.26 > 0.25, поэтому \(1\) остаётся в \(greater\)
		\end{itemize}
		
		\textbf{Шаг 2:}
		\begin{itemize}
			\item Берём \(l = 0\) (0.18), \(g = 1\) (0.26)
			\item \(base[0] = 0\), \(prob[0] = 4 \cdot 0.18 = 0.72\)
			\item \(p_1 = 0.26 - (0.25 - 0.18) = 0.26 - 0.07 = 0.19\)
			\item 0.19 < 0.25, перемещаем \(1\) в \(less\)
		\end{itemize}
		
		\textbf{Шаг 3:}
		\begin{itemize}
			\item Берём \(l = 1\) (0.19), \(g = 2\) (0.31)
			\item \(base[1] = 1\), \(prob[1] = 4 \cdot 0.19 = 0.76\)
			\item \(p_2 = 0.31 - (0.25 - 0.19) = 0.31 - 0.06 = 0.25\)
			\item 0.25 = 0.25, это граничный случай
		\end{itemize}
		
		\textbf{Шаг 4:} Обрабатываем оставшийся \(2\) с вероятностью ровно 0.25:
		\begin{itemize}
			\item \(base[2] = 2\), \(prob[2] = 1\)
		\end{itemize}
		
		Итоговая таблица:
		\[
		\begin{array}{c|cccc}
			i & 0 & 1 & 2 & 3 \\ \hline
			base[i] & 0 & 1 & 2 & 3 \\
			prob[i] & 0.72 & 0.76 & 1.0 & 0.12
		\end{array}
		\]
	\end{examplebox}
	
	\begin{center}
		\begin{tikzpicture}[scale=1.2]
			\draw[thick, blue] (0,0) -- (8,0);
			\foreach \x in {0,2,4,6,8} {
				\draw (\x,0.1) -- (\x,-0.1);
			}
			\node[below] at (0,-0.1) {0};
			\node[below] at (2,-0.1) {0.25};
			\node[below] at (4,-0.1) {0.5};
			\node[below] at (6,-0.1) {0.75};
			\node[below] at (8,-0.1) {1};
			
			% Полуинтервал 0
			\draw[thick, red] (0,0.5) -- (2,0.5);
			\draw[dashed] (2*0.72,0.5) -- (2*0.72,0.3);
			\node[above] at (1,0.6) {база 0};
			\node[above] at (1.5,0.4) {донор 1};
			
			% Полуинтервал 1
			\draw[thick, red] (2,0.5) -- (4,0.5);
			\draw[dashed] (2+2*0.76,0.5) -- (2+2*0.76,0.3);
			\node[above] at (3,0.6) {база 1};
			\node[above] at (3.5,0.4) {донор 2};
			
			% Полуинтервал 2
			\draw[thick, red] (4,0.5) -- (6,0.5);
			\node[above] at (5,0.6) {база 2};
			
			% Полуинтервал 3
			\draw[thick, red] (6,0.5) -- (8,0.5);
			\draw[dashed] (6+2*0.12,0.5) -- (6+2*0.12,0.3);
			\node[above] at (7,0.6) {база 3};
			\node[above] at (6.7,0.4) {донор 1};
		\end{tikzpicture}
		\captionof{figure}{Визуализация разбиения для метода Уолкера}
	\end{center}
	
	\subsubsection{Генерация значений за O(1)}
	
	\begin{theorembox}{17.7: Алгоритм генерации}
		Пусть построены массивы \(base\) и \(prob\) длины \(n\). Для генерации одного значения ДСВ:
		
		\begin{enumerate}
			\item Генерируем \(r \sim U[0,1)\) от непрерывного генератора \(\alpha_0\).
			\item Вычисляем \(i = \lfloor r \cdot n \rfloor\) — номер базового полуинтервала.
			\item Вычисляем дробную часть: \(frac = \{ r \cdot n \} = r \cdot n - i\).
			\item Если \(frac < prob[i]\), возвращаем \(base[i]\) (попали в область базы).
			\item Иначе возвращаем донора для этого полуинтервала. Донора нужно где-то хранить — обычно для этого заводится третий массив \(alias[i]\).
		\end{enumerate}
	\end{theorembox}
	
	\begin{remarkbox}{17.8: Почему это работает за O(1)}
		\begin{itemize}
			\item Умножение на \(n\) и взятие целой части — \(O(1)\).
			\item Сравнение дробной части с \(prob[i]\) — \(O(1)\).
			\item Доступ к элементу массива по индексу — \(O(1)\).
		\end{itemize}
		Таким образом, каждая генерация требует константного количества операций, независимо от \(n\).
	\end{remarkbox}
	
	\begin{examplebox}{17.9: Пример генерации}
		Для построенной выше таблицы и \(n=4\):
		
		Пусть \(\alpha_0 = 0.8\).
		\begin{align*}
			r \cdot n &= 0.8 \cdot 4 = 3.2 \\
			i &= \lfloor 3.2 \rfloor = 3 \\
			frac &= 0.2
		\end{align*}
		
		Сравниваем: \(frac = 0.2 < prob[3] = 0.12\)? Нет, \(0.2 > 0.12\).
		
		Значит, попадаем в область донора. Для \(i=3\) донором был индекс 1 (из предыдущих вычислений). Возвращаем 1.
	\end{examplebox}
	
	\subsubsection{Теоретическое обоснование}
	
	\begin{lemmabox}{17.10: Лемма 1 (существование базы)}
		Если \(n > 1\), то всегда найдётся \(l \in \{0, \dots, n-1\}\) такой, что \(p_l \le \frac{1}{n}\).
		
		\textbf{Доказательство:} Предположим противное: \(\forall l, p_l > \frac{1}{n}\). Тогда сумма всех вероятностей была бы \(> n \cdot \frac{1}{n} = 1\), что противоречит \(\sum p_l = 1\). \qed
	\end{lemmabox}
	
	\begin{lemmabox}{17.11: Лемма 2 (существование донора)}
		Если \(n > 1\) и \(p_l \le \frac{1}{n}\), то всегда найдётся \(m \neq l\) такой, что \(p_l + p_m > \frac{1}{n}\).
		
		\textbf{Доказательство:} Предположим противное: \(\forall m \neq l, p_l + p_m \le \frac{1}{n}\). Тогда:
		\[ \sum_{m \neq l} (p_l + p_m) \le (n-1) \cdot \frac{1}{n} < 1 \]
		
		С другой стороны:
		\[ \sum_{m \neq l} (p_l + p_m) = (n-1)p_l + \sum_{m \neq l} p_m = (n-1)p_l + (1 - p_l) = 1 + (n-2)p_l \ge 1 \]
		
		Получили противоречие (\(< 1\) и \(\ge 1\) одновременно). \qed
	\end{lemmabox}
	
	\subsubsection{Формальное построение (общий случай)}
	
	\begin{definitionbox}{17.12: Рекурсивное построение}
		Определим \(\xi^{(0)}\) с вероятностями \(p_k^{(0)} = p_k\).
		
		На \(i\)-м шаге (\(0 \le i < n-1\)):
		\begin{enumerate}
			\item Выбираем \(l_i\) такой, что \(p_{l_i}^{(i)} \le \frac{1}{n-i}\) (база).
			\item Выбираем \(m_i \neq l_i\) такой, что \(p_{l_i}^{(i)} + p_{m_i}^{(i)} > \frac{1}{n-i}\) (донор).
			\item Определяем вспомогательную ДСВ \(\psi^{(i)}\):
			\[ \Pr\{\psi^{(i)} = l_i\} = (n-i) \cdot p_{l_i}^{(i)}, \quad \Pr\{\psi^{(i)} = m_i\} = 1 - (n-i) \cdot p_{l_i}^{(i)} \]
			\item Строим \(\xi^{(i+1)}\) с \(n-i-1\) возможными исходами:
			\[ p_{l_i}^{(i+1)} = 0, \quad p_{m_i}^{(i+1)} = \frac{n-i}{n-i-1} \left( p_{m_i}^{(i)} - \left(\frac{1}{n-i} - p_{l_i}^{(i)}\right) \right) \]
			\[ p_s^{(i+1)} = \frac{n-i}{n-i-1} p_s^{(i)} \quad \text{для остальных } s \]
		\end{enumerate}
		
		Тогда выполняется соотношение:
		\[ p_k^{(i)} = \frac{1}{n-i} \Pr\{\psi^{(i)} = k\} + \frac{n-i-1}{n-i} p_k^{(i+1)} \]
	\end{definitionbox}
	
	\begin{theorembox}{17.13: Основное равенство}
		После \(n-1\) шагов получаем:
		\[ p_k = p_k^{(0)} = \frac{1}{n} \sum_{j=0}^{n-1} \Pr\{\psi^{(j)} = k\} \]
		
		Это означает, что вероятность получить значение \(k\) равна сумме по всем полуинтервалам вероятностей того, что в \(j\)-м полуинтервале мы выберем \(k\) (как базу или донора), делённой на \(n\).
	\end{theorembox}
	
	\myproof
	Раскрывая рекурсию:
	\begin{align*}
		p_k &= p_k^{(0)} = \frac{1}{n} \Pr\{\psi^{(0)} = k\} + \frac{n-1}{n} p_k^{(1)} \\
		&= \frac{1}{n} \Pr\{\psi^{(0)} = k\} + \frac{n-1}{n} \left( \frac{1}{n-1} \Pr\{\psi^{(1)} = k\} + \frac{n-2}{n-1} p_k^{(2)} \right) \\
		&= \frac{1}{n} \Pr\{\psi^{(0)} = k\} + \frac{1}{n} \Pr\{\psi^{(1)} = k\} + \frac{n-2}{n} p_k^{(2)} \\
		&= \cdots = \frac{1}{n} \sum_{j=0}^{n-1} \Pr\{\psi^{(j)} = k\}
	\end{align*}
	\qed
	

	
	\begin{importantbox}{17.15: Преимущества метода Уолкера}
		\begin{itemize}
			\item \textbf{Скорость:} \(O(1)\) на одну генерацию после предварительной обработки.
			\item \textbf{Точность:} работает с вещественными числами, но можно адаптировать для рациональных вероятностей.
			\item \textbf{Универсальность:} подходит для любого дискретного распределения с конечным числом исходов.
		\end{itemize}
		
		\textbf{Недостатки:}
		\begin{itemize}
			\item Требует \(O(n)\) дополнительной памяти.
			\item Предварительная обработка занимает \(O(n)\) времени.
			\item Реализация сложнее табличного метода.
		\end{itemize}
	\end{importantbox}
	
	\begin{notebook}{17.16: Заключение}
		Метод Уолкера (также известный как alias method) является оптимальным способом генерации значений из дискретного распределения, когда требуется многократная генерация. Он широко применяется в компьютерном моделировании, статистических симуляциях и машинном обучении.
	\end{notebook}
	
	
	
	
	
	
	% ==================== РАЗДЕЛ 18: МОДЕЛИРОВАНИЕ С ПОМОЩЬЮ СЛУЧАЙНЫХ БИТ ====================
	\newpage
	\section{Моделирование с помощью случайных бит}
	
	\subsection{Моделирование дискретной случайной величины с помощью последовательности случайных бит (18 вопрос)}
	
	\begin{definitionbox}{18.1: Псевдослучайный бит}
		\textbf{Псевдослучайный бит} — это значение 0 или 1, которое появляется с равной вероятностью \(\frac{1}{2}\) (аналог броска монетки).
	\end{definitionbox}
	
	\myremark{Постановка задачи:} Имея источник случайных бит (бросков монетки), мы хотим моделировать дискретную случайную величину \(\xi\) с заданным распределением вероятностей \(\{p_0, p_1, \dots, p_{n-1}\}\), где \(\sum p_i = 1\). Будем строить двоичное дерево и сопоставлять каждой ДСВ с её вероятностью какой-то лист или несколько листьев дерева.
	
	\subsubsection{Основная идея}
	
	\begin{remarkbox}{18.2: Ключевая идея}
		Представим вероятности \(p_i\) в двоичной системе счисления. Тогда, если число имеет единичный бит на \(k\)-й позиции (после запятой), то ДСВ будет иметь лист на глубине \(k\) двоичного дерева.
	\end{remarkbox}
	
	\begin{examplebox}{18.3: Пример двоичного представления вероятностей}
		Пусть заданы вероятности:
		\[ p_0 = 0.101001_2, \quad p_1 = 0.000001_2, \quad p_2 = 0.001101_2, \quad p_3 = 0.001001_2 \]
		
		Для \(p_0 = 0.101001_2\) единичные биты находятся на позициях 1, 3 и 6 (считая после запятой). Это означает, что ДСВ \(p_0\) будет иметь листья на глубинах 1, 3 и 6 двоичного дерева.
	\end{examplebox}
	
	\begin{center}
		\begin{tikzpicture}[scale=0.8]
			% Рисуем дерево
			\node[circle, draw, blue] (root) at (0,0) {};
			
			% Уровень 1
			\node[circle, draw, blue] (a1) at (-2,-1) {};
			\node[circle, draw, blue] (a2) at (2,-1) {};
			\draw (root) -- (a1) node[midway, left] {0};
			\draw (root) -- (a2) node[midway, right] {1};
			
			% Уровень 2
			\node[circle, draw, blue] (b1) at (-3,-2) {};
			\node[circle, draw, blue] (b2) at (-1,-2) {};
			\draw (a1) -- (b1) node[midway, left] {0};
			\draw (a1) -- (b2) node[midway, right] {1};
			
			\node[circle, draw, blue] (b3) at (1,-2) {};
			\node[circle, draw, blue] (b4) at (3,-2) {};
			\draw (a2) -- (b3) node[midway, left] {0};
			\draw (a2) -- (b4) node[midway, right] {1};
			
			% Уровень 3
			\node[circle, draw, red] (c1) at (-3.5,-3) {\(p_0\)};
			\node[circle, draw, blue] (c2) at (-2.5,-3) {};
			\draw (b1) -- (c1) node[midway, left] {0};
			\draw (b1) -- (c2) node[midway, right] {1};
			
			% Помечаем листья
			\node[circle, draw, green] (c3) at (-1.5,-3) {\(p_2\)};
			\node[circle, draw, blue] (c4) at (-0.5,-3) {};
			\draw (b2) -- (c3) node[midway, left] {0};
			\draw (b2) -- (c4) node[midway, right] {1};
			
			\node[circle, draw, blue] (c5) at (0.5,-3) {};
			\node[circle, draw, blue] (c6) at (1.5,-3) {};
			\draw (b3) -- (c5) node[midway, left] {0};
			\draw (b3) -- (c6) node[midway, right] {1};
			
			\node[circle, draw, blue] (c7) at (2.5,-3) {};
			\node[circle, draw, blue] (c8) at (3.5,-3) {};
			\draw (b4) -- (c7) node[midway, left] {0};
			\draw (b4) -- (c8) node[midway, right] {1};
			
			\draw[decorate, decoration={brace, amplitude=10pt}] (4,-3.5) -- (4,0.5);
			\node[right] at (4.2,-1.5) {глубина};
		\end{tikzpicture}
		\captionof{figure}{Пример двоичного дерева для моделирования ДСВ}
	\end{center}
	
	\subsubsection{Формальная постановка}
	
	\begin{definitionbox}{18.4: Двоичное представление вероятностей}
		Пусть \(\xi\) — дискретная случайная величина, принимающая значения \(0, 1, \dots, n-1\) с вероятностями \(p_i\), где \(\sum_{i=0}^{n-1} p_i = 1\).
		
		Каждую вероятность \(p_i\) можно представить в двоичной системе счисления:
		\[ p_i = 0.\pi_1^i \pi_2^i \pi_3^i \dots \pi_m^i \]
		где \(\pi_k^i \in \{0, 1\}\), а \(m\) — максимальная длина двоичного представления среди всех \(p_i\).
		
		Это означает, что
		\[ p_i = \frac{\pi_1^i}{2} + \frac{\pi_2^i}{4} + \cdots + \frac{\pi_m^i}{2^m} \]
		или, умножая на \(2^m\):
		\[ 2^m p_i = 2^{m-1}\pi_1^i + 2^{m-2}\pi_2^i + \cdots + \pi_m^i \]
	\end{definitionbox}
	
	\begin{theorembox}{18.5: Моделирование с помощью \(m\) бросков монеты}
		Рассмотрим множество \(A\) всех возможных исходов \(m\) бросков монетки. Тогда \(|A| = 2^m\).
		
		Разобьём \(A\) на непересекающиеся подмножества:
		\[ A = A_0 \cup A_1 \cup \cdots \cup A_{n-1} \]
		где \(\forall i \in \{0, \dots, n-1\}: |A_i| = 2^m p_i\).
		
		Если исход эксперимента "\(m\) бросков монетки" лежит в \(A_i\), то регистрируем исход \(i\). Тогда:
		\[ \Pr\{\text{зарегистрирован исход } i\} = \frac{|A_i|}{2^m} = p_i \]
	\end{theorembox}
	
	\myproof
	Каждый из \(2^m\) исходов равновероятен с вероятностью \(\frac{1}{2^m}\). Событие \(A_i\) состоит из \(|A_i|\) элементарных исходов, поэтому его вероятность равна \(|A_i| \cdot \frac{1}{2^m} = p_i\). \qed
	
	\myremark{Интерпретация:} Мы просто разбили множество всех исходов бросков монетки на подмножества. Каждому подмножеству соответствует своё значение ДСВ, причём мощность подмножества пропорциональна вероятности этого значения.
	
	\subsubsection{Построение оптимального разбиения (метод последовательных бросков)}
	
	\begin{definitionbox}{18.6: Множества индексов}
		Для каждого уровня \(k \in \{1, 2, \dots, m\}\) определим множество:
		\[ I_k = \{ i \in \{0, 1, \dots, n-1\} \mid \pi_k^i = 1 \} \]
		то есть множество индексов тех ДСВ, у которых в двоичной записи вероятности на \(k\)-м месте (после запятой) стоит единица.
	\end{definitionbox}
	
	\begin{algorithmbox}{18.7: Алгоритм построения дерева решений}
		
		\textbf{Шаг 1:} Бросаем монету 1 раз. Множество возможных исходов \(B_1\) имеет размер 2 (\(\{0, 1\}\)).
		
		\textbf{Шаг 2:} Выбираем подмножество \(M_1 \subseteq B_1\) такое, что \(|M_1| = |I_1|\). Устанавливаем биекцию между \(M_1\) и \(I_1\). Для каждого \(i \in I_1\) сопоставляем уникальный исход из \(M_1\). В этих случаях исход эксперимента уже определён (мы получили значение ДСВ).
		
		\textbf{Шаг 3:} В оставшихся случаях \(B_1 \setminus M_1\) (их количество \(2 - |I_1|\)) нужно продолжать бросать монету.
		
		\textbf{Шаг 4:} После второго броска имеем множество исходов \(B_2\) размера \(2 \cdot |B_1 \setminus M_1|\). Убираем из этого множества те элементы, у которых префикс совпадает с путём до уже использованного листа (т.е. те, которые соответствуют уже определённым исходам).
		
		\textbf{Шаг 5:} Выбираем подмножество \(M_2 \subseteq B_2\) такое, что \(|M_2| = |I_2|\). Устанавливаем биекцию между \(M_2\) и \(I_2\) для ещё не определённых ДСВ.
		
		\textbf{Шаг 6:} Продолжаем аналогично для \(k = 3, 4, \dots\).
		
		\textbf{Шаг 7:} Если после \(k\) бросков не все ДСВ получили свои листья, бросаем монету ещё раз и получаем новое множество исходов \(B_{k+1}\). Поскольку запись вероятностей \(p_i\) конечна и их сумма равна 1, процесс обязательно завершится.
	\end{algorithmbox}
	
	\begin{remarkbox}{18.8: Важное замечание о размерах множеств}
		Заметим, что \(|I_1|\) (количество единиц на первой позиции) не может превышать 2, иначе сумма вероятностей была бы больше 1. Действительно:
		\[ 1 = \sum_{i=0}^{n-1} p_i = \sum_{i=0}^{n-1} \left( \frac{\pi_1^i}{2} + \frac{\pi_2^i}{4} + \cdots \right) = \frac{|I_1|}{2} + \frac{|I_2|}{4} + \cdots \]
		
		Аналогичные ограничения действуют и для других уровней.
	\end{remarkbox}
	
	\subsubsection{Пример для наглядности}
	
	\begin{examplebox}{18.9: Пошаговый пример}
		Рассмотрим ДСВ с четырьмя значениями и вероятностями:
		\[ p_0 = 0.101001_2, \; p_1 = 0.000001_2, \; p_2 = 0.001101_2, \; p_3 = 0.001001_2 \]
		
		Выпишем двоичные представления (до 6 знаков):
		\begin{align*}
			p_0 &= 0.101001 \quad (\pi_1^0=1, \pi_2^0=0, \pi_3^0=1, \pi_4^0=0, \pi_5^0=0, \pi_6^0=1) \\
			p_1 &= 0.000001 \quad (\pi_1^1=0, \pi_2^1=0, \pi_3^1=0, \pi_4^1=0, \pi_5^1=0, \pi_6^1=1) \\
			p_2 &= 0.001101 \quad (\pi_1^2=0, \pi_2^2=0, \pi_3^2=1, \pi_4^2=1, \pi_5^2=0, \pi_6^2=1) \\
			p_3 &= 0.001001 \quad (\pi_1^3=0, \pi_2^3=0, \pi_3^3=1, \pi_4^3=0, \pi_5^3=0, \pi_6^3=1)
		\end{align*}
		
		Определим множества \(I_k\):
		\begin{align*}
			I_1 &= \{ i \mid \pi_1^i = 1 \} = \{0\} \\
			I_2 &= \{ i \mid \pi_2^i = 1 \} = \emptyset \\
			I_3 &= \{ i \mid \pi_3^i = 1 \} = \{0, 2, 3\} \\
			I_4 &= \{ i \mid \pi_4^i = 1 \} = \{2\} \\
			I_5 &= \{ i \mid \pi_5^i = 1 \} = \emptyset \\
			I_6 &= \{ i \mid \pi_6^i = 1 \} = \{0, 1, 2, 3\}
		\end{align*}
		
		\textbf{Шаг 1:} Бросаем монету. \(B_1 = \{0, 1\}\). \(|I_1| = 1\), выбираем \(M_1 = \{0\}\). Сопоставляем исходу 0 значение \(p_0\).
		
		\textbf{Шаг 2:} Остался неиспользованный исход \(1\) из \(B_1\). Переходим ко второму броску. \(B_2\) — всевозможные последовательности из двух бросков, начинающиеся с 1: \(\{10, 11\}\). \(|I_2| = 0\), поэтому ничего не сопоставляем, продолжаем.
		
		\textbf{Шаг 3:} Третий бросок. \(B_3\) — последовательности из трёх бросков, начинающиеся с 1: \(\{100, 101, 110, 111\}\). \(|I_3| = 3\), выбираем \(M_3 = \{100, 101, 110\}\). Сопоставляем:
		\begin{itemize}
			\item 100 \(\to\) \(p_0\) (уже было, но теперь это второй лист для \(p_0\))
			\item 101 \(\to\) \(p_2\)
			\item 110 \(\to\) \(p_3\)
		\end{itemize}
		
		\textbf{Шаг 4:} Остался неиспользованный исход 111. Четвёртый бросок: \(B_4 = \{1110, 1111\}\). \(|I_4| = 1\), выбираем \(M_4 = \{1110\}\). Сопоставляем 1110 \(\to\) \(p_2\) (второй лист для \(p_2\)).
		
		\textbf{Шаг 5:} Остался 1111. Пятый бросок: \(B_5 = \{11110, 11111\}\). \(|I_5| = 0\), продолжаем.
		
		\textbf{Шаг 6:} Шестой бросок: \(B_6 = \{111100, 111101, 111110, 111111\}\). \(|I_6| = 4\), выбираем все 4 исхода и сопоставляем:
		\begin{itemize}
			\item 111100 \(\to\) \(p_0\) (третий лист)
			\item 111101 \(\to\) \(p_1\) (первый и единственный лист для \(p_1\))
			\item 111110 \(\to\) \(p_2\) (третий лист)
			\item 111111 \(\to\) \(p_3\) (второй лист)
		\end{itemize}
		
		Все вероятности распределены. Проверим количество листьев для каждой ДСВ:
		\begin{itemize}
			\item \(p_0\): листья на глубинах 1, 3, 6 → всего 3 листа
			\item \(p_1\): лист на глубине 6 → 1 лист
			\item \(p_2\): листья на глубинах 3, 4, 6 → 3 листа
			\item \(p_3\): листья на глубинах 3, 6 → 2 листа
		\end{itemize}
		
		Всего листьев: \(3+1+3+2 = 9\), что соответствует \(2^6 -\) что-то? На самом деле дерево получилось неполным, но это нормально.
	\end{examplebox}
	
	\subsubsection{Общий случай}
	
	\begin{theorembox}{18.10: Общий алгоритм}
		Пусть \(B_k\) — множество исходов после \(k\) бросков, которые ещё не привели к определению значения ДСВ (т.е. все пути, не закончившиеся в листьях).
		
		На каждом шаге \(k\):
		\begin{enumerate}
			\item Выбираем подмножество \(M_k \subseteq B_k\) такое, что \(|M_k| = |I_k|\).
			\item Строим биекцию между \(M_k\) и множеством \(I_k\) (индексами ДСВ, у которых на \(k\)-м месте стоит 1, и которые ещё не получили всех своих листьев).
			\item Для каждого \(i \in I_k\) сопоставляем уникальный исход из \(M_k\), определяя значение ДСВ для этого пути.
			\item Оставшиеся исходы \(B_k \setminus M_k\) переходят в \(B_{k+1}\) с удвоением количества (после следующего броска).
		\end{enumerate}
		
		Процесс завершается, когда все ДСВ получат все свои листья (суммарное количество листьев равно \(\sum_i 2^m p_i = 2^m\)).
	\end{theorembox}
	
	\begin{remarkbox}{18.11: Почему это работает}
		\begin{itemize}
			\item На каждом шаге мы "отрезаем" ровно \(|I_k|\) листьев на глубине \(k\).
			\item Каждый лист на глубине \(k\) имеет вероятность \(\frac{1}{2^k}\).
			\item Суммарная вероятность всех листьев, отданных ДСВ \(i\), равна
			\[ \sum_{k: \pi_k^i=1} \frac{1}{2^k} = 0.\pi_1^i\pi_2^i\ldots = p_i \]
			\item Процесс всегда завершается, так как сумма вероятностей равна 1, а значит \(\sum_k \frac{|I_k|}{2^k} = 1\).
		\end{itemize}
	\end{remarkbox}
	
	\begin{importantbox}{18.12: Преимущества метода}
		\begin{itemize}
			\item \textbf{Эффективность:} В среднем требуется \(O(\log \frac{1}{p_{\min}})\) бросков монеты.
			\item \textbf{Точность:} Метод абсолютно точен для вероятностей, представимых в виде двоичных дробей конечной длины.
			\item \textbf{Универсальность:} Любую вероятность можно сколь угодно точно приблизить двоичной дробью.
		\end{itemize}
	\end{importantbox}
	
	\begin{notebook}{18.13: Связь с кодированием}
		Данный метод тесно связан с \textbf{префиксным кодированием} (кодами Хаффмана). Фактически, мы строим дерево, в котором каждому значению ДСВ соответствует набор листьев, а пути к этим листьям являются префиксными кодами. Это позволяет эффективно генерировать значения ДСВ, используя последовательность случайных бит.
	\end{notebook}
	
	\begin{center}
		\begin{tikzpicture}[scale=0.9]
			\draw[thick, blue] (0,0) -- (6,0);
			\foreach \x in {0,1,...,6} {
				\draw (\x,0.1) -- (\x,-0.1);
				\node[below] at (\x,-0.1) {\x};
			}
			\node[above] at (3,0.2) {Номер броска \(k\)};
			
			\foreach \x/\h in {1/1,2/0,3/3,4/1,5/0,6/4} {
				\draw[fill=red] (\x,\h/4) circle (2pt);
			}
			
			\node[right] at (6.2,1) {\(|I_k|\)};
		\end{tikzpicture}
		\captionof{figure}{График зависимости \(|I_k|\) от номера броска для примера 18.9}
	\end{center}
	
	
	
	
	\subsubsection{Наглядный пример: Пьяный Слава на экзамене}
	
	\begin{examplebox}{18.14: Пример из видео "Пьяный Слава"}
		Рассмотрим забавный пример, иллюстрирующий метод моделирования ДСВ с помощью случайных бит. Слава идёт на экзамен в нетрезвом состоянии, и вероятности его оценок заданы рациональными числами со знаменателем, являющимся степенью двойки (например, \(2^5 = 32\)).
		
		\textbf{Исходные вероятности:}
		\[
		\Pr\{\text{оценка } A\} = \frac{16}{32}, \quad
		\Pr\{\text{оценка } B\} = \frac{8}{32}, \quad
		\Pr\{\text{оценка } C\} = \frac{4}{32}, \quad
		\Pr\{\text{оценка } D\} = \frac{4}{32}
		\]
		
		В двоичном виде эти вероятности представляются как:
		\[
		p_A = 0.5 = 0.1_2, \quad
		p_B = 0.25 = 0.01_2, \quad
		p_C = 0.125 = 0.001_2, \quad
		p_D = 0.125 = 0.001_2
		\]
	\end{examplebox}
	
	\begin{center}
		\begin{tikzpicture}[scale=0.9, transform shape]
			% Рисуем дерево решений
			\node[circle, draw, blue] (root) at (0,0) {};
			
			% Уровень 1
			\node[circle, draw, blue] (l1) at (-2,-1.5) {};
			\node[circle, draw, blue] (r1) at (2,-1.5) {};
			\draw (root) -- (l1) node[midway, left] {0 (решка)};
			\draw (root) -- (r1) node[midway, right] {1 (орёл)};
			
			% Лист на уровне 1 (оценка A)
			\node[circle, draw, red, fill=red!20] (A) at (-2,-2.5) {A};
			\draw (l1) -- (A);
			\node at (-2,-3) {\(p = \frac{1}{2}\)};
			
			% Уровень 2 (продолжение для орла)
			\node[circle, draw, blue] (l2) at (1,-2.5) {};
			\node[circle, draw, blue] (r2) at (3,-2.5) {};
			\draw (r1) -- (l2) node[midway, left] {0};
			\draw (r1) -- (r2) node[midway, right] {1};
			
			% Лист на уровне 2 (оценка B)
			\node[circle, draw, green, fill=green!20] (B) at (1,-3.5) {B};
			\draw (l2) -- (B);
			\node at (1,-4) {\(p = \frac{1}{4}\)};
			
			% Уровень 3 (продолжение для второго орла)
			\node[circle, draw, blue] (l3) at (3,-3.5) {};
			\node[circle, draw, blue] (r3) at (4.5,-3.5) {};
			\draw (r2) -- (l3) node[midway, left] {0};
			\draw (r2) -- (r3) node[midway, right] {1};
			
			% Листья на уровне 3 (оценки C и D)
			\node[circle, draw, orange, fill=orange!20] (C) at (2.5,-4.5) {C};
			\draw (l3) -- (C);
			\node[circle, draw, purple, fill=purple!20] (D) at (4,-4.5) {D};
			\draw (r3) -- (D);
			
			\node at (2.5,-5) {\(p = \frac{1}{8}\)};
			\node at (4,-5) {\(p = \frac{1}{8}\)};
		\end{tikzpicture}
		\captionof{figure}{Дерево решений для примера с оценками Славы}
	\end{center}
	
	\begin{remarkbox}{18.15: Пошаговое моделирование}
		Процесс моделирования происходит следующим образом:
		
		\begin{enumerate}
			\item \textbf{Первый бросок:} 
			\begin{itemize}
				\item Если выпадает \textbf{решка (0)}, сразу попадаем в лист на первом уровне — Слава получает оценку \textbf{A}. Вероятность этого исхода \(\frac{1}{2}\).
				\item Если выпадает \textbf{орёл (1)}, переходим на второй уровень.
			\end{itemize}
			
			\item \textbf{Второй бросок:}
			\begin{itemize}
				\item Если выпадает \textbf{решка (0)}, попадаем в лист на втором уровне — Слава получает оценку \textbf{B}. Вероятность этого пути \(\frac{1}{2} \cdot \frac{1}{2} = \frac{1}{4}\).
				\item Если выпадает \textbf{орёл (1)}, переходим на третий уровень.
			\end{itemize}
			
			\item \textbf{Третий бросок:}
			\begin{itemize}
				\item Если выпадает \textbf{решка (0)}, попадаем в лист — Слава получает оценку \textbf{C}.
				\item Если выпадает \textbf{орёл (1)}, попадаем в другой лист — Слава получает оценку \textbf{D}.
			\end{itemize}
			Вероятность каждого из этих путей \(\frac{1}{2} \cdot \frac{1}{2} \cdot \frac{1}{2} = \frac{1}{8}\).
		\end{enumerate}
		
		Таким образом, для получения значения ДСВ требуется не более 3 бросков, что соответствует двоичному логарифму от знаменателя вероятностей (\(\log_2 8 = 3\)).
	\end{remarkbox}
	
	\begin{importantbox}{18.16: Преимущества метода на примере}
		\begin{itemize}
			\item \textbf{Эффективность:} Случайная величина определяется за \textbf{логарифмическое время} — количество бросков равно степени двойки в знаменателе (в данном случае максимум 3 броска, хотя знаменатель 32 потребовал бы до 5 бросков).
			
			\item \textbf{Простота:} Метод позволяет генерировать сложные распределения, используя простейший источник случайности — монетку (генератор случайных бит).
			
			\item \textbf{Наглядность:} Двоичное дерево наглядно показывает, как каждому значению ДСВ соответствуют один или несколько листьев, а сумма вероятностей всех листьев равна 1.
		\end{itemize}
	\end{importantbox}
	
	\begin{notebook}{18.17: Связь с префиксным кодированием}
		Заметим, что построенное дерево является \textbf{префиксным кодом} для значений ДСВ. Пути к листьям:
		\begin{itemize}
			\item A: \(0\) (длина 1)
			\item B: \(10\) (длина 2)
			\item C: \(110\) (длина 3)
			\item D: \(111\) (длина 3)
		\end{itemize}
		
		Ни один код не является префиксом другого, что гарантирует однозначность декодирования последовательности бросков. Это не случайно — метод моделирования с помощью случайных бит тесно связан с оптимальным кодированием Хаффмана.
	\end{notebook}
	
	\begin{center}
		\begin{tabular}{|c|c|c|c|}
			\hline
			\textbf{Оценка} & \textbf{Вероятность} & \textbf{Двоичный код} & \textbf{Длина кода} \\ \hline
			A & \(\frac{1}{2}\) & 0 & 1 \\
			B & \(\frac{1}{4}\) & 10 & 2 \\
			C & \(\frac{1}{8}\) & 110 & 3 \\
			D & \(\frac{1}{8}\) & 111 & 3 \\ \hline
		\end{tabular}
		\captionof{table}{Соответствие между оценками, вероятностями и двоичными кодами}
	\end{center}
	
\end{document}
